% Le raisonnement — Philosophie 1re Tech — Cours signé OnBuch+
% Programme officiel MINESEC. Compiler avec : tectonic N14-raisonnement.tex
\newcommand{\DOCMATIERE}{Philosophie}
\newcommand{\DOCNIVEAU}{1re Tech}
\newcommand{\DOCMODULE}{Philosophie – classe de Première technique}
\newcommand{\DOCLECON}{N14}
\newcommand{\DOCTITRE}{Le raisonnement}
% OnBuch+ — préambule « cours signés » ENRICHI (compilé avec tectonic / XeLaTeX)
% Système visuel « Pop » d'OnBuch+ : fond crème (jamais blanc), bordures encre
% épaisses, ombres « dures » décalées, titres Archivo Black en capitales.
% Version MATHÉMATIQUES : graphes (pgfplots/tikz), boîtes pédagogiques
% (définition, propriété, méthode, attention, le savais-tu, exercice résolu,
% exercice, TP, bilan), page de garde et signature OnBuch+ ancrée en bas.
%
% Macros attendues dans main.tex AVANT \input{preamble} :
%   \DOCMATIERE  \DOCNIVEAU  \DOCMODULE  \DOCLECON (numéro)  \DOCTITRE
\documentclass[11pt]{article}

\usepackage{fontspec}
\usepackage[french]{babel}
\usepackage[a4paper,margin=2cm,top=3.4cm,bottom=2.8cm,headheight=1.4cm,footskip=1.3cm]{geometry}
\usepackage{amsmath,amssymb,amsfonts}
\usepackage{mathtools} % \xmapsto, \coloneqq... souvent utilisés par les modèles
\usepackage{cancel} % simplifications barrées (bilans de forces)
% \abs{z} (module, valeur absolue) et \norm{u} (norme), utilisés par les modèles
\providecommand{\abs}{}\let\abs\relax\DeclarePairedDelimiter{\abs}{\lvert}{\rvert}
\providecommand{\norm}{}\let\norm\relax\DeclarePairedDelimiter{\norm}{\lVert}{\rVert}
\usepackage[table]{xcolor} % [table] : fournit aussi \rowcolors (alternance de lignes)
\usepackage{pagecolor}
\usepackage{tikz}
\usetikzlibrary{arrows.meta,positioning,calc,shapes.geometric,shapes.misc,shapes.arrows,decorations.pathmorphing,decorations.pathreplacing,decorations.markings,patterns,shadows,mindmap,trees,fit,backgrounds,matrix,intersections,angles,quotes}
\usepackage{pgfplots}
\usepackage[version=4]{mhchem} % équations nucléaires \ce{^{235}_{92}U}
\usepackage[european,siunitx]{circuitikz} % schémas électriques
\pgfplotsset{compat=1.18}
\usepgfplotslibrary{fillbetween,groupplots} % aires entre courbes (calcul intégral)
\usepackage{enumitem}
\usepackage{fancyhdr}
% [most] charge toutes les bibliothèques tcolorbox (listings, minted, raster,
% documentation...) et gonfle inutilement la mémoire TeX sur un document long
% (~300 boîtes) : on ne charge que ce qui est réellement utilisé.
\usepackage[skins,breakable]{tcolorbox}
\usepackage{graphicx}
\usepackage{colortbl}
\usepackage{booktabs}
\usepackage{tabularx}
\usepackage{multirow}
\usepackage{diagbox} % case d'en-tête diagonale des tableaux à double entrée
% Colonnes extensibles centrée / gauche / droite (C, L, R), souvent utilisées par les modèles
\newcolumntype{C}{>{\centering\arraybackslash}X}
\newcolumntype{L}{>{\raggedright\arraybackslash}X}
\newcolumntype{R}{>{\raggedleft\arraybackslash}X}
\usepackage{array}
\usepackage{multicol}
\usepackage{siunitx}
% parse-numbers=false : accepte aussi \SI{2,5 \times 10^{-3}}{...}
\sisetup{locale=FR,output-decimal-marker={,},group-minimum-digits=4,parse-numbers=false}
\DeclareSIUnit\ohm{\ensuremath{\symup{\Omega}}} % U+2126 absent de la police
\DeclareSIUnit\division{div} % réglages d'oscilloscope (ms/div, V/div)
\DeclareSIUnit\tr{tr} % tours (vitesse de rotation, ex. \SI{1500}{\tr\per\minute})
\DeclareSIUnit\molar{M} % concentration molaire (ex. \SI{3}{\molar}, \SI{1}{\micro\molar})
\DeclareSIUnit\an{an} % année (le modèle confond parfois avec \year, un compteur TeX) — ex. \SI{5}{\kilo\meter\per\an}
\DeclareSIUnit\FCFA{FCFA} % franc CFA (ex. \SI{500}{\FCFA\per\kilo\gram})
\DeclareSIUnit\degreeBrix{\textdegree Brix} % teneur en sucre d'un sirop/jus (réfractométrie)
\DeclareSIUnit\month{mois} % le modèle utilise parfois \month, un compteur TeX, comme unité de durée
\DeclareSIUnit\microliter{\micro\liter} % le modèle écrit parfois \microliter en un seul token
\DeclareSIUnit\million{millions} % dénombrement (ex. \SI{15}{\million\per\mL} spermatozoïdes)
\usepackage{qrcode}
% \degree : commande fréquemment utilisée par les modèles pour un angle ou une
% température en dehors de \SI{}{} (non fournie par les paquets chargés).
\providecommand{\degree}{^{\circ}}
% \qed (fin de démonstration), utilisé par les modèles sans amsthm
\providecommand{\qed}{\hfill$\square$}
% \barre : double barre des valeurs interdites dans un tableau de variations (tkz-tab)
\providecommand{\barre}{\ensuremath{\|}}
% environnement proof (amsthm non chargé) utilisé par les modèles
\ifdefined\proof\else\newenvironment{proof}[1][Démonstration]{\par\noindent\textit{#1.}\ }{\qed\par}\fi
\usepackage{lastpage}
\usepackage{hyperref}
\hypersetup{hidelinks}

% ============================================================
% Palette « Pop » OnBuch+ — mêmes hex que la classe P (plus_ui.dart)
% ============================================================
\definecolor{poporange}{HTML}{F59321}
\definecolor{popdark}{HTML}{E07A0C}
\definecolor{popcream}{HTML}{FFF6E8}
\definecolor{popink}{HTML}{1C1714}
\definecolor{poppurple}{HTML}{7A5AE0}
\definecolor{popgreen}{HTML}{2E9E6A}
\definecolor{poppink}{HTML}{E0457B}
\definecolor{popblue}{HTML}{2F7DE1}
\definecolor{popgold}{HTML}{C98A12}
\definecolor{popmuted}{HTML}{8A7F76}
\definecolor{popline}{HTML}{EADFD2}
\definecolor{popgray}{HTML}{8A7F76}
\colorlet{popgrey}{popgray}
% Alias tolérants (le modèle nomme parfois les couleurs différemment)
\colorlet{popwhite}{white}
\colorlet{popblack}{popink}
\colorlet{popred}{poppink}
\colorlet{popyellow}{popgold}
% Teintes claires (fonds de tableaux/encadrés) — le modèle nomme parfois
% les couleurs avec un suffixe L (ex. popblueL) sans les avoir définies.
\colorlet{poporangeL}{poporange!20}
\colorlet{popdarkL}{popdark!20}
\colorlet{poppurpleL}{poppurple!20}
\colorlet{popgreenL}{popgreen!20}
\colorlet{poppinkL}{poppink!20}
\colorlet{popblueL}{popblue!20}
\colorlet{popgoldL}{popgold!20}
\colorlet{popredL}{popred!20}
\colorlet{popgrayL}{popgray!20}
% Teintes claires pour fonds de tableaux / zones
\colorlet{popblueL}{popblue!10!white}
\colorlet{poporangeL}{poporange!14!white}
\colorlet{popgreenL}{popgreen!12!white}
\colorlet{poppurpleL}{poppurple!12!white}
\colorlet{poppinkL}{poppink!12!white}
\colorlet{popgoldL}{popgold!14!white}

\pagecolor{popcream}

% ============================================================
% Typographie
% ============================================================
\defaultfontfeatures{Ligatures=TeX}
\setmainfont{PlusJakartaSans-Medium.ttf}[Path=../../../fonts/,BoldFont=PlusJakartaSans-Bold.ttf]
% Maths en sans-serif (Fira Math) avec chiffres et lettres droites de Plus
% Jakarta, pour que les formules se fondent dans le texte.
\usepackage[math-style=ISO,bold-style=ISO]{unicode-math}
\setmathfont{FiraMath-Regular.otf}[Path=../../../fonts/]
\setmathfont{PlusJakartaSans-Medium.ttf}[Path=../../../fonts/,range={up/num,up/{Latin,latin},\mathup}]
\usepackage{icomma} % virgule décimale sans espace en mode math
% Caractères mathématiques Unicode absents des polices de texte (Plus Jakarta,
% Archivo Black) : rendus par leur équivalent mathématique, en texte comme en
% formule — sinon ils s'affichent en carré vide (ex. « Arithmétique dans ℤ »).
\usepackage{newunicodechar}
\newunicodechar{ℝ}{\ensuremath{\mathbb{R}}}
\newunicodechar{ℤ}{\ensuremath{\mathbb{Z}}}
\newunicodechar{ℂ}{\ensuremath{\mathbb{C}}}
\newunicodechar{ℕ}{\ensuremath{\mathbb{N}}}
\newunicodechar{ℚ}{\ensuremath{\mathbb{Q}}}
\newunicodechar{𝔻}{\ensuremath{\mathbb{D}}}
\newunicodechar{α}{\ensuremath{\alpha}}
\newunicodechar{β}{\ensuremath{\beta}}
\newunicodechar{γ}{\ensuremath{\gamma}}
\newunicodechar{δ}{\ensuremath{\delta}}
\newunicodechar{ε}{\ensuremath{\varepsilon}}
\newunicodechar{θ}{\ensuremath{\theta}}
\newunicodechar{λ}{\ensuremath{\lambda}}
\newunicodechar{μ}{\ensuremath{\mu}}
\newunicodechar{σ}{\ensuremath{\sigma}}
\newunicodechar{τ}{\ensuremath{\tau}}
\newunicodechar{φ}{\ensuremath{\varphi}}
\newunicodechar{ψ}{\ensuremath{\psi}}
\newunicodechar{ω}{\ensuremath{\omega}}
\newunicodechar{Σ}{\ensuremath{\Sigma}}
% majuscules produites par \MakeUppercase dans les titres (α-aminés -> Α)
\newunicodechar{Α}{\ensuremath{\alpha}}
\newunicodechar{Β}{\ensuremath{\beta}}
\newunicodechar{Γ}{\ensuremath{\Gamma}}
\newunicodechar{Δ}{\ensuremath{\Delta}}
\newunicodechar{Ε}{\ensuremath{\varepsilon}}
\newunicodechar{Θ}{\ensuremath{\Theta}}
\newunicodechar{Λ}{\ensuremath{\Lambda}}
\newunicodechar{Μ}{\ensuremath{\mu}}
\newunicodechar{Τ}{\ensuremath{\tau}}
\newunicodechar{Φ}{\ensuremath{\Phi}}
\newunicodechar{Ψ}{\ensuremath{\Psi}}
\newunicodechar{Ω}{\ensuremath{\Omega}}
\newunicodechar{↦}{\ensuremath{\mapsto}}
\newunicodechar{∈}{\ensuremath{\in}}
\newunicodechar{∉}{\ensuremath{\notin}}
\newunicodechar{∧}{\ensuremath{\wedge}}
\newunicodechar{⇔}{\ensuremath{\Leftrightarrow}}
\newunicodechar{⟺}{\ensuremath{\Longleftrightarrow}}
\newunicodechar{⇒}{\ensuremath{\Rightarrow}}
\newunicodechar{⟹}{\ensuremath{\Longrightarrow}}
\newunicodechar{∘}{\ensuremath{\circ}}
\newunicodechar{∪}{\ensuremath{\cup}}
\newunicodechar{∩}{\ensuremath{\cap}}
\newunicodechar{≡}{\ensuremath{\equiv}}
\newunicodechar{⊂}{\ensuremath{\subset}}
\newunicodechar{⊥}{\ensuremath{\perp}}
\newunicodechar{∃}{\ensuremath{\exists}}
\newunicodechar{∀}{\ensuremath{\forall}}
\newunicodechar{∛}{\ensuremath{\sqrt[3]{\,}}}
\newunicodechar{∜}{\ensuremath{\sqrt[4]{\,}}}
\newunicodechar{⃗}{\ensuremath{{}^{\to}}}
\newunicodechar{ⁿ}{\ifmmode^{n}\else\textsuperscript{\ensuremath{n}}\fi}
\newunicodechar{ˣ}{\ifmmode^{x}\else\textsuperscript{\ensuremath{x}}\fi}
\newunicodechar{ᵇ}{\ifmmode^{b}\else\textsuperscript{\ensuremath{b}}\fi}
\newunicodechar{ʳ}{\ifmmode^{r}\else\textsuperscript{\ensuremath{r}}\fi}
\newunicodechar{ᵘ}{\ifmmode^{u}\else\textsuperscript{\ensuremath{u}}\fi}
\newunicodechar{ʸ}{\ifmmode^{y}\else\textsuperscript{\ensuremath{y}}\fi}
\newunicodechar{ᵃ}{\ifmmode^{a}\else\textsuperscript{\ensuremath{a}}\fi}
\newunicodechar{ᵉ}{\ifmmode^{e}\else\textsuperscript{\ensuremath{e}}\fi}
\newunicodechar{ᵏ}{\ifmmode^{k}\else\textsuperscript{\ensuremath{k}}\fi}
\newunicodechar{ᵖ}{\ifmmode^{p}\else\textsuperscript{\ensuremath{p}}\fi}
\newunicodechar{ᴺ}{\ifmmode^{N}\else\textsuperscript{\ensuremath{N}}\fi}
\newunicodechar{ᶜ}{\ifmmode^{c}\else\textsuperscript{\ensuremath{c}}\fi}
\newunicodechar{ʰ}{\ifmmode^{h}\else\textsuperscript{\ensuremath{h}}\fi}
\newunicodechar{ᵠ}{\ifmmode^{q}\else\textsuperscript{\ensuremath{q}}\fi}
\newunicodechar{ₙ}{\ifmmode_{n}\else\textsubscript{\ensuremath{n}}\fi}
\newunicodechar{ₐ}{\ifmmode_{a}\else\textsubscript{\ensuremath{a}}\fi}
\newunicodechar{ᵢ}{\ifmmode_{i}\else\textsubscript{\ensuremath{i}}\fi}
\newunicodechar{ᵣ}{\ifmmode_{r}\else\textsubscript{\ensuremath{r}}\fi}
\newunicodechar{ₖ}{\ifmmode_{k}\else\textsubscript{\ensuremath{k}}\fi}
\newunicodechar{ᵦ}{\ifmmode_{\beta}\else\textsubscript{\ensuremath{\beta}}\fi}
\newunicodechar{ₓ}{\ifmmode_{x}\else\textsubscript{\ensuremath{x}}\fi}
\newfontfamily\popdisplay{ArchivoBlack-Regular.ttf}[Path=../../../fonts/]
\newfontfamily\poplogo{SpaceGrotesk-Bold.ttf}[Path=../../../fonts/]
\newfontfamily\popsemi{PlusJakartaSans-SemiBold.ttf}[Path=../../../fonts/]
\newfontfamily\popxbold{PlusJakartaSans-ExtraBold.ttf}[Path=../../../fonts/]
\newfontfamily\popxboldls{PlusJakartaSans-ExtraBold.ttf}[Path=../../../fonts/,LetterSpace=3.0]

\setlist[enumerate]{leftmargin=1.7em,itemsep=5pt,topsep=4pt}
\setlist[itemize]{leftmargin=1.7em,itemsep=4pt,topsep=4pt,label={\color{poporange}\textbullet}}
\setlength{\parindent}{0pt}
\setlength{\parskip}{5pt}
\renewcommand{\arraystretch}{1.25}

% pgfplots : style Pop par défaut
\pgfplotsset{
  every axis/.append style={axis line style={popink,line width=1pt},tick style={popink},
    grid style={popline,line width=0.5pt},label style={font=\small},tick label style={font=\footnotesize}},
  popcurve/.style={poporange,line width=1.8pt,no markers,smooth},
}
% Même style disponible dans un \draw TikZ simple (hors pgfplots)
\tikzset{popcurve/.style={poporange,line width=1.8pt}}
\tikzset{popgrid/.style={popline,very thin}}
\colorlet{popcurve}{poporange} % le nom du style est parfois utilisé comme couleur

% ============================================================
% Pill / squiggle
% ============================================================
\newcommand{\poppill}[3]{%
  \tikz[baseline=-0.55ex]\node[draw=popink,line width=1.2pt,rounded corners=2.6pt,%
    fill=#2,inner xsep=6.5pt,inner ysep=3pt]{%
    \popxboldls\fontsize{7}{7}\selectfont\color{#3}\MakeUppercase{#1}};%
}
\newcommand{\popsquiggle}[1]{%
  \tikz[baseline=-0.4ex]{
    \draw[#1,line width=2.6pt,line cap=round]
      plot [smooth] coordinates {(0,0.09) (0.34,0.01) (0.68,0.21) (1.02,0.01) (1.36,0.10)};
  }%
}
% Mot-clé surligné (marqueur orange sous le texte)
\newcommand{\cle}[1]{{\popxbold\color{popdark}#1}}

% ============================================================
% En-tête / pied
% ============================================================
\pagestyle{fancy}
\fancyhf{}
\renewcommand{\headrulewidth}{0pt}
\renewcommand{\footrulewidth}{0pt}
\lhead{\raisebox{-0.15ex}{\poplogo\fontsize{13}{13}\selectfont\color{popink}OnBuch\color{poporange}+}}
\rhead{\poppill{\DOCMATIERE\ \textbullet\ \DOCNIVEAU\ \textbullet\ Leçon \DOCLECON}{white}{popink}}
\lfoot{\popxbold\fontsize{7.6}{7.6}\selectfont\color{popmuted}\MakeUppercase{Cours signé OnBuch+}\ \textbullet\ {\popsemi\DOCTITRE}}
\rfoot{\tikz[baseline=-0.6ex]\node[rounded corners=8.5pt,fill=popink,minimum height=17pt,minimum width=17pt,inner xsep=5pt,inner sep=0]{\popxbold\fontsize{8}{8}\selectfont\color{popcream}\thepage\,/\,\pageref*{LastPage}};}

% ============================================================
% Titres
% ============================================================
\newcommand{\chaptitre}[2]{%
  \begin{tcolorbox}[enhanced,colback=poporange,colframe=popink,boxrule=2.6pt,arc=11pt,
      drop shadow=popink,left=15pt,right=15pt,top=12pt,bottom=11pt,before skip=3pt,after skip=18pt]
    \poppill{#2}{popink}{popcream}\par\vspace{9pt}
    {\raggedright\hyphenpenalty=10000\exhyphenpenalty=10000\popdisplay\fontsize{22}{24}\selectfont\color{popcream}\MakeUppercase{#1}\par}
  \end{tcolorbox}
}
% Section numérotée : pastille orange avec le numéro + titre Archivo
\newcounter{popsec}
\newcommand{\coursec}[1]{%
  \stepcounter{popsec}\setcounter{popsub}{0}%
  \par\vspace{16pt}\noindent
  \begin{minipage}{\linewidth}
  \tikz[baseline=-0.7ex]\node[draw=popink,line width=1.6pt,fill=poporange,rounded corners=4pt,minimum size=22pt,inner sep=0]{\popdisplay\fontsize{12}{12}\selectfont\color{popcream}\thepopsec};%
  \hspace{8pt}{\popdisplay\fontsize{15}{17}\selectfont\color{popink}\MakeUppercase{#1}}\par
  \vspace{4pt}\noindent{\color{poporange}\rule{36pt}{3.6pt}}
  \end{minipage}\par\nopagebreak\vspace{8pt}
}
\newcounter{popsub}
\newcommand{\courssub}[1]{%
  \stepcounter{popsub}%
  \par\vspace{9pt}\noindent{\popxbold\fontsize{11.5}{13}\selectfont\color{popink}{\color{poporange}\thepopsec.\thepopsub}\hspace{6pt}#1}\par\nopagebreak\vspace{4pt}
}

% ============================================================
% Boîtes pédagogiques — même recette : fond blanc, bordure épaisse colorée,
% ombre dure encre, badge plein. Titre optionnel [#1].
% ============================================================
\tcbset{popbase/.style={breakable,enhanced,colback=white,boxrule=2.3pt,arc=9pt,
  shadow={3pt}{-3pt}{0pt}{popink},left=14pt,right=14pt,top=11pt,bottom=11pt,before skip=11pt,after skip=15pt,
  fontupper=\popsemi\fontsize{10.6}{14.5}\selectfont\color{popink},
  fontlower=\popsemi\fontsize{10.6}{14.5}\selectfont\color{popink}}}
% Coche dessinée (le glyphe ✓ n'existe pas dans les polices du document)
\newcommand{\popcheck}[1]{\tikz[baseline=-0.5ex]\draw[#1,line width=1.6pt,line cap=round,line join=round](-0.13,0.0)--(-0.04,-0.09)--(0.13,0.1);}
\providecommand{\coche}{\popcheck{popgreen}}
\providecommand{\resultat}[1]{\par\smallskip\noindent\fcolorbox{popgreen}{popgreenL}{\parbox{\dimexpr\linewidth-2\fboxsep-2\fboxrule\relax}{\textbf{Résultat.} #1}}\par\smallskip}
\newcommand{\popboxhead}[3]{\poppill{#1}{#2}{white}\ifx\relax#3\relax\else\hspace{6pt}{\popxbold\fontsize{10.6}{12}\selectfont\color{popink}#3}\fi\par\vspace{7pt}}

\newtcolorbox{aretenir}[1][]{popbase,colframe=popgreen,before upper={\popboxhead{À retenir}{popgreen}{#1}}}
\newtcolorbox{exemplebox}[1][]{popbase,colframe=poporange,before upper={\popboxhead{Exemple}{poporange}{#1}}}
\newtcolorbox{definition}[1][]{popbase,colframe=popblue,before upper={\popboxhead{Définition}{popblue}{#1}}}
\newtcolorbox{propriete}[1][]{popbase,colframe=poppurple,before upper={\popboxhead{Propriété}{poppurple}{#1}}}
\newtcolorbox{methode}[1][]{popbase,colframe=popgold,before upper={\popboxhead{Méthode}{popgold}{#1}}}
\newtcolorbox{attention}[1][]{popbase,colframe=poppink,before upper={\popboxhead{Attention}{poppink}{#1}}}
\newtcolorbox{experience}[1][]{popbase,colframe=popblue,colback=popblueL,before upper={\popboxhead{Expérience}{popblue}{#1}}}
% « Le savais-tu ? » : carte violette pleine (contexte, culture, Cameroun)
\newtcolorbox{savaistu}[1][]{popbase,colback=poppurple,colframe=popink,
  fontupper=\popsemi\fontsize{10.4}{14.2}\selectfont\color{white},
  before upper={\poppill{Le savais-tu\,?}{white}{poppurple}\ifx\relax#1\relax\else\hspace{6pt}{\popxbold\color{white}#1}\fi\par\vspace{7pt}}}
% Exercice résolu : énoncé en haut, \tcblower puis la solution sur fond crème
\newcounter{popexo}\newcounter{popexr}
\newtcolorbox{exoresolu}[1][]{popbase,colframe=poporange,colbacklower=popcream,
  segmentation style={popink,line width=1.2pt,dashed},
  before upper={\stepcounter{popexr}\popboxhead{Exercice résolu \thepopexr}{poporange}{#1}},
  before lower={\poppill{Solution}{popink}{popcream}\par\vspace{6pt}}}
% Exercice (non résolu en place) — la correction est dans la partie Corrigés
\newtcolorbox{exercice}[1][]{popbase,colframe=popink,
  before upper={\stepcounter{popexo}\popboxhead{Exercice \thepopexo}{popink}{#1}}}
% Corrigé
\newtcolorbox{corrige}[1][]{popbase,colframe=popgreen,colback=popgreenL,
  before upper={\popboxhead{Corrigé}{popgreen}{#1}}}
% Objectifs / prérequis / bilan
\newtcolorbox{objectifs}{popbase,colframe=popink,colback=poporangeL,
  before upper={\popboxhead{Objectifs de la leçon}{popink}{}}}
\newtcolorbox{prerequis}{popbase,colframe=popink,colback=popblueL,
  before upper={\popboxhead{Prérequis}{popblue}{}}}
\newtcolorbox{bilan}{popbase,colframe=popink,colback=white,boxrule=2.8pt,
  before upper={\popboxhead{Fiche bilan}{poporange}{L'essentiel à connaître pour l'examen}}}
% Figure encadrée avec légende
\newtcolorbox{popfigure}[1][]{enhanced,breakable=false,colback=white,colframe=popline,boxrule=1.4pt,arc=8pt,
  left=10pt,right=10pt,top=10pt,bottom=8pt,before skip=10pt,after skip=12pt,halign=center,
  lower separated=false,halign lower=center,
  fontlower=\popsemi\fontsize{9}{11.5}\selectfont\color{popmuted},
  #1}
\newcommand{\legende}[1]{\par\vspace{6pt}{\popsemi\fontsize{9}{11.5}\selectfont\color{popmuted}#1}}

% ============================================================
% Signature OnBuch+
% ============================================================
\newcommand{\poplogoinline}[1]{{\poplogo\fontsize{#1}{#1}\selectfont\color{popink}OnBuch\color{poporange}+}}
\newcommand{\signatureonbuch}{%
  \begin{tcolorbox}[enhanced,colback=popink,colframe=popink,boxrule=2.2pt,arc=10pt,
      drop shadow=poporange,left=14pt,right=14pt,top=10pt,bottom=10pt,before skip=0pt,after skip=0pt]
    \tikz[baseline=-0.7ex]\node[circle,fill=popgreen,draw=popcream,line width=1.3pt,minimum size=22pt,inner sep=0]{\popcheck{white}};%
    \hspace{9pt}\begin{minipage}[c]{0.62\linewidth}
      {\popxbold\fontsize{12}{13}\selectfont\color{popcream}Signé\ \textbullet\ {\poplogo OnBuch\color{poporange}+}}\par\vspace{2pt}
      {\raggedright\popsemi\fontsize{8}{10.5}\selectfont\color{popline}\MakeUppercase{Équipe pédagogique OnBuch+}\\\MakeUppercase{Conforme au programme officiel MINESEC}\par}
    \end{minipage}\hfill
    {\poplogo\fontsize{22}{22}\selectfont\color{popcream}OnBuch\color{poporange}+}
  \end{tcolorbox}%
}

% ============================================================
% Page de garde
% #1 titre · #2 matière · #3 niveau · #4 module · #5 n° leçon · #6 durée ·
% #7 description / objectifs (texte ou itemize)
% ============================================================
\newcommand{\pagegarde}[7]{%
  \thispagestyle{empty}
  \newgeometry{a4paper,margin=1.7cm,top=1.6cm,bottom=1.4cm}%
  \begin{tikzpicture}[remember picture,overlay]
    \fill[poporange!18] ([xshift=-3.2cm,yshift=-3.6cm]current page.north east) circle (4.6cm);
    \fill[poppurple!14] ([xshift=2.4cm,yshift=7.6cm]current page.south west) circle (3.8cm);
  \end{tikzpicture}%
  {\poplogo\fontsize{30}{30}\selectfont\color{popink}OnBuch\color{poporange}+}\hfill
  \raisebox{0.6ex}{\poppill{Cours signé \textbullet\ Premium}{popink}{popcream}}\par
  \vspace{1pt}\popsquiggle{poppurple}\par
  \vspace{8pt}
  \begin{tcolorbox}[enhanced,colback=poporange,colframe=popink,boxrule=2.8pt,arc=13pt,
      drop shadow=popink,left=18pt,right=18pt,top=12pt,bottom=13pt,after skip=10pt]
    \poppill{#2 \textbullet\ #3}{popink}{popcream}\hspace{5pt}\poppill{#4}{white}{popink}\par\vspace{8pt}
    {\popdisplay\fontsize{14}{14}\selectfont\color{popink}LEÇON #5}\par\vspace{4pt}
    {\raggedright\hyphenpenalty=10000\exhyphenpenalty=10000\popdisplay\fontsize{25}{27}\selectfont\color{popcream}\MakeUppercase{#1}\par}
    \vspace{8pt}
    {\popxbold\fontsize{9.5}{9.5}\selectfont\color{popink}\MakeUppercase{Durée conseillée : #6}}
  \end{tcolorbox}
  \begin{tcolorbox}[enhanced,colback=white,colframe=popink,boxrule=2.2pt,arc=10pt,
      drop shadow=popink,left=16pt,right=16pt,top=10pt,bottom=10pt,after skip=8pt]
    \poppill{Dans cette leçon}{poppurple}{white}\par\vspace{5pt}
    \popsemi\fontsize{9.3}{12.6}\selectfont\color{popink}#7
  \end{tcolorbox}
  \noindent\begin{minipage}[t]{0.487\linewidth}
    \begin{tcolorbox}[enhanced,colback=popgreen,colframe=popink,boxrule=2.2pt,arc=10pt,
        drop shadow=popink,left=11pt,right=11pt,top=10pt,bottom=10pt,height=3.1cm,valign=center]
      \tcbox[colback=white,colframe=popink,boxrule=1.3pt,arc=5pt,boxsep=0pt,nobeforeafter,
        left=3pt,right=3pt,top=3pt,bottom=3pt,valign=center]{\qrcode[height=1.9cm]{https://play.google.com/store/apps/details?id=com.luvvix.onbuch&hl=fr}}%
      \hspace{8pt}\begin{minipage}[c]{0.5\linewidth}
        {\popdisplay\fontsize{11}{12.5}\selectfont\color{white}\MakeUppercase{Télécharge OnBuch}}\par\vspace{4pt}
        {\raggedright\popsemi\fontsize{8.4}{11}\selectfont\color{white}Gratuit \textbullet\ Google Play\\Tuteur IA, annales, concours, résultats\par}
      \end{minipage}
    \end{tcolorbox}
  \end{minipage}\hfill
  \begin{minipage}[t]{0.487\linewidth}
    \begin{tcolorbox}[enhanced,colback=poppurple,colframe=popink,boxrule=2.2pt,arc=10pt,
        drop shadow=popink,left=11pt,right=11pt,top=10pt,bottom=10pt,height=3.1cm,valign=center]
      \tikz[baseline=-0.7ex]\node[circle,fill=white,draw=popink,line width=1.3pt,minimum size=1.6cm,inner sep=0]{\popdisplay\fontsize{24}{24}\selectfont\color{poppurple}+};%
      \hspace{8pt}\begin{minipage}[c]{0.6\linewidth}
        {\popdisplay\fontsize{11}{12.5}\selectfont\color{white}\MakeUppercase{Passe à OnBuch+}}\par\vspace{4pt}
        {\raggedright\popsemi\fontsize{8.4}{11}\selectfont\color{white}Tous les cours signés, Tuteur IA illimité, fiches PDF — onglet OnBuch+ de l'app\par}
      \end{minipage}
    \end{tcolorbox}
  \end{minipage}
  \vspace{8pt}
  \popcontenu
  \vfill
  \signatureonbuch
  \restoregeometry
  \newpage
}

% Rangée « Ce cours contient » (page de garde)
\newcommand{\poptile}[3]{% #1 couleur #2 gros texte #3 légende
  \begin{tcolorbox}[enhanced,colback=#1,colframe=popink,boxrule=2pt,arc=9pt,shadow={2.5pt}{-2.5pt}{0pt}{popink},
      left=6pt,right=6pt,top=9pt,bottom=9pt,halign=center,valign=center,height=1.75cm,width=\linewidth,nobeforeafter]
    {\popdisplay\fontsize{12.5}{13}\selectfont\color{white}\MakeUppercase{#2}}\par\vspace{3pt}
    {\centering\popsemi\fontsize{7.8}{9.5}\selectfont\color{white}#3\par}
  \end{tcolorbox}}
\newcommand{\popcontenu}{%
  \par\noindent{\popxboldls\fontsize{7.5}{7.5}\selectfont\color{popmuted}CE COURS SIGNÉ CONTIENT}\par\vspace{7pt}
  \noindent\begin{minipage}[t]{0.235\linewidth}\poptile{popblue}{Cours}{complet, illustré, pas à pas}\end{minipage}\hfill
  \begin{minipage}[t]{0.235\linewidth}\poptile{poporange}{Méthodes}{et exercices résolus}\end{minipage}\hfill
  \begin{minipage}[t]{0.235\linewidth}\poptile{poppink}{Exercices}{type examen + corrigés}\end{minipage}\hfill
  \begin{minipage}[t]{0.235\linewidth}\poptile{popgreen}{Fiche bilan}{l'essentiel à retenir}\end{minipage}\par}

% Sommaire
\newtcolorbox{sommaire}{enhanced,colback=white,colframe=popink,boxrule=2.2pt,arc=10pt,
  drop shadow=popink,left=18pt,right=18pt,top=13pt,bottom=13pt,before skip=6pt,after skip=20pt,
  before upper={\poppill{Sommaire}{poppurple}{white}\par\vspace{9pt}\popsemi\fontsize{10.6}{14.5}\selectfont\color{popink}}}

% Signature de fin de cours — poussée tout en bas de la dernière page
\newcommand{\findecours}{%
  \par\vfill\nopagebreak
  \begin{center}\popsquiggle{poporange}\end{center}\vspace{4pt}
  \signatureonbuch
}
\AtBeginDocument{\def\llbracket{\mathopen{[\mkern-3.5mu[}}\def\rrbracket{\mathclose{]\mkern-3.5mu]}}} % intervalles d'entiers (glyphe ⟦ absent de la police)
% Glyphes absents de Fira Math : redéfinis à partir de glyphes disponibles
\AtBeginDocument{%
  \def\setminus{\mathbin{\backslash}}\let\smallsetminus\setminus
  \def\vdots{\mathord{\vbox{\baselineskip3.2pt\lineskiplimit0pt\kern4pt\hbox{.}\hbox{.}\hbox{.}}}}%
  \def\ddots{\mathinner{\mkern1mu\raise6.5pt\hbox{.}\mkern2mu\raise3.6pt\hbox{.}\mkern2mu\raise0.7pt\hbox{.}\mkern1mu}}%
  \def\triangle{\mathord{\scalebox{0.9}{$\symup{\Delta}$}}}%
  \def\nabla{\mathord{\raisebox{\depth}{\scalebox{1}[-1]{$\symup{\Delta}$}}}}%
  \def\checkmark{\popcheck{popdark}}%
  \def\star{\mathbin{\ast}}\def\bigstar{\mathord{\ast}}%
  \def\diamond{\mathbin{\scalebox{0.6}{$\Diamond$}}}%
  \def\bigcup{\mathop{\mathchoice{\raisebox{-0.3em}{\scalebox{1.7}{$\cup$}}}{\scalebox{1.25}{$\cup$}}{\cup}{\cup}}\displaylimits}%
  \def\bigcap{\mathop{\mathchoice{\raisebox{-0.3em}{\scalebox{1.7}{$\cap$}}}{\scalebox{1.25}{$\cap$}}{\cap}{\cap}}\displaylimits}%
  \def\lesseqqgtr{\mathrel{\lessgtr}}\def\bigoplus{\mathop{\oplus}\displaylimits}%
}
\providecommand{\ssi}{si et seulement si} % abréviation fréquente
\AtBeginDocument{\providecommand{\wideparen}{\overparen}} % arc de cercle

%% FIGFIX-BEGIN v4 — correctifs globaux des figures (docs/onbuch-generation/tools/figfix)
\usepackage{adjustbox}\usepackage{etoolbox}
% toute figure TikZ/pgfplots plus large que la ligne est réduite à la largeur disponible
\BeforeBeginEnvironment{tikzpicture}{\begin{adjustbox}{max width=\linewidth}}
\AfterEndEnvironment{tikzpicture}{\end{adjustbox}}
% légendes pgfplots lisibles par-dessus les courbes
\pgfplotsset{every axis/.append style={legend style={fill=white,fill opacity=0.92,text opacity=1,draw=black!15,inner sep=3pt}}}
% étiquettes d'arêtes (syntaxe edge["..."]) avec fond blanc
\tikzset{every edge quotes/.append style={fill=white,inner sep=1.2pt}}
%% FIGFIX-END
\begin{document}

\pagegarde{Le raisonnement}{Philosophie}{1re Tech}{Philosophie – classe de Première technique}{N14}{3 séances (fiche de progression harmonisée)}{Cette leçon étudie le raisonnement comme opération fondamentale de la pensée : structure et validité du syllogisme, distinction entre raisonnement déductif, inductif et analogique, puis règles de validité et erreurs (paralogismes et sophismes) à éviter. L'élève apprend à analyser, construire et justifier des arguments dans des situations concrètes de la vie courante camerounaise.
\vspace{4pt}
\begin{multicols}{2}\small
\begin{itemize}
\item À la fin de cette leçon, je sais définir le raisonnement et identifier ses éléments (prémisses, conclusion)
\item Décrire la structure d'un syllogisme (majeure, mineure, conclusion ; grand, moyen et petit terme)
\item Vérifier la validité formelle d'un syllogisme à partir d'exemples concrets
\item Distinguer raisonnement déductif, inductif et par analogie dans des discours de la vie courante
\item Énoncer les principales règles du raisonnement correct
\item Repérer et nommer les erreurs de raisonnement (paralogismes, sophismes) dans un discours
\end{itemize}\end{multicols}}

\begin{sommaire}
\begin{multicols}{2}
\begin{enumerate}
\item Qu'est-ce que raisonner ?
\item Le syllogisme : structure et termes
\item Les types de raisonnement
\item Les règles du raisonnement correct
\item Les erreurs de raisonnement : paralogismes et sophismes
\item Raisonner dans la vie courante : méthode et rédaction
\item Activité d'intégration
\item Méthodes et exercices résolus
\item Exercices
\item Corrigés des exercices
\item Fiche bilan
\end{enumerate}
\end{multicols}
\end{sommaire}

\begin{objectifs}
\begin{itemize}
\item À la fin de cette leçon, je sais définir le raisonnement et identifier ses éléments (prémisses, conclusion)
\item Décrire la structure d'un syllogisme (majeure, mineure, conclusion ; grand, moyen et petit terme)
\item Vérifier la validité formelle d'un syllogisme à partir d'exemples concrets
\item Distinguer raisonnement déductif, inductif et par analogie dans des discours de la vie courante
\item Énoncer les principales règles du raisonnement correct
\item Repérer et nommer les erreurs de raisonnement (paralogismes, sophismes) dans un discours
\item Appliquer les notions de l'unité à des situations concrètes de la vie courante (débats, marchés, médias, réseaux sociaux)
\item Rédiger et justifier une démarche argumentative, une réponse ou une conclusion
\end{itemize}
\end{objectifs}

\begin{prerequis}
\begin{itemize}
\item Notion d'argument et d'argumentation (leçons antérieures sur le discours et la vérité)
\item Distinction entre opinion et connaissance (Seconde / début d'année)
\item Notions de proposition, de jugement vrai ou faux
\item Vocabulaire de base : prémisse, conclusion, conséquence logique
\item Habitude d'analyser un texte philosophique court
\end{itemize}
\end{prerequis}

\newpage

\coursec{Qu'est-ce que raisonner ?}

\textbf{Situation d'accroche.} Au marché Mokolo de Yaoundé, entre les étals de tomates et de poissons fumés, Mama Ngo Bell lance à voix haute : « Tous les clients qui marchandent sont des étrangers ; or ce jeune homme marchande ; donc c'est un étranger. » Son voisin de comptoir, père de famille originaire de Bafoussam et grand amateur de marchandage, proteste aussitôt : « Moi je marchande tous les jours, et je ne suis pas un étranger ! » Qui a raison ? Le raisonnement de Mama Ngo Bell est-il solide ? Plus tard, le soir même, un élève de Première technique entend à la radio un député qui « prouve » la culpabilité d'un accusé en déclarant : « Il est coupable, car tout le monde le dit au quartier. » Le lendemain, sur WhatsApp, une rumeur circule : « Ce remède guérit tout, puisque ma tante l'a pris et elle va mieux. » Dans ces trois situations, on prétend tirer une conclusion à partir de ce qu'on affirme. Mais toutes les conclusions se valent-elles ? Comment distinguer un raisonnement solide d'un raisonnement trompeur ? C'est toute la question de cette leçon : \emph{qu'est-ce que raisonner correctement ?}

\courssub{Définition du raisonnement}

Commençons par l'intuition. Quand Mama Ngo Bell parle, elle ne se contente pas d'énoncer des faits séparés : elle \emph{enchaîne} des affirmations. Elle pose deux affirmations (« tous les clients qui marchandent sont des étrangers », « ce jeune homme marchande »), puis elle en tire une troisième (« donc c'est un étranger »). Ce mouvement de l'esprit, qui passe d'affirmations admises à une affirmation nouvelle, c'est précisément le raisonnement.

\begin{definition}[Le raisonnement]
Le \cle{raisonnement} est une opération de l'esprit qui consiste à tirer une \cle{conclusion} d'une ou plusieurs propositions admises, appelées \cle{prémisses}. C'est une suite ordonnée de propositions dont la dernière (la conclusion) découle des précédentes (les prémisses).
\end{definition}

Précisons le vocabulaire. Une \cle{proposition} est un énoncé qui peut être vrai ou faux : « ce jeune homme marchande » est une proposition, car on peut vérifier si elle est vraie ou fausse. Une question (« marchande-t-il ? ») ou un ordre (« arrête de marchander ! ») ne sont pas des propositions, car ils ne sont ni vrais ni faux. Le raisonnement est donc un enchaînement de propositions, et non un simple bavardage.

\begin{aretenir}[Les deux piliers du raisonnement]
Tout raisonnement comporte deux éléments :
\begin{itemize}
\item les \cle{prémisses} : les propositions de départ, admises comme points d'appui ;
\item la \cle{conclusion} : la proposition nouvelle, obtenue à partir des prémisses.
\end{itemize}
Le passage des prémisses à la conclusion s'appelle l'\cle{inférence}.
\end{aretenir}

\courssub{La structure minimale : des prémisses à la conclusion}

Un raisonnement possède donc une structure minimale que l'on peut schématiser ainsi : \textbf{prémisses} $\rightarrow$ \textbf{conclusion}. Apprenons à repérer cette structure dans trois situations camerounaises familières.

\begin{exemplebox}[Débat au village]
Lors d'une réunion de la communauté villageoise à Bafut, un notable déclare : « Tous ceux qui ont participé aux travaux communautaires recevront une part du harvest ; or Mbah a participé aux travaux ; donc Mbah recevra une part. »
\begin{itemize}
\item Prémisse 1 : tous les participants aux travaux recevront une part.
\item Prémisse 2 : Mbah a participé aux travaux.
\item Conclusion : donc Mbah recevra une part.
\end{itemize}
\end{exemplebox}

\begin{exemplebox}[Plaidoirie au tribunal]
Devant le tribunal de première instance de Douala, l'avocate plaide : « Nul ne peut être condamné sans preuve ; or aucune preuve n'a été apportée contre mon client ; par conséquent, mon client ne peut pas être condamné. » La conclusion est introduite par « par conséquent » ; les prémisses sont les deux affirmations qui la précèdent.
\end{exemplebox}

\begin{exemplebox}[Spot publicitaire à la radio]
Une publicité vantant une boisson énergisante affirme : « Les champions boivent cette boisson ; tu veux devenir champion ; alors bois-la ! » Ici, le raisonnement est incomplet et discutable (nous y reviendrons à la section 5), mais on y retrouve bien la structure : deux affirmations, puis une conclusion invitant à l'achat.
\end{exemplebox}

\begin{popfigure}
\begin{center}
\begin{tikzpicture}[
  prem/.style={draw=popblue, fill=popblueL, rounded corners, thick, text width=4.6cm, align=center, font=\small},
  concl/.style={draw=poporange, fill=poporangeL, rounded corners, thick, text width=5.2cm, align=center, font=\small},
  fleche/.style={-{Stealth[length=3mm]}, very thick, popdark}
]
\node[prem, align=center] (p1) at (0,2.2){\textbf{Prémisse 1}\\ Tous les clients qui marchandent sont des étrangers.};
\node[prem, align=center] (p2) at (0,0){\textbf{Prémisse 2}\\ Or ce jeune homme marchande.};
\node[concl, align=center] (c) at (8.6,1.1){\textbf{Conclusion}\\ Donc ce jeune homme est un étranger.};
\draw[fleche] (p1.east) -- (c.170);
\draw[fleche] (p2.east) -- (c.190);
\node[font=\small\itshape, popmuted] at (4.3,2.6) {inférence};
\end{tikzpicture}
\end{center}
\legende{Structure du raisonnement de Mama Ngo Bell : deux prémisses convergent vers une conclusion. La flèche symbolise l'inférence, c'est-à-dire le passage logique des prémisses à la conclusion.}
\end{popfigure}

\begin{attention}[Piège : la conclusion n'est pas toujours à la fin]
Dans le langage courant, la conclusion peut être énoncée \emph{avant} les prémisses : « Ce jeune homme est un étranger, \textbf{car} il marchande, et tous ceux qui marchandent sont des étrangers. » L'ordre des phrases dans le discours n'est pas toujours l'ordre logique. C'est le sens, et non la place des phrases, qui permet d'identifier la conclusion.
\end{attention}

\courssub{Raisonnement valide et raisonnement vrai : une distinction capitale}

Voici maintenant la distinction la plus importante de cette section, et celle qui permet de répondre au voisin de comptoir de Mama Ngo Bell. Il faut distinguer deux choses que le langage courant confond :

\begin{definition}[Validité et vérité]
\begin{itemize}
\item Un raisonnement est \cle{valide} lorsque sa \emph{forme} est correcte : si les prémisses étaient vraies, la conclusion en découlerait nécessairement. La validité concerne le \emph{lien logique} entre les propositions.
\item Un raisonnement (ou une conclusion) est \cle{vrai} lorsque son \emph{contenu} est conforme à la réalité. La vérité concerne le \emph{rapport des propositions aux faits}.
\end{itemize}
\end{definition}

Reprenons l'exemple du marché Mokolo. Le raisonnement de Mama Ngo Bell est \emph{formellement valide} : si vraiment tous les clients qui marchandent étaient des étrangers, et si ce jeune homme marchande, alors il faudrait nécessairement conclure qu'il est étranger. La forme est correcte. Mais la prémisse 1 est \emph{fausse} : beaucoup de Camerounais, comme le voisin de comptoir, marchandent sans être étrangers. La conclusion est donc fausse en l'occurrence, alors que le raisonnement est valide en forme. C'est exactement ce que le voisin a senti, sans forcément savoir le formuler.

\begin{attention}[Erreur fréquente à éviter absolument]
Ne confonds jamais « raisonnement valide » et « conclusion vraie ». Un raisonnement peut être \textbf{valide avec des prémisses fausses} : la forme est correcte, mais le résultat ne dit rien de la réalité. Inversement, une conclusion peut être vraie alors que le raisonnement qui y conduit est invalide (on arrive « juste par chance »). Au Probatoire et au Baccalauréat technique, cette distinction est souvent demandée : apprends à dire clairement « ce raisonnement est valide mais sa prémisse est fausse » ou « ce raisonnement est invalide ».
\end{attention}

\begin{exemplebox}[L'exemple d'Amina]
« Tous les élèves de 1re Tech sont des philosophes ; or Amina est en 1re Tech ; donc Amina est philosophe. » Ce raisonnement est \textbf{valide en forme} : si les deux prémisses étaient vraies, la conclusion s'imposerait. Mais la première prémisse est \textbf{à discuter en contenu} : tous les élèves de Première technique étudient la philosophie, certes, mais sont-ils pour autant des « philosophes » au sens fort (des penseurs qui produisent une réflexion originale) ? Voilà un beau sujet de débat en classe.
\end{exemplebox}

Le tableau suivant croise les deux critères et montre les quatre cas possibles.

\begin{center}
\begin{tabularx}{\linewidth}{|l|X|X|}
\hline
\rowcolor{popblueL} & \textbf{Prémisses vraies} & \textbf{Prémisses fausses} \\
\hline
\textbf{Raisonnement valide} & Conclusion nécessairement vraie. Ex. : « Tous les élèves inscrits passent l'examen ; or Amina est inscrite ; donc Amina passera l'examen. » & Conclusion non garantie. Ex. : « Tous ceux qui marchandent sont étrangers ; or ce jeune homme marchande ; donc il est étranger » (forme correcte, contenu faux). \\
\hline
\textbf{Raisonnement invalide} & Conclusion vraie par hasard. Ex. : « Il pleut à Douala ; or Douala est au Cameroun ; donc il pleut partout au Cameroun » (la conclusion peut être vraie un jour de pluie nationale, mais le lien logique est défectueux). & Conclusion non garantie et souvent fausse. Ex. : « Tous les riches sont ministres ; or ce commerçant est riche ; donc il est ministre. » \\
\hline
\end{tabularx}
\end{center}

\begin{aretenir}[La règle d'or]
Seule la combinaison « \textbf{forme valide + prémisses vraies} » garantit une conclusion vraie. Un bon raisonnement doit donc être examiné deux fois : une fois pour sa forme (la conclusion suit-elle ?), une fois pour son contenu (les prémisses sont-elles vraies ?).
\end{aretenir}

\courssub{Raisonnement et pensée : l'instrument de la recherche de la vérité}

Pourquoi raisonnons-nous ? Parce que nous ne pouvons pas tout vérifier directement par nos sens. Je ne peux pas voir si le jeune homme du marché est étranger rien qu'en le regardant ; je ne peux pas voyager dans le futur pour savoir si ma candidature à l'examen sera acceptée. Le raisonnement est l'instrument qui permet à la pensée d'aller \emph{au-delà de ce qu'elle perçoit immédiatement} : à partir de ce que je sais déjà (les prémisses), je découvre ce que je ne savais pas encore (la conclusion).

Le philosophe français René \cle{Descartes} (1596-1650), dans son \emph{Discours de la méthode}, affirme que « le bon sens est la chose du monde la mieux partagée », c'est-à-dire que tout être humain possède la capacité de raisonner. Mais il ajoute aussitôt que l'essentiel est de \emph{bien conduire} sa raison : ce n'est pas assez d'avoir un instrument, il faut apprendre à s'en servir. Le raisonnement est ainsi l'outil par excellence de la recherche de la \cle{vérité} : bien utilisé, il nous conduit de connaissances sûres à des connaissances nouvelles ; mal utilisé, il nous égare dans l'erreur, le préjugé et la rumeur.

C'est pourquoi la philosophie accorde tant d'importance à la logique : raisonner correctement, c'est se donner les moyens de ne pas être trompé — ni par les autres, ni par soi-même.

\courssub{Raisonnement et langage : le rôle des connecteurs logiques}

Le raisonnement est une opération de l'esprit, mais il s'exprime dans le \cle{langage}. Or le langage possède des mots-outils qui signalent les articulations du raisonnement : ce sont les \cle{connecteurs logiques} (ou articulations du discours). Les repérer, c'est déjà analyser le raisonnement.

\begin{center}
\begin{tabularx}{\linewidth}{|l|X|}
\hline
\rowcolor{popblueL} \textbf{Fonction} & \textbf{Connecteurs courants} \\
\hline
Introduire une prémisse (une raison) & car, parce que, puisque, en effet, vu que, or \\
\hline
Introduire la conclusion & donc, par conséquent, ainsi, c'est pourquoi, alors, d'où, de sorte que \\
\hline
Ajouter une prémisse & de plus, en outre, d'ailleurs, et \\
\hline
Marquer une opposition & mais, cependant, pourtant, néanmoins \\
\hline
Marquer une condition & si... alors, à condition que, sauf si \\
\hline
\end{tabularx}
\end{center}

\begin{exemplebox}[Repérage sur un discours politique]
Pendant une campagne électorale à Maroua, un candidat déclare : « Notre commune a besoin de routes, \textbf{car} sans routes les produits agricoles pourrissent au village. \textbf{Or} mon adversaire n'a construit aucune route en dix ans. \textbf{Par conséquent}, il faut me choisir. » Les connecteurs « car » et « or » signalent les prémisses ; « par conséquent » signale la conclusion. Une fois la structure dégagée, on peut discuter : la conclusion suit-elle vraiment ? (Voter pour un candidat est-il la seule solution au problème des routes ?)
\end{exemplebox}

\begin{savaistu}[D'où vient le mot « logique » ?]
Le mot « logique » vient du grec ancien \emph{logos}, qui signifie à la fois la \textbf{parole}, la \textbf{raison} et le \textbf{calcul}. Ce triple sens n'est pas un hasard : pour les Grecs, parler correctement, raisonner correctement et calculer correctement relevaient d'une même exigence de cohérence. C'est le philosophe \cle{Aristote} (384-322 av. J.-C.) qui, le premier, a étudié systématiquement les règles du raisonnement dans des ouvrages réunis sous le nom d'\emph{Organon} (« instrument » en grec) : la logique est bien l'instrument de la pensée.
\end{savaistu}

\courssub{Méthode : analyser un raisonnement en trois étapes}

Face à n'importe quel discours — un éditorial à la radio, un message WhatsApp, une décision de conseil de famille —, l'élève philosophe doit pouvoir démonter le raisonnement pour le soumettre à l'examen critique.

\begin{methode}[Analyser un raisonnement]
\begin{enumerate}
\item \textbf{Repérer la conclusion} : que cherche-t-on à me faire admettre ? Cherche les connecteurs de conclusion (donc, par conséquent, c'est pourquoi) ou demande-toi quelle est la thèse défendue.
\item \textbf{Relever les prémisses} : quelles raisons avance-t-on pour justifier cette conclusion ? Cherche les connecteurs de prémisse (car, puisque, or) et reformule chaque prémisse clairement.
\item \textbf{Vérifier le lien logique} : la conclusion découle-t-elle nécessairement des prémisses (validité) ? Et les prémisses sont-elles vraies (vérité) ? Un raisonnement n'est solide que s'il passe ces deux tests.
\end{enumerate}
\end{methode}

\begin{exoresolu}[La rumeur sur WhatsApp]
\textbf{Énoncé.} Un message largement partagé sur les groupes WhatsApp de Yaoundé affirme : « Ce remède traditionnel guérit la typhoïde, puisque la tante de mon voisin l'a pris et elle va mieux maintenant. Partagez massivement pour sauver des vies ! » Analyse ce raisonnement en trois étapes et dis s'il est solide.
\tcblower
\textbf{Solution détaillée.}
\begin{enumerate}
\item \emph{Conclusion} : « ce remède guérit la typhoïde » (c'est ce que le message veut faire admettre, avant d'inviter au partage).
\item \emph{Prémisse} : « la tante de mon voisin l'a pris et elle va mieux maintenant » (signalée par « puisque »).
\item \emph{Vérification du lien logique} : le raisonnement est \textbf{invalide}. D'un seul cas de guérison, on ne peut pas conclure que le remède guérit en général : la tante aurait pu guérir d'elle-même, ou prendre en même temps un traitement prescrit à l'hôpital. On passe abusivement d'un cas particulier à une loi générale. De plus, la prémisse elle-même est invérifiable (qui est cette tante ? avait-elle vraiment la typhoïde ?). Le raisonnement échoue donc aux deux tests : il n'est ni valide ni fondé sur des prémisses établies. Conclusion de l'analyse : ce message est un raisonnement trompeur ; il ne faut ni le croire ni le partager sans vérification auprès des sources médicales fiables.
\end{enumerate}
\end{exoresolu}

\begin{exemplebox}[Décision d'un conseil de famille]
À Bafoussam, un conseil de famille délibère : « Notre fils Junior réussit bien en mécanique, \textbf{car} il a obtenu les meilleures notes de sa classe en atelier. \textbf{Or} le lycée technique de la ville offre une excellente formation dans cette filière. \textbf{Donc} nous l'inscrirons en Première technique plutôt qu'au lycée général. » Ici, la structure est claire, les prémisses sont vérifiables (les notes, la qualité de l'établissement) et la conclusion est raisonnablement liée aux prémisses. Le raisonnement est discutable (on pourrait ajouter d'autres critères : les goûts de Junior, les débouchés), mais il est \emph{solide} dans sa forme : c'est un bon exemple de raisonnement pratique dans la vie courante.
\end{exemplebox}

\begin{exercice}[À toi de raisonner]
Pour chacun des énoncés suivants, repère la conclusion, relève les prémisses, puis dis si le raisonnement est valide et si les prémisses sont vraies :
\begin{enumerate}
\item « Tous les camerounais parlent français ; or ce commerçant de Buea est camerounais ; donc il parle français. »
\item « Il a réussi son Probatoire technique, donc il avait travaillé sérieusement, car tous ceux qui réussissent ont travaillé sérieusement. »
\item « Ce téléphone est de bonne qualité, puisqu'il coûte cher. »
\end{enumerate}
\end{exercice}

\begin{corrige}[Exercice : À toi de raisonner]
\begin{enumerate}
\item \textbf{Conclusion} : « ce commerçant parle français ». \textbf{Prémisses} : « tous les Camerounais parlent français » et « ce commerçant est camerounais ». Le raisonnement est \textbf{valide en forme} (modus ponens), mais la première prémisse est \textbf{fausse} : de nombreux Camerounais parlent l'anglais, ou principalement une langue nationale, sans maîtriser le français. La conclusion n'est donc pas garantie.
\item \textbf{Conclusion} : « il avait travaillé sérieusement ». \textbf{Prémisses} : « il a réussi » (Q) et « tous ceux qui réussissent ont travaillé sérieusement » (Si Q alors P). La forme logique est : Q et (Q $\rightarrow$ P), donc P. Il s'agit de l'affirmation du conséquent, qui est un raisonnement \textbf{invalide} (on ne peut pas déduire l'antécédent du conséquent). Par ailleurs, la prémisse générale est discutable : certains réussissent grâce à la chance, à une mémoire exceptionnelle ou à des facilités. Le raisonnement est donc invalide \emph{et} repose sur une généralisation abusive.
\item \textbf{Conclusion} : « ce téléphone est de bonne qualité ». \textbf{Prémisse explicite} : « il coûte cher ». C'est un enthymème : la prémisse majeure implicite est « tout ce qui coûte cher est de bonne qualité ». Si on la rend explicite, le raisonnement prend la forme valide : P (cher) et (P $\rightarrow$ Q), donc Q (qualité). La \textbf{forme est donc valide} (modus ponens), mais la prémisse majeure implicite est \textbf{fausse} (le prix élevé n'entraîne pas nécessairement la qualité : contrefaçons, marketing, défauts). Le raisonnement est valide mais non solide (prémisse fausse).
\end{enumerate}
\end{corrige}

\begin{aretenir}[L'essentiel de la section]
\begin{itemize}
\item Le \cle{raisonnement} est une suite ordonnée de propositions dont la dernière (la \cle{conclusion}) découle des précédentes (les \cle{prémisses}).
\item Il faut distinguer la \cle{validité} (correction de la forme) et la \cle{vérité} (conformité du contenu à la réalité) : seul un raisonnement valide à prémisses vraies garantit une conclusion vraie.
\item Le raisonnement est l'instrument de la recherche de la vérité : il permet de passer du connu à l'inconnu.
\item Les \cle{connecteurs logiques} (donc, car, or, par conséquent...) balisent le raisonnement dans le langage et aident à l'analyser.
\item Analyser un raisonnement se fait en trois étapes : repérer la conclusion, relever les prémisses, vérifier le lien logique et la vérité des prémisses.
\end{itemize}
\end{aretenir}

Nous savons désormais ce qu'est raisonner et comment reconnaître la structure d'un raisonnement. Mais tous les raisonnements n'ont pas la même forme. La section suivante étudiera la forme la plus célèbre et la plus rigoureuse d'entre eux : le \emph{syllogisme}, hérité d'Aristote.

\coursec{Le syllogisme : structure et termes}

Dans la section précédente, nous avons vu que raisonner, c'est tirer une conclusion d'une ou plusieurs affirmations appelées \cle{prémisses}. Mais comment une conclusion peut-elle être \emph{garantie} par ses prémisses ? C'est toute la question de la \cle{déduction}. Le philosophe grec \textbf{Aristote} a donné à la déduction sa forme la plus célèbre et la plus rigoureuse : le \cle{syllogisme}. Cette section va nous apprendre à démonter un syllogisme comme un mécanicien démonte un moteur : pièce par pièce, afin de comprendre comment chaque élément contribue à la solidité du raisonnement.

\courssub{Définition du syllogisme}

Commençons par une intuition simple. Supposez que vous sachiez deux choses : d'une part, que \emph{tous les élèves inscrits au lycée technique paient la cotisation scolaire} ; d'autre part, que \emph{Manga est un élève inscrit au lycée technique}. Vous concluez immédiatement, sans même y réfléchir : \emph{Manga paie la cotisation scolaire}. Votre esprit vient de produire un syllogisme. Vous n'avez pas eu besoin d'aller vérifier dans le registre de la comptabilité : la conclusion s'impose \emph{nécessairement}, par la seule force de la forme du raisonnement.

\begin{definition}[Le syllogisme]
Le \cle{syllogisme} est un \textbf{raisonnement déductif} composé de \textbf{deux prémisses} et d'\textbf{une conclusion}, qui relie \textbf{trois termes} et trois termes seulement. Si la forme est correcte (le syllogisme est \emph{valide}), la conclusion découle nécessairement des prémisses ; si, en plus, les prémisses sont vraies, la conclusion est nécessairement vraie (le syllogisme est \emph{solide}). C'est Aristote qui, le premier, en a donné la théorie complète.
\end{definition}

Le mot « syllogisme » vient du grec \emph{syllogismos}, qui signifie « mise en relation de discours », « calcul de propositions ». L'idée est précisément celle d'une \emph{liaison} : deux affirmations séparées produisent, lorsqu'on les rapproche, une troisième affirmation qui n'était écrite nulle part mais qui était contenue dans les deux premières.

\begin{savaistu}[Aristote, le père de la logique]
\textbf{Aristote} (384--322 av. J.-C.), disciple de Platon et précepteur d'Alexandre le Grand, est le premier penseur à avoir codifié le syllogisme dans un ensemble d'ouvrages que la tradition a appelé l'\emph{Organon} (l'« instrument » de la pensée). Sa logique a dominé la pensée occidentale pendant plus de \textbf{2000 ans} : au Moyen Âge, on disait simplement « le Philosophe » pour le désigner. Même la logique moderne, née au XIX\textsuperscript{e} siècle, reconnaît en lui le fondateur de la discipline.
\end{savaistu}

\courssub{Les trois propositions : majeure, mineure, conclusion}

Un syllogisme comporte exactement \textbf{trois propositions}, c'est-à-dire trois phrases qui affirment quelque chose et qui peuvent être vraies ou fausses.

\begin{aretenir}[Les trois propositions du syllogisme]
\begin{enumerate}
\item La \cle{majeure} : la première prémisse. Elle contient le \textbf{grand terme} et le \textbf{moyen terme}. Elle énonce généralement une vérité générale (« Tous les hommes sont mortels »).
\item La \cle{mineure} : la seconde prémisse. Elle contient le \textbf{petit terme} et le \textbf{moyen terme}. Elle applique la généralité à un cas particulier (« Or Socrate est un homme »).
\item La \cle{conclusion} : introduite par « donc », « par conséquent », « c'est pourquoi ». Elle relie le \textbf{petit terme} au \textbf{grand terme} (« Donc Socrate est mortel »).
\end{enumerate}
\end{aretenir}

On reconnaît facilement la conclusion aux petits mots qui l'annoncent : \emph{donc}, \emph{ainsi}, \emph{par suite}, \emph{il s'ensuit que}. Les prémisses, elles, sont souvent introduites par \emph{or}, \emph{car}, \emph{puisque}. Apprendre à repérer ces mots-outils dans un texte est déjà un premier pas vers l'analyse des raisonnements, compétence exigée au Probatoire et au Baccalauréat technique.

\courssub{Les trois termes : grand terme, petit terme, moyen terme}

Regardons maintenant l'exemple canonique, celui que tous les manuels de philosophie reprennent depuis l'Antiquité :

\begin{exemplebox}[Le syllogisme de Socrate]
\begin{itemize}
\item \textbf{Majeure} : Tous les hommes sont mortels.
\item \textbf{Mineure} : Or Socrate est un homme.
\item \textbf{Conclusion} : Donc Socrate est mortel.
\end{itemize}
\end{exemplebox}

Ce raisonnement ne met en jeu que \textbf{trois termes}, c'est-à-dire trois notions :

\begin{itemize}
\item Le \cle{grand terme} : « mortel ». C'est l'\emph{attribut} de la conclusion (ce qu'on affirme du sujet). On le note souvent $G$.
\item Le \cle{petit terme} : « Socrate ». C'est le \emph{sujet} de la conclusion (celui dont on parle). On le note $P$.
\item Le \cle{moyen terme} : « homme ». C'est le terme \textbf{commun aux deux prémisses}, et qui \textbf{disparaît de la conclusion}. On le note $M$.
\end{itemize}

\begin{definition}[Le moyen terme]
Le \cle{moyen terme} est le terme commun aux deux prémisses, absent de la conclusion. Il joue le rôle de \textbf{charnière} ou de \textbf{pont} : c'est lui qui relie le petit terme au grand terme, comme un passeur qui fait traverser une rivière puis s'efface. Sans moyen terme correctement employé, les deux prémisses restent juxtaposées et rien ne peut être conclu.
\end{definition}

On peut comparer le moyen terme à un \emph{intermédiaire de commerce} : le commerçant de Mokolo qui achète le manioc au producteur de Mbalmayo et le revend à la ménagère de Yaoundé. Une fois la transaction faite, le producteur et l'acheteuse sont reliés, et l'intermédiaire sort de la scène. De même, « homme » relie « Socrate » à « mortel », puis s'efface de la conclusion.

\begin{popfigure}
\begin{center}
\begin{tikzpicture}[scale=0.9]
% Grand cercle : mortels
\draw[thick, popblue, fill=popblueL] (0,0) circle (3.2);
\node[popdark, font=\bfseries] at (0,2.75) {Les mortels ($G$)};
% Cercle moyen : hommes
\draw[thick, poporange, fill=poporangeL] (-0.6,-0.4) circle (1.8);
\node[popdark, font=\bfseries] at (-0.6,0.85) {Les hommes ($M$)};
% Point Socrate
\fill[poppink] (-0.6,-0.5) circle (0.09);
\node[popdark] at (-0.6,-0.85) {\textbf{Socrate} ($P$)};
% Flèches d'inclusion
\draw[-{Stealth}, thick, popmuted] (1.6,-1.9) -- (2.6,-2.6);
\node[popmuted, align=left, font=\small] at (1.4,-2.9) {Tous les hommes\\ sont mortels};
\end{tikzpicture}
\end{center}
\legende{Diagramme de Venn du syllogisme de Socrate : la classe des hommes ($M$) est entièrement incluse dans celle des mortels ($G$) ; Socrate ($P$), étant dans la classe des hommes, est nécessairement dans celle des mortels.}
\end{popfigure}

Ce diagramme rend la nécessité de la conclusion \emph{visible} : si le cercle des hommes est tout entier à l'intérieur du cercle des mortels, alors tout point placé dans le petit cercle se trouve forcément dans le grand. La déduction n'est rien d'autre que cette inclusion des classes.

\courssub{Le schéma général du syllogisme}

On peut résumer la structure du syllogisme par un schéma étiqueté, à retenir pour l'examen :

\begin{popfigure}
\begin{center}
\begin{tikzpicture}[
box/.style={draw=popblue, thick, rounded corners, fill=popcream, text width=4.6cm, align=center, minimum height=1.5cm},
lbl/.style={font=\small\itshape, popmuted}]
\node[box, align=center] (maj) at (0,3.4){\textbf{Majeure}\\ Grand terme ($G$) + Moyen terme ($M$)\\ « Tous les hommes sont mortels »};
\node[box, align=center] (min) at (0,1.1){\textbf{Mineure}\\ Petit terme ($P$) + Moyen terme ($M$)\\ « Or Socrate est un homme »};
\node[box, draw=poporange, fill=poporangeL, align=center] (ccl) at (0,-1.4){\textbf{Conclusion}\\ Petit terme ($P$) + Grand terme ($G$)\\ « Donc Socrate est mortel »};
\draw[-{Stealth}, very thick, popblue] (maj) -- (min);
\draw[-{Stealth}, very thick, popblue] (min) -- (ccl);
\node[lbl, right=0.25cm of maj] {vérité générale};
\node[lbl, right=0.25cm of min] {cas particulier};
\node[lbl, right=0.25cm of ccl] {résultat nécessaire};
\end{tikzpicture}
\end{center}
\legende{Schéma étiqueté du syllogisme : le moyen terme $M$ figure dans les deux prémisses mais disparaît de la conclusion, qui relie directement $P$ à $G$.}
\end{popfigure}

\begin{attention}[Le piège du quatrième terme]
Un syllogisme ne doit contenir que \textbf{trois termes}. Si le moyen terme change de sens entre les deux prémisses, on obtient en réalité \emph{quatre} termes et le raisonnement s'effondre. Exemple : « La \emph{fin} justifie les moyens ; or la mort est la \emph{fin} de la vie ; donc la mort justifie les moyens. » Le mot « fin » signifie d'abord « but », puis « terme final » : ce n'est plus le même terme. Ce type d'erreur, appelé \emph{équivoque}, sera étudié dans la section sur les sophismes.
\end{attention}

\courssub{Les syllogismes catégoriques : A, E, I, O}

Les syllogismes étudiés par Aristote sont dits \cle{catégoriques} : leurs propositions relient deux termes sans condition (« si... alors »). Or une proposition catégorique peut varier selon deux dimensions :

\begin{itemize}
\item la \textbf{quantité} : elle est \emph{universelle} si elle porte sur tous les membres d'une classe (« \textbf{tous} les élèves »), \emph{particulière} si elle n'en vise que quelques-uns (« \textbf{certains} élèves ») ;
\item la \textbf{qualité} : elle est \emph{affirmative} si elle attribue (« sont »), \emph{négative} si elle refuse (« ne sont pas »).
\end{itemize}

En croisant ces deux dimensions, on obtient \textbf{quatre types de propositions}, que la tradition scolastique a désignés par les voyelles \textbf{A, E, I, O} (tirées des mots latins \emph{affirmo} « j'affirme » et \emph{nego} « je nie ») :

\begin{center}
\begin{tabularx}{\linewidth}{|c|l|X|}
\hline
\rowcolor{popblueL} \textbf{Lettre} & \textbf{Type} & \textbf{Exemple} \\
\hline
\textbf{A} & Universelle affirmative & \textbf{Tous} les élèves de Première \textbf{sont} assidus. \\
\hline
\textbf{E} & Universelle négative & \textbf{Aucun} élève de Première \textbf{n'est} dispensé d'examen. \\
\hline
\textbf{I} & Particulière affirmative & \textbf{Certains} élèves de Première \textbf{sont} internes. \\
\hline
\textbf{O} & Particulière négative & \textbf{Certains} élèves de Première \textbf{ne sont pas} internes. \\
\hline
\end{tabularx}
\end{center}

Le syllogisme de Socrate est ainsi composé de trois propositions de type \textbf{A} : c'est le schéma le plus simple et le plus solide. Mais attention : mélanger les types sans précaution produit des raisonnements qui \emph{semblent} corrects et ne le sont pas. Par exemple, de « Tous les élèves sérieux réussissent » et « Certains élèves de cette classe sont sérieux », on ne peut conclure que « \emph{certains} élèves de cette classe réussissent » (conclusion particulière), et non « tous ». Nous verrons les règles précises de validité dans la section 4.

\begin{exemplebox}[Exercice simulé : la forme trompeuse est fréquente]
Lors d'un exercice simulé en classe, on a soumis \textbf{5 syllogismes} extraits de discours publics à l'analyse des termes. Résultat : \textbf{3 syllogismes sur 5 se révèlent invalides}, soit 60 \% des raisonnements examinés. Erreur la plus fréquente : le moyen terme n'est pas pris dans toute son extension, ou bien la conclusion affirme plus que les prémisses ne le permettent. Leçon pratique : un raisonnement peut avoir l'\emph{allure} d'un syllogisme (deux prémisses, un « donc ») sans en avoir la \emph{validité}. Seule l'analyse de la structure permet de trancher.
\end{exemplebox}

\courssub{Application : un syllogisme camerounais}

Transposons maintenant le syllogisme de Socrate dans notre contexte national :

\begin{exoresolu}[Le syllogisme du contribuable]
\textbf{Énoncé.} Analyser le raisonnement suivant : identifier la majeure, la mineure, la conclusion, ainsi que le grand terme, le petit terme et le moyen terme. « Tous les contribuables doivent payer leurs impôts ; or les commerçants de Douala sont des contribuables ; donc les commerçants de Douala doivent payer leurs impôts. »
\tcblower
\textbf{Solution détaillée.}
\begin{enumerate}
\item \textbf{Repérage de la conclusion} : elle est introduite par « donc » : \emph{les commerçants de Douala doivent payer leurs impôts}.
\item \textbf{Identification des termes de la conclusion} : le sujet est « les commerçants de Douala » : c'est le \textbf{petit terme} ($P$). L'attribut est « doivent payer leurs impôts » : c'est le \textbf{grand terme} ($G$).
\item \textbf{Recherche du moyen terme} : le terme qui figure dans les deux prémisses mais pas dans la conclusion est « contribuables » : c'est le \textbf{moyen terme} ($M$).
\item \textbf{Classement des prémisses} : la prémisse qui contient le grand terme ($G$ + $M$) est la \textbf{majeure} : « Tous les contribuables doivent payer leurs impôts ». Celle qui contient le petit terme ($P$ + $M$) est la \textbf{mineure} : « Or les commerçants de Douala sont des contribuables ».
\item \textbf{Vérification} : le moyen terme « contribuables » relie bien les deux prémisses et disparaît de la conclusion ; il n'y a que trois termes ; la forme est celle du syllogisme de Socrate (trois propositions de type A). Le raisonnement est donc \textbf{valide}.
\end{enumerate}
\end{exoresolu}

\courssub{Méthode de vérification d'un syllogisme}

\begin{methode}[Vérifier un syllogisme en trois étapes]
\begin{enumerate}
\item \textbf{Identifier les trois termes} : repérer d'abord la conclusion (grâce à « donc », « par conséquent »), y relever le petit terme (sujet) et le grand terme (attribut), puis trouver dans les prémisses le terme commun : le moyen terme.
\item \textbf{Vérifier le rôle du moyen terme} : s'assurer qu'il figure bien dans les deux prémisses, qu'il garde le même sens dans les deux, et qu'il relie effectivement le petit terme au grand terme. S'il est absent d'une prémisse ou s'il change de sens, le raisonnement est invalide.
\item \textbf{Contrôler la conclusion} : vérifier qu'elle ne contient que le petit terme et le grand terme, et qu'elle n'affirme pas plus que ce que les prémisses permettent (par exemple, ne pas conclure « tous » quand les prémisses ne garantissent que « certains »).
\end{enumerate}
\end{methode}

\courssub{TP guidé : construire un syllogisme à partir du règlement intérieur}

\begin{experience}[Atelier : fabriquer un syllogisme valide sur la vie scolaire]
\textbf{Matériel} : un extrait du règlement intérieur du lycée (par exemple : « Tout élève en retard répété est convoqué avec ses parents »), une feuille, le tableau.
\textbf{Protocole} :
\begin{enumerate}
\item Choisir dans le règlement une règle \emph{générale} : elle fournira la \textbf{majeure} (type A). Exemple : « Tous les élèves en tenue incorrecte sont refoulés à l'entrée. »
\item Constater un \emph{cas particulier} qui relève de cette règle : ce sera la \textbf{mineure}. Exemple : « Or Amina est une élève en tenue incorrecte. »
\item Tirer la \textbf{conclusion} en reliant le petit terme au grand terme, sans le moyen terme : « Donc Amina est refoulée à l'entrée. »
\item Vérifier avec la méthode en trois étapes : trois termes seulement (« élève en tenue incorrecte » = $M$, « Amina » = $P$, « refoulée à l'entrée » = $G$), moyen terme présent dans les deux prémisses, conclusion limitée à $P$ et $G$.
\end{enumerate}
\textbf{Observations} : les groupes qui introduisent un quatrième terme (« Amina est une élève de Terminale » au lieu de « élève en tenue incorrecte ») obtiennent une conclusion qui ne suit plus : le lien logique est rompu.
\textbf{Interprétation} : la validité d'un syllogisme dépend uniquement de sa \emph{forme} (la bonne disposition des trois termes), non du contenu concret. C'est pourquoi on peut construire des syllogismes corrects sur n'importe quel sujet, du règlement intérieur au code de la route.
\end{experience}

\begin{exercice}[À faire en classe]
Pour chacun des raisonnements suivants, identifier les trois termes et dire si la conclusion est correctement reliée aux prémisses :
\begin{enumerate}
\item « Tous les chauffeurs prudents respectent le code de la route ; or Etienne est un chauffeur prudent ; donc Etienne respecte le code de la route. »
\item « Tous les élèves boursiers sont assidus ; or certains élèves de Première technique sont boursiers ; donc tous les élèves de Première technique sont assidus. »
\end{enumerate}
\end{exercice}

\begin{corrige}[Exercice]
\begin{enumerate}
\item Petit terme : « Etienne » ; grand terme : « respecte le code de la route » ; moyen terme : « chauffeur prudent », présent dans les deux prémisses et absent de la conclusion. Le raisonnement est \textbf{valide} (même forme que le syllogisme de Socrate).
\item Petit terme : « les élèves de Première technique » ; grand terme : « assidus » ; moyen terme : « boursiers ». Mais la mineure est \emph{particulière} (« certains ») : elle ne garantit le lien que pour une partie des élèves. La conclusion « tous » affirme donc plus que les prémisses ne le permettent : le raisonnement est \textbf{invalide}. La conclusion correcte serait : « certains élèves de Première technique sont assidus ».
\end{enumerate}
\end{corrige}

\begin{aretenir}[L'essentiel de la section]
Un syllogisme est un raisonnement déductif à deux prémisses (majeure et mineure) et une conclusion, ne mettant en jeu que trois termes : le petit terme (sujet de la conclusion), le grand terme (attribut de la conclusion) et le moyen terme (terme charnière commun aux prémisses, absent de la conclusion). Les propositions catégoriques se classent en quatre types (A, E, I, O) selon leur quantité (universelle ou particulière) et leur qualité (affirmative ou négative). La validité dépend de la forme, pas du contenu : c'est ce que nous préciserons en étudiant les règles du raisonnement correct dans la section 4.
\end{aretenir}

\coursec{Les types de raisonnement}

Dans la section précédente, nous avons étudié le syllogisme : nous avons vu qu'il est composé de deux prémisses et d'une conclusion, et que sa validité dépend de la bonne disposition des termes (majeur, moyen, mineur). Mais le syllogisme n'est qu'\emph{une} forme de raisonnement parmi d'autres. Dans la vie quotidienne comme dans les sciences, nous raisonnons de plusieurs manières : parfois nous appliquons une règle générale à un cas précis, parfois nous tirons une loi générale de plusieurs observations, parfois encore nous transposons ce qui a marché dans un cas à un cas semblable. Ces trois grandes manières de raisonner ne se valent pas : elles n'offrent pas le même degré de certitude. Les distinguer, c'est se donner les moyens d'évaluer la solidité de nos propres conclusions et de celles des autres.

\courssub{Le raisonnement déductif : du général au particulier}

\textbf{Intuition.} Un agent de la circulation à Yaoundé arrête un moto-taxi qui vient de griller un feu rouge au carrefour Poste Centrale. Il n'a pas besoin d'hésiter : le code de la route dit que \emph{tout} conducteur doit s'arrêter au feu rouge ; ce conducteur est un conducteur ; donc il devait s'arrêter. L'agent part d'une règle qui vaut pour tous et l'applique à un cas particulier. C'est cela, déduire.

\begin{definition}[La déduction]
Le raisonnement \cle{déductif} est un raisonnement qui va \textbf{du général au particulier} : on part d'une ou plusieurs vérités générales (les prémisses) et on en tire une conclusion particulière. Si les prémisses sont \textbf{vraies} et si la forme du raisonnement est \textbf{valide}, alors la conclusion est \textbf{certaine} : elle ne peut pas être fausse.
\end{definition}

Le syllogisme étudié dans la section précédente est la forme typique de la déduction. Reprenons son mécanisme :

\begin{enumerate}
\item Prémisse majeure (générale) : tout conducteur qui grille un feu rouge commet une infraction.
\item Prémisse mineure (particulière) : or ce motocycliste a grillé un feu rouge.
\item Conclusion (particulière) : donc ce motocycliste a commis une infraction.
\end{enumerate}

La conclusion est \emph{contenue} dans les prémisses : la déduction ne fait qu'expliciter ce qui était déjà implicitement posé. C'est pourquoi elle est certaine : elle ne découvre rien de nouveau sur le monde, elle déploie la règle. Toute la mathématique procède ainsi : à partir d'axiomes et de théorèmes généraux, on démontre des résultats particuliers qui sont vrais de façon nécessaire.

\begin{exemplebox}[La déduction dans le calcul de la TVA]
Un commerçant de Douala vend un appareil à \num{100000} FCFA hors taxe. La loi fiscale camerounaise fixe le taux général de TVA à 19,25~\%. Le raisonnement est déductif :
\begin{enumerate}
\item Toute vente taxable est soumise à une TVA de 19,25~\% (règle générale).
\item Or cette vente est une vente taxable (cas particulier).
\item Donc cette vente est soumise à une TVA de 19,25~\%, soit $100000 \times 0{,}1925 = 19250$ FCFA. Le prix TTC est de \num{119250} FCFA.
\end{enumerate}
Si la règle fiscale est correctement citée et le calcul exact, la conclusion est \textbf{certaine} : personne ne peut la contester sans contester la loi elle-même.
\end{exemplebox}

\begin{attention}[Les deux conditions de la certitude déductive]
Une déduction ne garantit sa conclusion que si \textbf{deux} conditions sont réunies : (1) les prémisses sont vraies ; (2) la forme est valide. Si la majeure est fausse (« tous les élèves de Première technique aiment les mathématiques »), la conclusion sera peut-être fausse même si le raisonnement est bien construit. Une forme valide ne sauve pas une prémisse fausse : \emph{garbage in, garbage out}, disent les logiciens.
\end{attention}

\courssub{Le raisonnement inductif : du particulier au général}

\textbf{Intuition.} Un cultivateur de Bafoussam constate que les pluies sont arrivées en mars en 2023, en mars en 2024 et en mars en 2025. Il en conclut : « les pluies arriveront encore en mars cette année », et il prépare ses champs en conséquence. Son raisonnement part de cas observés (particuliers) pour aboutir à une loi générale. C'est une induction.

\begin{definition}[L'induction]
Le raisonnement \cle{inductif} est un raisonnement qui va \textbf{du particulier au général} : à partir d'un certain nombre de cas observés, on formule une conclusion générale qui dépasse les cas observés. Parce qu'elle va \emph{au-delà} de ce qui a été constaté, la conclusion inductive n'est jamais certaine : elle est seulement \textbf{probable}, et elle reste toujours \textbf{révisable}.
\end{definition}

La structure de l'induction est la suivante :
\begin{enumerate}
\item Le cas 1 présente le caractère A ; le cas 2 présente le caractère A ; \ldots{} le cas $n$ présente le caractère A.
\item Donc (probablement) tous les cas de ce type présentent le caractère A.
\end{enumerate}

L'induction est le moteur des sciences expérimentales : le physicien répète une expérience, observe une régularité, puis formule une loi. Elle est aussi le moteur des enquêtes statistiques.

\begin{exemplebox}[L'induction statistique à l'INS]
L'Institut National de la Statistique (INS) ne peut pas interroger les quelque dix millions de ménages camerounais pour connaître le taux de chômage. Il interroge un \textbf{échantillon} de ménages, choisi avec soin dans les dix régions, en milieu urbain et rural. Si 12~\% des actifs de l'échantillon sont au chômage, l'INS \emph{généralise} : le taux de chômage national est d'environ 12~\%. C'est une induction : la conclusion porte sur la population entière alors que l'observation ne porte que sur une partie. Elle est probable, d'autant plus fiable que l'échantillon est grand et représentatif, mais jamais absolument certaine : on parle d'\emph{estimation} avec une marge d'erreur.
\end{exemplebox}

\begin{savaistu}[Le problème de l'induction selon David Hume]
Le philosophe écossais David Hume (1711--1776) a montré que l'induction repose sur un pari : nous supposons que la nature est uniforme, que l'avenir ressemblera au passé. Or rien ne le \emph{garantit logiquement}. Le soleil s'est levé chaque jour depuis toujours, mais cela ne prouve pas avec certitude absolue qu'il se lèvera demain : c'est seulement extrêmement probable. L'induction est donc indispensable à la science et à la vie pratique, mais elle ne donne jamais la certitude nécessaire de la déduction.
\end{savaistu}

\begin{attention}[L'induction hâtive : généraliser trop vite]
L'erreur la plus fréquente consiste à tirer une loi générale d'un nombre de cas trop faible ou non représentatif. Exemple : « J'ai vu deux taxis rouler vite ce matin à Douala, donc tous les taximen sont des chauffards. » Deux observations ne suffisent pas à conclure sur des milliers de chauffeurs. Pour qu'une induction soit solide, il faut : des cas \textbf{nombreux}, \textbf{variés} et \textbf{représentatifs}, et l'absence de \textbf{contre-exemple} connu. Un seul contre-exemple suffit à ruiner une généralisation : il suffit d'un cygne noir pour réfuter « tous les cygnes sont blancs ».
\end{attention}

\courssub{Le raisonnement par analogie : d'un cas semblable à un cas semblable}

\textbf{Intuition.} Une habitante de Bafang voit son voisin guérir d'une fièvre après avoir pris une décoction traditionnelle. Elle se dit : « ce remède a guéri mon voisin, il me guérira aussi. » Elle ne part ni d'une loi générale (ce n'est pas une déduction), ni d'une série d'observations nombreuses (ce n'est pas une vraie induction) : elle transpose un résultat d'\emph{un cas} à \emph{un autre cas} jugé semblable. C'est le raisonnement par analogie.

\begin{definition}[L'analogie]
Le raisonnement par \cle{analogie} (ou par comparaison) va \textbf{d'un cas particulier à un autre cas particulier} : parce que deux situations se ressemblent sur plusieurs points, on conclut qu'elles se ressembleront aussi sur un point supplémentaire. Sa conclusion n'est que \textbf{plausible} : elle dépend entièrement de la pertinence des ressemblances invoquées.
\end{definition}

Sa structure est la suivante :
\begin{enumerate}
\item Le cas A possède les caractères $a$, $b$, $c$ et le caractère $d$.
\item Le cas B possède les caractères $a$, $b$, $c$.
\item Donc (vraisemblablement) le cas B possède aussi le caractère $d$.
\end{enumerate}

\begin{exemplebox}[Analogie entre deux exploitations agricoles de l'Ouest]
Un agriculteur possède deux champs dans les Hauts-Plateaux de l'Ouest. Sur le premier champ (sol volcanique, altitude élevée, climat frais), la pomme de terre a donné d'excellents rendements. Il raisonne : « mon second champ a le même sol volcanique, la même altitude et le même climat ; la pomme de terre y réussira donc aussi. » L'analogie est ici \textbf{forte}, car les ressemblances portent sur des facteurs réellement déterminants pour la culture. Elle serait faible s'il disait : « les deux champs sont clôturés de la même façon, donc le rendement sera le même » : la clôture n'a aucun rapport avec le rendement.
\end{exemplebox}

\textbf{Force et limites de l'analogie.} L'analogie est précieuse pour \emph{inventer}, \emph{imaginer des hypothèses} et \emph{faire comprendre} (les comparaisons du professeur sont des analogies pédagogiques). Mais elle ne prouve rien à elle seule : deux cas peuvent se ressembler sur dix points et différer précisément sur le point essentiel. Le remède qui a guéri le voisin peut être inefficace, voire dangereux, pour moi si ma maladie, mon âge ou mon état de santé diffèrent. L'analogie doit donc être vérifiée par l'expérience ou complétée par d'autres raisonnements.

\courssub{Synthèse comparée des trois types de raisonnement}

\begin{popfigure}
\begin{center}
\begin{tikzpicture}[
  box/.style={draw=popblue, thick, rounded corners, fill=popblueL, text width=3.1cm, align=center, minimum height=1.1cm, font=\small},
  gbox/.style={draw=popdark, thick, rounded corners, fill=popcream, text width=2.6cm, align=center, minimum height=0.9cm, font=\small\bfseries},
  arr/.style={-{Stealth[length=3mm]}, very thick, poporange},
  lab/.style={font=\small\itshape, popdark}
]
% Déduction
\node[box] (d1) at (0,2.2) {Tout conducteur doit s'arrêter au feu rouge};
\node[gbox] (dg) at (0,3.6) {Général};
\node[gbox] (dp) at (0,0.4) {Particulier};
\node[box] (d2) at (0,-0.9) {Ce motocycliste devait s'arrêter};
\draw[arr] (d1) -- (d2);
\node[lab] at (0,-1.9) {\textbf{Déduction} : conclusion certaine};
% Induction
\node[box] (i1) at (5.4,-0.9) {Pluies en mars 2023, 2024, 2025 à Bafoussam};
\node[gbox] (ip) at (5.4,0.4) {Particulier};
\node[gbox] (ig) at (5.4,3.6) {Général};
\node[box] (i2) at (5.4,2.2) {Les pluies arrivent en mars à Bafoussam};
\draw[arr] (i1) -- (i2);
\node[lab] at (5.4,-1.9) {\textbf{Induction} : conclusion probable};
% Analogie
\node[box] (a1) at (10.8,2.2) {Champ 1 : sol volcanique $\rightarrow$ bon rendement};
\node[gbox] (ap1) at (10.8,3.6) {Cas particulier};
\node[gbox] (ap2) at (10.8,0.4) {Cas semblable};
\node[box] (a2) at (10.8,-0.9) {Champ 2 : même sol $\rightarrow$ bon rendement ?};
\draw[arr] (a1) -- (a2);
\node[lab] at (10.8,-1.9) {\textbf{Analogie} : conclusion plausible};
\end{tikzpicture}
\end{center}
\legende{Les trois directions du raisonnement : la déduction descend du général au particulier, l'induction monte du particulier au général, l'analogie passe horizontalement d'un cas à un cas semblable.}
\end{popfigure}

\begin{center}
\rowcolors{1}{popcream}{white}
\begin{tabularx}{\linewidth}{|l|X|X|X|}
\hline
\rowcolor{popblueL} \textbf{Critère} & \textbf{Déduction} & \textbf{Induction} & \textbf{Analogie} \\
\hline
Point de départ & Vérité(s) générale(s) & Cas particuliers observés & Un cas particulier \\
\hline
Point d'arrivée & Cas particulier & Loi générale & Un autre cas particulier \\
\hline
Direction & Général $\rightarrow$ particulier & Particulier $\rightarrow$ général & Particulier $\rightarrow$ particulier \\
\hline
Degré de certitude & Certaine (si prémisses vraies et forme valide) & Probable (révisable) & Plausible (à vérifier) \\
\hline
Domaine typique & Mathématiques, droit, application de règles & Sciences expérimentales, statistiques, sondages & Invention, hypothèses, médecine de première intention, vie quotidienne \\
\hline
\end{tabularx}
\end{center}

\begin{popfigure}
\begin{center}
\begin{tikzpicture}
\begin{axis}[
  ybar, bar width=1.1cm,
  width=0.85\linewidth, height=6cm,
  symbolic x coords={Déduction, Induction, Analogie},
  xtick=data,
  ymin=0, ymax=100,
  ylabel={Degré de certitude (échelle indicative)},
  ytick={0,25,50,75,100},
  yticklabels={0\,\%, 25\,\%, 50\,\%, 75\,\%, 100\,\%},
  nodes near coords,
  every node near coord/.append style={font=\small\bfseries},
  axis line style={popline},
  tick style={popline},
]
\addplot[fill=popblue, draw=popdark] coordinates {(Déduction,100) (Induction,70) (Analogie,40)};
\end{axis}
\end{tikzpicture}
\end{center}
\legende{Degré de certitude relatif des trois types de raisonnement (échelle indicative) : la déduction atteint la certitude, l'induction la probabilité, l'analogie la simple plausibilité.}
\end{popfigure}

\courssub{Complément utile au Bac : l'abduction ou inférence à la meilleure explication}

Il existe un quatrième mode d'inférence, signalé ici comme complément utile pour l'épreuve : l'\cle{abduction}, appelée aussi \emph{inférence à la meilleure explication}. Elle consiste, face à un fait surprenant, à adopter l'hypothèse qui l'explique le mieux.

\begin{itemize}
\item \textbf{Le médecin} : le patient a de la fièvre, des frissons et des maux de tête ; la meilleure explication, dans une zone endémique, est le paludisme ; il prescrit un test pour confirmer.
\item \textbf{Le mécanicien} : le moteur ne démarre pas et le voyant de batterie reste éteint ; la meilleure explication est une batterie déchargée ; il la teste.
\item \textbf{Le détective} : la fenêtre est brisée de l'extérieur et le coffre est vide ; la meilleure explication est un cambriolage.
\end{itemize}

L'abduction ne prouve pas : elle propose l'hypothèse la plus vraisemblable, qu'il faudra ensuite \emph{vérifier} (par un test, une expérience, une enquête). Elle est au départ de toute démarche de diagnostic et de toute hypothèse scientifique.

\begin{methode}[Identifier le type d'un raisonnement]
Face à un raisonnement donné (dans un texte, un discours, une conversation), pose-toi trois questions dans l'ordre :
\begin{enumerate}
\item \textbf{Va-t-on du général au particulier ?} (une règle, une loi appliquée à un cas) $\rightarrow$ c'est une \textbf{déduction}. Conclusion certaine si prémisses vraies et forme valide.
\item \textbf{Va-t-on du particulier au général ?} (plusieurs observations transformées en loi) $\rightarrow$ c'est une \textbf{induction}. Conclusion probable ; vérifier le nombre et la représentativité des cas.
\item \textbf{Va-t-on d'un cas à un cas semblable ?} (comparaison, transposition) $\rightarrow$ c'est une \textbf{analogie}. Conclusion plausible ; vérifier la pertinence des ressemblances.
\end{enumerate}
Conclue toujours en évaluant le degré de certitude que le raisonnement autorise.
\end{methode}

\begin{exoresolu}[Identifier et évaluer trois raisonnements]
\textbf{Énoncé.} Pour chacun des raisonnements suivants, indiquez le type de raisonnement et évaluez sa solidité.
\begin{enumerate}
\item « Tout commerçant inscrit au registre de commerce doit payer la patente. Or Mama Fotso est une commerçante inscrite. Donc elle doit payer la patente. »
\item « Les trois élèves que j'ai interrogés n'ont pas compris la leçon. Donc toute la classe n'a rien compris. »
\item « Mon frère a réussi son Probatoire en révisant deux heures par jour. Moi aussi je réussirai en révisant deux heures par jour. »
\end{enumerate}
\tcblower
\textbf{Solution détaillée.}
\begin{enumerate}
\item On applique une règle générale (la loi fiscale) à un cas particulier (Mama Fotso) : c'est une \textbf{déduction}. La forme est valide (syllogisme correct) ; si la prémisse légale est exacte, la conclusion est \textbf{certaine}.
\item On passe de trois cas observés à une conclusion générale sur toute la classe : c'est une \textbf{induction}, mais elle est \textbf{hâtive} : trois élèves ne constituent pas un échantillon représentatif d'une classe entière. La conclusion est très faiblement probable et doit être révisée après interrogation d'un plus grand nombre d'élèves.
\item On transpose ce qui a réussi à une personne à une autre personne jugée semblable : c'est une \textbf{analogie}. Sa solidité dépend des ressemblances réelles (même classe, mêmes difficultés, même assiduité). La conclusion est \textbf{plausible} mais non garantie : la réussite dépend aussi de facteurs que la comparaison ne mentionne pas.
\end{enumerate}
\end{exoresolu}

\begin{aretenir}[L'essentiel de la section]
\begin{itemize}
\item La \textbf{déduction} va du général au particulier ; sa conclusion est \textbf{certaine} si les prémisses sont vraies et la forme valide (ex. : application du code de la route, calcul de la TVA).
\item L'\textbf{induction} va du particulier au général ; sa conclusion est \textbf{probable} et révisable (ex. : prévisions météo, sondages, estimations de l'INS). Piège : l'induction hâtive.
\item L'\textbf{analogie} va d'un cas particulier à un cas semblable ; sa conclusion est \textbf{plausible} et dépend de la pertinence des ressemblances.
\item Complément : l'\textbf{abduction} choisit la meilleure explication d'un fait (médecin, mécanicien, détective) ; elle doit être vérifiée.
\end{itemize}
\end{aretenir}

\begin{exercice}[À vous de raisonner]
\begin{enumerate}
\item Construisez un raisonnement déductif complet (deux prémisses, une conclusion) à partir de la règle : « tout élève non inscrit ne peut pas composer au Probatoire ».
\item Un candidat affirme : « J'ai visité deux villages du Nord et il n'y avait pas d'électricité ; donc tout le Nord-Cameroun n'a pas d'électricité. » Identifiez le type de raisonnement et critiquez-le.
\item Expliquez pourquoi le raisonnement du mécanicien qui diagnostique une panne est une abduction et non une déduction.
\end{enumerate}
\end{exercice}

\begin{corrige}[Exercice : À vous de raisonner]
\begin{enumerate}
\item Tout élève non inscrit ne peut pas composer au Probatoire (majeure générale). Or cet élève n'est pas inscrit (mineure particulière). Donc cet élève ne peut pas composer au Probatoire (conclusion). Déduction valide ; conclusion certaine si les prémisses sont vraies.
\item C'est une induction (du particulier au général), et elle est hâtive : deux villages ne représentent pas l'ensemble d'une vaste région. Il faudrait un échantillon nombreux et varié, ou des données officielles, pour généraliser.
\item Le mécanicien n'applique pas une règle certaine à un cas : face à un symptôme (le moteur ne démarre pas), plusieurs causes sont possibles (batterie, démarreur, carburant). Il choisit l'explication la plus vraisemblable, qu'il devra vérifier par un test. C'est le schéma de l'abduction : du fait observé vers la meilleure explication, et non d'une loi vers un cas.
\end{enumerate}
\end{corrige}

\coursec{Les règles du raisonnement correct}

Nous venons de distinguer les grands types de raisonnement : la déduction, qui va du général au particulier, et l'induction, qui va du particulier au général. Mais connaître les types de raisonnement ne suffit pas : encore faut-il raisonner \emph{correctement}. De même qu'un match de football ne se joue bien que si les joueurs respectent les règles du jeu, un raisonnement n'est recevable que s'il obéit à des règles précises. Cette section présente les six règles fondamentales du raisonnement correct, telles que la logique classique, héritée d'Aristote, les a formulées.

\courssub{Les deux conditions d'un bon raisonnement}

\begin{propriete}[Vérité et validité]
Un raisonnement correct exige à la fois la \cle{vérité} des prémisses et la \cle{validité} de la forme. Ces deux conditions sont nécessaires : si les prémisses sont fausses, la conclusion n'a aucune valeur, même si la forme est bonne ; si la forme est invalide, la conclusion ne suit pas, même si les prémisses sont vraies.
\end{propriete}

Il faut bien distinguer ces deux exigences. La \textbf{validité} concerne la \emph{forme} du raisonnement : la conclusion découle-t-elle nécessairement des prémisses ? La \textbf{vérité} concerne le \emph{contenu} : ce que disent les prémisses correspond-il à la réalité ? Un raisonnement peut être valide mais fondé sur du faux : « Tous les élèves de Première technique sont des enseignants ; or Jean est un élève de Première technique ; donc Jean est un enseignant. » La forme est irréprochable, mais la première prémisse est fausse, donc la conclusion est fausse. Inversement, on peut partir de prémisses vraies et mal conclure : « Tous les chiens sont des animaux ; or mon chat est un animal ; donc mon chat est un chien. »

\courssub{Règle 1 : partir de prémisses vraies ou admises}

\textbf{Intuition.} Un bâtiment construit sur du sable s'effondre, aussi belle soit son architecture. De même, on ne peut rien conclure de solide à partir du faux.

\textbf{Énoncé.} Les prémisses d'un raisonnement doivent être vraies, ou du moins admises par tous les participants à la discussion. Si l'on part d'une prémisse fausse, la conclusion, même obtenue par une forme valide, ne prouve rien.

\begin{exemplebox}[Au marché de Mokolo]
« Tous les vendeurs de ce marché trichent sur les poids ; or cette dame est vendeuse ici ; donc elle triche. » La forme est valide, mais la première prémisse est une généralisation abusive (un préjugé) : la conclusion est donc sans valeur. Il fallait d'abord vérifier la prémisse.
\end{exemplebox}

\courssub{Règle 2 : respecter les principes logiques fondamentaux}

Tout raisonnement repose sur trois principes que la tradition philosophique appelle les \cle{principes de la pensée}.

\begin{definition}[Principe de non-contradiction]
Le \cle{principe de non-contradiction} interdit d'affirmer et de nier la même chose, au même moment et sous le même rapport. On ne peut pas dire à la fois « ce sac contient dix kilogrammes de riz » et « ce sac ne contient pas dix kilogrammes de riz ».
\end{definition}

À ce principe s'ajoutent deux autres :

\begin{itemize}
\item Le \textbf{principe d'identité} : une chose est ce qu'elle est ($A$ est $A$). Dans un raisonnement, chaque terme désigne toujours la même réalité.
\item Le \textbf{principe du tiers exclu} : entre une affirmation et sa négation, il n'y a pas de troisième possibilité. Un élève est présent ou absent ; il n'existe pas de position intermédiaire entre « présent » et « absent ».
\end{itemize}

\begin{attention}[Attention aux apparences de contradiction]
« Ce jeune homme est un enfant pour sa mère, mais un adulte pour l'état civil » n'est pas une contradiction, car on ne parle pas \emph{sous le même rapport}. La contradiction n'existe que si l'on affirme et nie la même chose, du même sujet, au même moment et sous le même rapport.
\end{attention}

\courssub{Règle 3 : le moyen terme doit être pris au moins une fois dans toute son extension}

\textbf{Intuition.} Dans un syllogisme, le moyen terme est le « pont » qui relie les deux autres termes. Si ce pont ne porte pas sur la totalité de ce qu'il désigne au moins une fois, la liaison peut échouer : les deux termes extrêmes peuvent toucher des parties différentes du moyen terme sans jamais se rencontrer.

\textbf{Énoncé.} On dit qu'un terme est \cle{distribué} (ou pris dans toute son extension) lorsqu'il est pris universellement : « tous les hommes », « aucun voleur ». Le moyen terme doit être distribué au moins une fois dans les prémisses.

\begin{exemplebox}[Le sophisme du stade]
« Tous les voleurs sont des hommes ; or Paul est un homme ; donc Paul est un voleur. » Le moyen terme « homme » n'est jamais pris dans toute son extension : on dit seulement que les voleurs sont \emph{parmi} les hommes, et que Paul est \emph{parmi} les hommes. Paul peut être dans la partie des hommes qui ne sont pas voleurs. La conclusion ne suit pas.
\end{exemplebox}

\begin{popfigure}
\begin{center}
\begin{tikzpicture}[scale=0.9]
\draw[thick, popblue] (0,0) ellipse (3.4 and 2.1);
\node[popblue, font=\bfseries] at (-2.3,1.6) {Hommes};
\draw[thick, poporange, fill=poporangeL] (-1.1,0.2) circle (0.85);
\node[popdark, font=\small] at (-1.1,0.2) {Voleurs};
\fill[popgreen] (1.6,-0.6) circle (0.07);
\node[popgreen, font=\small, below] at (1.6,-0.75) {Paul (possible)};
\fill[popgreen] (-1.1,0.9) circle (0.07);
\node[popgreen, font=\small, above] at (-1.1,1.0) {Paul (possible)};
\end{tikzpicture}
\end{center}
\legende{Erreur du moyen terme non distribué. Les voleurs sont inclus dans les hommes, et Paul est un homme ; mais rien ne dit où se situe Paul : il peut être dans le cercle des voleurs ou en dehors. La conclusion « Paul est un voleur » n'est donc pas nécessaire.}
\end{popfigure}

\courssub{Règle 4 : de deux prémisses négatives ou particulières, on ne peut rien conclure}

\textbf{Énoncé.} Deux cas d'impossibilité :
\begin{enumerate}
\item De \textbf{deux prémisses négatives}, on ne peut rien conclure : deux négations séparent les termes au lieu de les relier.
\item De \textbf{deux prémisses particulières} (« quelques », « certains »), on ne peut rien conclure : aucun terme n'est alors pris dans toute son extension pour assurer la liaison.
\end{enumerate}

\begin{attention}[Erreur fréquente : conclure à partir de deux négations]
« Les chèvres ne sont pas des bovins ; or cet animal n'est pas une chèvre ; donc cet animal est un bovin. » Faux raisonnement : l'animal pourrait être un mouton, un porc, un poulet... De deux prémisses négatives, \emph{rien ne suit}. Dire ce qu'une chose n'est pas ne dit jamais ce qu'elle est.
\end{attention}

De même : « Certains élèves sont ponctuels ; or certains ponctuels sont récompensés ; donc certains élèves sont récompensés » est invalide : les « certains » des deux prémisses peuvent ne désigner personne en commun.

\courssub{Règle 5 : la conclusion suit toujours la prémisse la plus faible}

\textbf{Intuition.} Une chaîne n'est jamais plus solide que son maillon le plus faible. La conclusion d'un raisonnement ne peut pas prétendre davantage que ce que les prémisses autorisent.

\textbf{Énoncé.} Si l'une des prémisses est \cle{particulière}, la conclusion doit être particulière ; si l'une des prémisses est \cle{négative}, la conclusion doit être négative.

\begin{exemplebox}[Exemple chiffré : la conclusion probable]
« 80 \% des élèves interrogés réussissent quand ils révisent ; or Jean révise ; donc Jean a de fortes chances de réussir. » Ici la première prémisse est statistique, donc particulière : la conclusion ne peut être que \emph{probable}, non certaine. Affirmer « donc Jean réussira sûrement » serait violer la règle 5 : on conclurait plus que ce que les prémisses permettent.
\end{exemplebox}

\courssub{Règle 6 : rigueur du vocabulaire}

\textbf{Énoncé.} Un même terme doit garder le \emph{même sens} du début à la fin du raisonnement. Sinon, on commet le sophisme d'\cle{équivocation} : on joue sur deux sens d'un mot pour faire croire à une conclusion.

\begin{exemplebox}[Le mot « libre »]
« Un pays libre ne connaît pas la contrainte ; or le Cameroun est un pays libre ; donc au Cameroun aucune loi ne contraint les citoyens. » Le mot « libre » change de sens : il signifie d'abord « indépendant » (souveraineté), puis « sans aucune règle » (anarchie). Le raisonnement est invalide par équivocation.
\end{exemplebox}

\courssub{Tableau récapitulatif des six règles}

\begin{popfigure}
\begin{center}
\begin{tikzpicture}
\node[anchor=north west, font=\small, text width=\linewidth, inner sep=4pt] {
\begin{tabular}{|p{0.15\linewidth}|p{0.4\linewidth}|p{0.4\linewidth}|}
\hline
\rowcolor{popblueL} \textbf{Règle} & \textbf{Exemple valide} & \textbf{Contre-exemple (invalide)} \\
\hline
1. Prémisses vraies & Tous les élèves présents ont vu le corrigé ; Amina est présente ; donc elle l'a vu. & Tous les poissons volent ; le capitaine est un poisson ; donc il vole (prémisse fausse). \\
\hline
2. Principes logiques & Ce cahier est rouge ; il ne peut pas être en même temps non rouge. & « Je suis d'accord et pas d'accord en même temps sur le même point. » \\
\hline
3. Moyen terme distribué & Tous les hommes sont mortels ; Socrate est un homme ; donc Socrate est mortel. & Tous les voleurs sont des hommes ; Paul est un homme ; donc Paul est un voleur. \\
\hline
4. Pas deux négatives ni deux particulières & Aucun menteur n'est fiable ; or cet homme est un menteur ; donc il n'est pas fiable. & Les chèvres ne sont pas des bovins ; cet animal n'est pas une chèvre ; donc c'est un bovin. \\
\hline
5. Conclusion la plus faible & Certains élèves révisent ; or tous ceux qui révisent progressent ; donc certains élèves progressent. & 80 \% réussissent en révisant ; Jean révise ; donc Jean réussira sûrement. \\
\hline
6. Même sens des termes & « Droit » au sens de loi partout dans le raisonnement. & Jouer sur « libre » (indépendant / sans règles) pour conclure. \\
\hline
\end{tabular}
};
\end{tikzpicture}
\end{center}
\legende{Les six règles du raisonnement correct, avec pour chacune un exemple valide et un contre-exemple à reconnaître à l'examen.}
\end{popfigure}

\courssub{Application guidée : tester la validité de quatre syllogismes}

\begin{methode}[Fiche de contrôle d'un syllogisme en cinq questions]
Pour vérifier un syllogisme, pose-toi dans l'ordre :
\begin{enumerate}
\item \textbf{Termes} : y a-t-il exactement trois termes, et gardent-ils le même sens (règle 6) ?
\item \textbf{Négations} : les deux prémisses sont-elles négatives ? Si oui, rien ne suit (règle 4).
\item \textbf{Particularité} : les deux prémisses sont-elles particulières ? Si oui, rien ne suit (règle 4).
\item \textbf{Distribution} : le moyen terme est-il pris au moins une fois dans toute son extension (règle 3) ?
\item \textbf{Conclusion} : suit-elle la prémisse la plus faible, sans conclure plus que les prémisses (règle 5) ? Les prémisses sont-elles vraies ou admises (règle 1) ?
\end{enumerate}
\end{methode}

\begin{exoresolu}[Quatre syllogismes de la vie courante]
Teste la validité des quatre syllogismes suivants à l'aide de la fiche de contrôle.
\begin{enumerate}
\item \emph{Marché} : « Toutes les marchandes de ce quartier vendent des légumes frais ; or Mama Rose est marchande dans ce quartier ; donc elle vend des légumes frais. »
\item \emph{Stade} : « Tous les supporters des Lions Indomptables aiment le football ; or ce jeune aime le football ; donc il est supporter des Lions. »
\item \emph{Réseaux sociaux} : « Aucune information vérifiée n'est un mensonge ; or cette publication n'est pas une information vérifiée ; donc c'est un mensonge. »
\item \emph{Salle de classe} : « Certains élèves de la classe sont en retard ; or certains élèves en retard seront sanctionnés ; donc certains élèves de la classe seront sanctionnés. »
\end{enumerate}
\tcblower
\textbf{Solutions détaillées.}
\begin{enumerate}
\item \emph{Marché} : trois termes, même sens ; prémisse majeure universelle ; le moyen terme « marchandes de ce quartier » est distribué (« toutes ») ; conclusion universelle comme les prémisses. \textbf{Valide} (si la majeure est vraie, ce qu'il faut vérifier).
\item \emph{Stade} : moyen terme « ceux qui aiment le football » jamais distribué : on dit seulement que les supporters sont \emph{parmi} les amateurs de football. Le jeune peut aimer le football sans supporter les Lions. \textbf{Invalide} (violation de la règle 3).
\item \emph{Réseaux sociaux} : deux prémisses négatives (« aucune... n'est », « n'est pas »). Une publication non vérifiée peut être vraie sans être encore vérifiée. \textbf{Invalide} (violation de la règle 4) : de deux négations, rien ne suit.
\item \emph{Salle de classe} : deux prémisses particulières (« certains ») ; les élèves en retard sanctionnés peuvent n'être aucun de ceux de la classe. \textbf{Invalide} (violation de la règle 4).
\end{enumerate}
\end{exoresolu}

\begin{aretenir}[L'essentiel]
Un raisonnement correct exige la vérité des prémisses et la validité de la forme. Six règles garantissent la validité : prémisses vraies ou admises ; respect de l'identité, de la non-contradiction et du tiers exclu ; moyen terme distribué au moins une fois ; impossibilité de conclure de deux prémisses négatives ou particulières ; conclusion toujours aussi faible que la prémisse la plus faible ; constance du sens des termes. La fiche de contrôle en cinq questions permet de tester n'importe quel syllogisme, y compris à l'examen.
\end{aretenir}

\begin{exercice}[Pour s'entraîner]
Construis un syllogisme valide sur la vie de ton établissement scolaire, puis un sophisme volontaire qui viole l'une des six règles ; échange avec ton voisin et identifie la règle violée dans son sophisme.
\end{exercice}

\begin{corrige}[Exercice]
Exemple de syllogisme valide : « Tous les élèves inscrits au Probatoire technique passent une épreuve pratique ; or les élèves de notre classe sont inscrits au Probatoire technique ; donc ils passeront une épreuve pratique. » Exemple de sophisme : « Tous les grands footballeurs sont célèbres ; or mon frère est célèbre dans le quartier ; donc c'est un grand footballeur » : violation de la règle 3, le moyen terme « célèbre » n'est jamais distribué.
\end{corrige}

\coursec{Les erreurs de raisonnement : paralogismes et sophismes}

Dans la section précédente, nous avons établi les \cle{règles} qui rendent un raisonnement correct : validité formelle, vérité des prémisses, pertinence des arguments. Mais savoir ce qu'est un bon raisonnement ne suffit pas : il faut aussi apprendre à reconnaître les \emph{mauvais} raisonnements, car ils abondent dans la vie quotidienne — débats télévisés, réunions de quartier, publicités, messages WhatsApp. Or un raisonnement faux peut être séduisant, éloquent, émotionnel, et pourtant conduire à une conclusion inacceptable. Cette section nous donne les outils pour démasquer ces pièges.

\courssub{Paralogisme et sophisme : deux sortes d'erreurs}

\begin{definition}[Paralogisme et sophisme]
Un \cle{paralogisme} est une erreur de raisonnement commise \emph{involontairement}, de bonne foi : celui qui raisonne se trompe sans le vouloir, par inattention, précipitation ou manque de rigueur. Un \cle{sophisme} est un raisonnement faux présenté avec l'\emph{intention de tromper} : celui qui l'utilise sait (ou devine) que son argument est invalide, mais il cherche à convaincre ou à gagner le débat à tout prix.
\end{definition}

L'intuition est simple : imagine un élève qui, pressé, conclut « tous les élèves de la classe aiment le football, car Jean et Paul l'aiment ». Il se trompe de bonne foi : c'est un paralogisme (généralisation hâtive). Imagine maintenant un vendeur qui sait pertinemment que son produit est médiocre, mais qui affirme : « Une star célèbre l'utilise, donc il est excellent » pour te le vendre : c'est un sophisme (argument d'autorité mal employé). La \emph{forme} de l'erreur peut être identique ; c'est l'\emph{intention} qui distingue le paralogisme du sophisme.

\begin{popfigure}
\begin{center}
\begin{tikzpicture}[
  grow=down,
  level 1/.style={sibling distance=7cm, level distance=1.6cm},
  level 2/.style={sibling distance=3.6cm, level distance=1.5cm},
  every node/.style={draw, rounded corners=2pt, align=center, font=\small, inner sep=3pt},
  edge from parent/.style={draw=popmuted, thick}
]
\node[fill=popblueL, font=\small\bfseries] {Erreurs de raisonnement}
  child { node[fill=popgreenL, align=center]{\textbf{Paralogismes}\\ (involontaires, bonne foi)}
    child { node[fill=white, align=center]{Paralogismes formels\\ (erreur de structure)} }
    child { node[fill=white, align=center]{Paralogismes informels\\ (précipitation, émotion)} }
  }
  child { node[fill=poporangeL, align=center]{\textbf{Sophismes}\\ (volontaires, pour tromper)}
    child { node[fill=white, align=center]{Homme de paille\\ Argument d'autorité\\ Pente glissante} }
    child { node[fill=white, align=center]{Attaque personnelle\\ Faux dilemme\\ Pétition de principe} }
  };
\end{tikzpicture}
\end{center}
\legende{Classification des erreurs de raisonnement. Le critère de partage est l'intention : le paralogisme est une faute commise de bonne foi, le sophisme est une manœuvre destinée à tromper l'auditoire. On distingue aussi les paralogismes formels (structure invalide) des paralogismes informels (contenu trompeur par mégarde).}
\end{popfigure}

\begin{savaistu}[Les Sophistes de l'Antiquité]
Dans l'Athènes du \textsc{v}\textsuperscript{e} siècle avant J.-C., les \emph{Sophistes} (Protagoras, Gorgias) étaient des maîtres payés pour enseigner l'art de persuader, c'est-à-dire de « rendre forte la cause faible ». Platon et Aristote les ont vivement combattus : pour eux, la philosophie cherche la \emph{vérité}, tandis que la sophistique ne cherche que la \emph{victoire} dans le débat. C'est d'eux que vient le mot « sophisme ». Aristote a répertorié les premiers sophismes dans ses \emph{Réfutations sophistiques}.
\end{savaistu}

\courssub{Les paralogismes formels : erreurs de structure}

\begin{definition}[Paralogismes formels courants]
Un paralogisme formel viole les règles de la logique déductive (syllogisme). Les deux plus classiques sont :
\begin{itemize}
\item \textbf{L'affirmation du conséquent} : « Si P alors Q. Q est vrai. Donc P est vrai. » (Ex. : « Si il pleut, le sol est mouillé. Le sol est mouillé. Donc il a plu. » — Faux : un arrosage peut mouiller le sol).
\item \textbf{La négation de l'antécédent} : « Si P alors Q. P est faux. Donc Q est faux. » (Ex. : « Si je travaille, je réussis. Je ne travaille pas. Donc je ne réussis pas. » — Faux : on peut avoir de la chance ou du talent).
\end{itemize}
Ces erreurs transforment une implication ($P \Rightarrow Q$) en équivalence ($P \Leftrightarrow Q$) sans justification.
\end{definition}

\begin{exemplebox}[Paralogisme formel dans la vie courante]
« Si ce médicament est efficace, le patient guérit. Le patient a guéri. Donc le médicament est efficace. » C'est une \emph{affirmation du conséquent} : la guérison peut avoir une autre cause (immunité naturelle, effet placebo, autre traitement). Ce paralogisme est fréquent dans les témoignages de produits « miracles ».
\end{exemplebox}

\courssub{Les principaux sophismes (erreurs intentionnelles)}

\textbf{1. L'homme de paille.}\par
Le sophisme de l'\cle{homme de paille} consiste à \emph{déformer} la thèse de l'adversaire pour la réfuter plus facilement : au lieu d'attaquer l'argument réel, on fabrique un épouvantail (un « homme de paille ») que l'on renverse sans effort. Exemple dans un débat politique télévisé : un candidat propose « réduire les dépenses de cérémonies officielles pour financer les écoles » ; son adversaire réplique : « Mon concurrent veut détruire les traditions du Cameroun ! » La proposition réelle (réduire certaines dépenses) a été remplacée par une caricature (détruire les traditions), beaucoup plus facile à condamner.

\begin{popfigure}
\begin{center}
\begin{tikzpicture}[
  case/.style={draw=popdark, thick, rounded corners=3pt, fill=popcream, minimum width=4.6cm, minimum height=3.4cm, align=center},
  bulle/.style={draw=popblue, fill=white, rounded corners=6pt, align=center, font=\footnotesize, inner sep=4pt}
]
% Case 1
\node[case] (c1) at (0,0) {};
\node[font=\footnotesize\bfseries, text=popdark] at (0,1.45) {Case 1 : la thèse réelle};
\node[circle, draw=popdark, fill=popblueL, minimum size=0.5cm] (a1) at (-1.4,-0.9) {};
\node[bulle, text width=2.6cm] at (0.6,0.1) {« Je propose de réduire les heures de cours du samedi. »};
% Case 2
\node[case] (c2) at (5.4,0) {};
\node[font=\footnotesize\bfseries, text=popdark] at (5.4,1.45) {Case 2 : la déformation};
\node[circle, draw=popdark, fill=poporangeL, minimum size=0.5cm] (b2) at (6.8,-0.9) {};
\node[bulle, draw=poporange, text width=2.6cm] at (4.8,0.1) {« Tu veux donc supprimer l'école et que les élèves restent ignorants ! »};
% Case 3
\node[case] (c3) at (10.8,0) {};
\node[font=\footnotesize\bfseries, text=popdark] at (10.8,1.45) {Case 3 : la fausse victoire};
\node[circle, draw=popdark, fill=poporangeL, minimum size=0.5cm] (b3) at (12.2,-0.9) {};
\node[bulle, draw=poporange, text width=2.6cm] at (10.2,0.1) {« Voyez ! Il veut l'ignorance des enfants. J'ai gagné le débat ! »};
\end{tikzpicture}
\end{center}
\legende{L'homme de paille en trois temps : (1) un élève émet une proposition précise ; (2) son camarade la remplace par une caricature exagérée ; (3) il réfute la caricature et croit avoir réfuté la thèse réelle.}
\end{popfigure}

\textbf{2. L'argument d'autorité mal employé.}\par
Invoquer une autorité \emph{compétente dans son domaine} est légitime (un médecin sur la santé, un ingénieur sur un pont). Le sophisme consiste à invoquer une autorité \emph{hors de son domaine} : « Un footballeur célèbre dit que cette boisson donne de l'intelligence, donc c'est vrai. » La célébrité du joueur ne prouve rien sur la boisson. Ce procédé domine la publicité et les réseaux sociaux, où des vedettes font la promotion de produits qu'elles ne connaissent pas.

\textbf{3. La pente glissante.}\par
La \cle{pente glissante} enchaîne des conséquences de plus en plus catastrophiques \emph{sans prouver} que chaque étape entraîne la suivante : « Si on autorise les élèves à poser des questions au professeur, bientôt ils ne respecteront plus personne, puis l'école deviendra un champ de bataille, puis le pays s'effondrera. » Chaque maillon de la chaîne est présenté comme évident alors qu'aucun n'est démontré.

\textbf{4. L'attaque personnelle (ad hominem).}\par
L'argument \cle{ad hominem} (« contre l'homme ») répond à la \emph{personne} au lieu de répondre à ses \emph{arguments}.

\begin{exemplebox}[Attaque personnelle]
« Tu critiques la gestion de la commune, mais toi-même tu jettes des papiers par terre. » Le fait que le critique jette des papiers ne rend pas sa critique fausse : la gestion de la commune peut être mauvaise \emph{même si} celui qui le dit est lui-même malpropre. On a attaqué la personne, pas l'argument.
\end{exemplebox}

\textbf{5. La fausse généralisation et le faux dilemme.}\par
La \cle{fausse généralisation} tire une conclusion universelle de quelques cas : « J'ai vu deux conducteurs de moto imprudents, donc tous les motards sont des dangers publics. » Le \cle{faux dilemme} réduit abusivement la situation à deux options alors qu'il en existe d'autres : « Soit tu es avec nous, soit tu es contre nous » — alors qu'on peut être partiellement d'accord, neutre, ou favorable à une troisième solution.

\textbf{6. Le cercle vicieux (pétition de principe).}\par
La \cle{pétition de principe} prend pour prémisse ce qu'on prétend démontrer : « Ce médium est fiable parce que ses prédictions sont exactes ; et ses prédictions sont exactes parce qu'il est fiable. » Le raisonnement tourne en rond : la conclusion était déjà cachée dans les prémisses.

\courssub{Méthode de détection et pièges à éviter}

\begin{methode}[Démasquer un sophisme en trois étapes]
\begin{enumerate}
\item \textbf{Reformuler} le raisonnement sous forme standard : identifier clairement les prémisses et la conclusion (« P1, P2, donc C »).
\item \textbf{Appliquer les règles de validité} : la conclusion découle-t-elle des prémisses ? Les prémisses sont-elles vraies et pertinentes ?
\item \textbf{Nommer l'erreur} : si le lien entre prémisses et conclusion repose sur une déformation, une émotion, une attaque ou une autorité hors sujet, identifier le sophisme précis commis.
\end{enumerate}
\end{methode}

\begin{attention}[L'éloquence n'est pas la validité]
Erreur fréquente : croire qu'un raisonnement \emph{bien présenté} — éloquent, émotionnel, applaudi par la foule — est forcément valide. Un sophisme est précisément un raisonnement \emph{invalide mais persuasif}. Au Probatoire et au Baccalauréat technique, le correcteur sanctionne les sophismes même s'ils sont joliment écrits : juge toujours la \emph{structure} de l'argument, jamais son emballage.
\end{attention}

\begin{exoresolu}[Identifier un sophisme]
Énoncé. Un candidat déclare à la télévision : « Mon adversaire affirme qu'il faut moderniser l'agriculture. Vous l'avez entendu : il méprise nos ancêtres et veut vendre nos terres aux étrangers ! » Nommer le sophisme et le réfuter.
\tcblower
Solution. Étape 1 — Reformulation : prémisse réelle de l'adversaire = « il faut moderniser l'agriculture » ; conclusion du candidat = « il méprise nos ancêtres et veut vendre nos terres ». Étape 2 — Test de validité : moderniser (améliorer les techniques) ne signifie ni mépriser les ancêtres ni vendre les terres ; le lien n'est pas démontré. Étape 3 — Nom : c'est un \textbf{homme de paille} (la thèse réelle est remplacée par une caricature), doublé d'une \textbf{pente glissante} (moderniser $\rightarrow$ vendre les terres, sans preuve). Réfutation : il faut répondre à la thèse réelle — discuter les avantages et les coûts de la modernisation — et non à la caricature.
\end{exoresolu}

\courssub{Travaux pratiques : décrypter des messages réels}

\begin{experience}[TP guidé : la chasse aux sophismes]
\textbf{Matériel} : 5 messages publicitaires ou posts WhatsApp circulant au Cameroun (relevés par les élèves ou fournis par le professeur). \textbf{Protocole} : pour chaque message, appliquer la méthode en trois étapes (reformuler, tester, nommer). \textbf{Exemples de messages à analyser} :
\begin{enumerate}
\item « Un célèbre artiste musicien recommande cette tisane : elle guérit tout ! » $\rightarrow$ argument d'autorité mal employé.
\item « Si tu ne partages pas ce message, tu n'aimes pas le Cameroun. » $\rightarrow$ faux dilemme (et pression émotionnelle).
\item « Ce voyant est puissant car tout le monde dit qu'il est puissant. » $\rightarrow$ pétition de principe.
\item « Mon voisin a acheté ce produit et il est tombé malade : ce produit rend malade. » $\rightarrow$ fausse généralisation (un seul cas, sans lien prouvé).
\item « Tu critiques mon projet ? Toi qui as redoublé, tais-toi ! » $\rightarrow$ attaque personnelle (ad hominem).
\end{enumerate}
\textbf{Interprétation} : rédiger pour chaque message une phrase de réfutation qui répond à l'argument réel et non au piège.
\end{experience}

\begin{savaistu}[Sophismes et théories du complot]
Les théories du complot qui circulent sur les réseaux sociaux combinent souvent plusieurs sophismes à la fois : le \emph{faux dilemme} (« soit tu crois la version officielle, soit tu es éveillé »), l'\emph{argument d'autorité} (« un docteur très connu l'a dit dans une vidéo ») et la \emph{pente glissante} (« si on accepte ceci, bientôt on contrôlera toutes nos vies »). Savoir nommer ces procédés est une protection concrète du citoyen.
\end{savaistu}

\begin{aretenir}[L'essentiel]
Un \cle{paralogisme} est une erreur de raisonnement involontaire (paralogisme formel : structure invalide ; paralogisme informel : précipitation, émotion) ; un \cle{sophisme} est un raisonnement fallacieux volontaire, destiné à tromper. Les sophismes principaux sont : l'homme de paille, l'argument d'autorité mal employé, la pente glissante, l'attaque personnelle, la fausse généralisation, le faux dilemme et la pétition de principe. Pour les démasquer : reformuler le raisonnement sous forme standard, tester sa validité, puis nommer l'erreur. Un discours séduisant n'est jamais, par cela seul, un discours vrai.
\end{aretenir}

\coursec{Raisonner dans la vie courante : méthode et rédaction}

Savoir repérer les paralogismes et les sophismes, comme nous l'avons appris dans la section précédente, c'est savoir \emph{déjouer} les mauvais raisonnements. Mais l'esprit critique ne consiste pas seulement à détruire : il faut aussi savoir \emph{construire} un raisonnement correct, l'exposer clairement et le justifier. C'est ce que nous apprend cette dernière section : comment passer de la théorie du raisonnement à sa pratique effective, dans la vie quotidienne, dans les métiers techniques et à l'examen.

\courssub{Construire un argumentaire : thèse, arguments, exemples, conclusion}

\begin{definition}[Argumentaire]
Un \cle{argumentaire} est un ensemble organisé d'énoncés destiné à faire admettre une \cle{thèse} (la conclusion) à partir d'\cle{arguments} (les prémisses), appuyés si possible par des \cle{exemples} qui les rendent concrets et vérifiables.
\end{definition}

Tout argumentaire bien construit comporte quatre éléments :

\begin{enumerate}
\item \textbf{La thèse} : ce que je veux établir ou défendre. Elle doit être formulée clairement, en une phrase.
\item \textbf{Les arguments (prémisses)} : les raisons que je donne d'accepter la thèse. Ils doivent être vrais (ou admis par mon interlocuteur) et pertinents.
\item \textbf{Les exemples} : des cas concrets qui illustrent et renforcent les arguments.
\item \textbf{La conclusion} : elle reprend la thèse, désormais justifiée, avec la force qui correspond au raisonnement employé.
\end{enumerate}

\begin{exemplebox}[Un argumentaire complet]
Thèse : « Il faut réparer le pont de la Wouri avant la saison des pluies. »
Arguments : (1) Les pluies diluviennes fragilisent les structures déjà fissurées ; (2) le pont présente des fissures visibles depuis mars ; (3) des milliers de camions l'empruntent chaque jour.
Exemple : l'effondrement d'un ouvrage similaire dans une autre ville après des pluies exceptionnelles.
Conclusion : donc, reporter les travaux expose les usagers à un danger réel ; il faut agir maintenant.
\end{exemplebox}

\begin{methode}[Rédiger une réponse argumentée en 4 temps]
\begin{enumerate}
\item \textbf{J'expose le problème} : je reformule la question posée et j'annonce ma thèse.
\item \textbf{Je pose mes prémisses} : je présente mes arguments, du plus simple au plus décisif, en les appuyant d'exemples.
\item \textbf{Je déduis ma conclusion} : je montre comment elle découle des prémisses (« donc », « par conséquent », « il s'ensuit que »).
\item \textbf{Je vérifie la validité} : je relis en me demandant : mes prémisses sont-elles vraies ? Ma conclusion est-elle bien contenue dans mes prémisses ? N'ai-je commis aucun sophisme ?
\end{enumerate}
\end{methode}

\courssub{Justifier une démarche : nommer le raisonnement et montrer qu'il respecte les règles}

Justifier une démarche, ce n'est pas seulement donner un résultat : c'est \emph{montrer le chemin parcouru}. Une bonne justification comporte deux gestes :

\begin{enumerate}
\item \textbf{Annoncer le type de raisonnement utilisé} : « Je raisonne ici par déduction », « Mon raisonnement est inductif », « Je procède par analogie ».
\item \textbf{Montrer que les règles sont respectées} : pour une déduction, vérifier que le moyen terme est bien pris universellement au moins une fois et que la conclusion ne contient rien de plus que les prémisses ; pour une induction, préciser que la généralisation repose sur des observations nombreuses et variées ; pour une analogie, indiquer sur quelles ressemblances pertinentes elle s'appuie.
\end{enumerate}

\begin{exoresolu}[Justifier une démarche]
Énoncé. Un élève affirme : « Tous les élèves sérieux révisent régulièrement. Or Amina révise régulièrement. Donc Amina est une élève sérieuse. » Identifiez le type de raisonnement et dites s'il est valide.
\tcblower
Solution. Il s'agit d'un syllogisme déductif : deux prémisses et une conclusion. Mais il est \emph{invalide} : le moyen terme « réviser régulièrement » n'est jamais pris universellement (dans la majeure, il est attribut ; dans la mineure, il est encore attribut). Rien n'interdit qu'Amina révise sans être sérieuse, ou que d'autres révisent sans sérieux. La démarche est donc une déduction \emph{mal construite} : la forme est celle du syllogisme, mais la règle du moyen terme est violée. La conclusion peut être vraie par hasard, elle n'est pas démontrée.
\end{exoresolu}

\courssub{Rédiger une conclusion proportionnée aux prémisses}

\begin{propriete}[Proportionnalité de la conclusion]
Une conclusion doit toujours être \cle{proportionnée} à la force du raisonnement qui la produit : \textbf{certaine} en \cle{déduction} (si les prémisses sont vraies et la forme valide, la conclusion est nécessaire) ; \textbf{probable} en \cle{induction} (la généralisation reste révisable) ; \textbf{plausible} en \cle{analogie} (la ressemblance suggère, sans garantir).
\end{propriete}

Concrètement, cela se traduit dans la rédaction : après une déduction valide, on écrit « donc » ; après une induction, « il est donc probable que », « cela laisse penser que » ; après une analogie, « on peut supposer que », « il est plausible que ». Affirmer avec certitude ce qui n'est que probable est une faute de raisonnement.

\begin{attention}[Erreur fréquente à l'examen]
Beaucoup de candidats affirment leur conclusion sans montrer les étapes du raisonnement. Or le correcteur évalue la \emph{démarche}, pas seulement le résultat : une conclusion juste mal justifiée vaut moins qu'une démarche rigoureuse. Autre piège : conclure « donc » après une simple induction ou une analogie. Réservez le « donc » sec à la déduction.
\end{attention}

\courssub{Esprit critique face aux informations}

À l'ère des réseaux sociaux, des \emph{fake news} et des rumeurs de quartier, raisonner juste devient une hygiène quotidienne. Trois réflexes :

\begin{enumerate}
\item \textbf{Vérifier les sources} : qui parle ? Est-ce un témoin direct, un spécialiste, un inconnu anonyme ? L'information est-elle confirmée par plusieurs sources indépendantes ?
\item \textbf{Repérer les sophismes} : attaque contre la personne (\emph{ad hominem}), appel à la peur ou à la pitié, argument d'autorité abusive, généralisation hâtive à partir d'un seul cas.
\item \textbf{Exiger des preuves} : une affirmation extraordinaire exige des preuves solides. « On dit que\ldots » n'est pas une preuve.
\end{enumerate}

\begin{exemplebox}[La rumeur de quartier]
Une rumeur circule à Douala : « L'eau du quartier est empoisonnée, quelqu'un est tombé malade. » L'esprit critique demande : combien de malades ? Le lien avec l'eau est-il établi par une analyse ? Qui a lancé l'information ? Une seule personne malade ne prouve rien (généralisation hâtive) ; sans analyse de l'eau, l'accusation reste une rumeur.
\end{exemplebox}

\begin{savaistu}[Raisonner juste, une compétence démocratique]
L'esprit critique est la première défense du citoyen contre la manipulation. Une démocratie suppose des électeurs capables d'évaluer des discours, de démasquer les sophismes des démagogues et d'exiger des preuves. Apprendre à raisonner, c'est donc aussi apprendre à être libre.
\end{savaistu}

\courssub{Le raisonnement dans les métiers techniques : la démarche hypothético-déductive}

Les professionnels techniques raisonnent chaque jour selon la démarche \cle{hypothético-déductive} : face à un problème observé, ils formulent une hypothèse, en déduisent des conséquences vérifiables, testent, puis concluent.

\begin{popfigure}
\begin{center}
\begin{tikzpicture}[node distance=0.4cm,
etape/.style={draw=popblue, fill=popblueL, rounded corners, thick, text width=2.7cm, align=center, minimum height=1cm, font=\small},
fleche/.style={-{Stealth[length=3mm]}, thick, popdark}]
\node[etape] (obs) {\textbf{1. Observation}\par Un fait, un symptôme};
\node[etape, right=of obs] (hyp) {\textbf{2. Hypothèse}\par Explication possible};
\node[etape, right=of hyp] (ded) {\textbf{3. Déduction}\par « Si l'hypothèse est vraie, alors\ldots »};
\node[etape, right=of ded] (ver) {\textbf{4. Vérification}\par Test, mesure, examen};
\node[etape, right=of ver] (ccl) {\textbf{5. Conclusion}\par Hypothèse confirmée ou rejetée};
\draw[fleche] (obs) -- (hyp);
\draw[fleche] (hyp) -- (ded);
\draw[fleche] (ded) -- (ver);
\draw[fleche] (ver) -- (ccl);
\draw[fleche, poporange, dashed] (ccl.south) to[bend right=25] node[below, font=\scriptsize, text=poporange]{si rejetée : nouvelle hypothèse} (hyp.south);
\end{tikzpicture}
\end{center}
\legende{Les cinq étapes de la démarche hypothético-déductive. Si le test échoue, on formule une nouvelle hypothèse et on recommence.}
\end{popfigure}

\begin{exemplebox}[Le mécanicien de Bonabéri]
Un mécanicien de Bonabéri entend un sifflement sous le capot (observation). Il formule une hypothèse : la courroie est usée. Il en déduit une conséquence testable : « Si la courroie est usée, elle doit présenter des fissures et le sifflement doit cesser quand je la remplace. » Il vérifie : la courroie est fissurée ; après remplacement, le bruit disparaît. Conclusion : l'hypothèse est confirmée. Ce mécanicien \emph{raisonne} : il ne bricole pas au hasard, il suit une démarche logique.
\end{exemplebox}

La même démarche guide l'\textbf{infirmier} (fièvre + toux $\rightarrow$ hypothèse : infection $\rightarrow$ prise de température et examen $\rightarrow$ diagnostic), le \textbf{comptable} (écart de caisse $\rightarrow$ hypothèse : erreur de saisie $\rightarrow$ rapprochement des pièces $\rightarrow$ correction) et l'\textbf{agronome} (feuilles jaunies $\rightarrow$ hypothèse : carence en azote $\rightarrow$ analyse du sol $\rightarrow$ fertilisation adaptée).

\courssub{Dossier guidé : analyser un discours}

\begin{exoresolu}[Analyse d'un extrait de débat]
Énoncé. Lors d'un débat, un intervenant déclare : « Ce candidat est jeune ; or les jeunes sont tous impulsifs ; un impulsif ne peut pas diriger. D'ailleurs, ses adversaires disent du mal de lui, c'est bien la preuve qu'il est dangereux. Voter pour lui, c'est donc ruiner le pays. » Analysez ce discours : prémisses, conclusion, types de raisonnement, erreurs éventuelles.
\tcblower
Solution. \textbf{Prémisses} : (1) le candidat est jeune ; (2) « les jeunes sont tous impulsifs » ; (3) « un impulsif ne peut pas diriger » ; (4) « ses adversaires disent du mal de lui ». \textbf{Conclusion} : voter pour lui ruinerait le pays. \textbf{Types} : un syllogisme apparent (1+2+3) et un argument fallacieux d'adversaires (4). \textbf{Erreurs} : la prémisse (2) est une \emph{généralisation hâtive} (faux universel, paralogisme) ; (4) est un sophisme : que des adversaires critiquent un rival ne prouve rien (argument \emph{ad hominem} inversé / homme de paille). La conclusion est donc non démontrée. \textbf{Paragraphe de justification} : « Ce discours a la forme d'une déduction, mais il est invalide dans son contenu : sa majeure repose sur un stéréotype, non sur un fait universel ; il ajoute un sophisme d'autorité détournée. Sa conclusion, affirmée avec certitude, n'est donc même pas plausible : le locuteur abuse de la forme logique pour masquer l'absence de preuves. »
\end{exoresolu}

\courssub{Bilan de synthèse de la leçon}

\begin{popfigure}
\begin{center}
\begin{tikzpicture}[mindmap, grow cyclic,
every node/.style={concept, font=\small},
concept color=popblue,
level 1 concept/.append style={level distance=3.4cm, sibling angle=72, font=\small},
level 2 concept/.append style={level distance=2.2cm, font=\scriptsize}]
\node[concept, fill=popblue, text=white]{\textbf{Le raisonnement}}
child[concept color=poporange] { node[concept]{Syllogisme}
  child { node[concept]{majeure, mineure, conclusion} }
  child { node[concept]{moyen terme} }
}
child[concept color=popgreen] { node[concept]{Types}
  child { node[concept]{déduction} }
  child { node[concept]{induction} }
  child { node[concept]{analogie} }
}
child[concept color=poppurple] { node[concept]{Règles}
  child { node[concept]{moyen terme universel} }
  child { node[concept]{conclusion proportionnée} }
}
child[concept color=poppink] { node[concept]{Erreurs}
  child { node[concept]{paralogismes} }
  child { node[concept]{sophismes} }
}
child[concept color=popgold] { node[concept]{Vie courante}
  child { node[concept]{argumentaire} }
  child { node[concept]{esprit critique} }
};
\end{tikzpicture}
\end{center}
\legende{Carte mentale de synthèse : du raisonnement au syllogisme, aux types, aux règles, aux erreurs et à la pratique quotidienne.}
\end{popfigure}

\begin{aretenir}[L'essentiel]
Raisonner dans la vie courante, c'est construire un argumentaire (thèse, prémisses, exemples, conclusion), justifier sa démarche en nommant le type de raisonnement employé, et proportionner la conclusion à sa force : certaine en déduction, probable en induction, plausible en analogie. Face aux informations, l'esprit critique vérifie les sources, repère les sophismes et exige des preuves. Les métiers techniques appliquent la démarche hypothético-déductive : observer, supposer, déduire, vérifier, conclure.
\end{aretenir}

\begin{exercice}[Pour s'entraîner]
Votre téléphone ne s'allume plus. Rédigez en cinq phrases la démarche hypothético-déductive complète que vous suivriez, en nommant chaque étape, puis rédigez une conclusion correctement proportionnée.
\end{exercice}

\begin{corrige}[Exercice 1]
(1) Observation : l'écran reste noir malgré la pression du bouton. (2) Hypothèse : la batterie est vide. (3) Déduction : si la batterie est vide, le téléphone doit s'allumer après branchement au chargeur. (4) Vérification : je branche ; le témoin de charge apparaît et l'appareil démarre. (5) Conclusion : mon hypothèse est confirmée, la panne venait de la batterie déchargée. La conclusion est ici pratiquement certaine car le test décisif a réussi ; si le test avait échoué, j'aurais dû formuler une nouvelle hypothèse (chargeur défectueux, carte mère, etc.).
\end{corrige}

\newpage

\coursec{Activité d'intégration}

\courssub{Objectifs de l'activité}

Cette activité d'intégration mobilise l'ensemble des savoirs et savoir-faire de la leçon \emph{Le raisonnement} dans une situation authentique de la vie sociale camerounaise. À l'issue de l'activité, l'élève doit être capable de :

\begin{itemize}
\item identifier, dans un discours réel, les \cle{prémisses} et la \cle{conclusion} d'un argument ;
\item classer un raisonnement selon son type : \cle{déductif}, \cle{inductif} ou par \cle{analogie} ;
\item analyser la structure d'un \cle{syllogisme} (majeure, mineure, conclusion ; grand terme, moyen terme, petit terme) et tester sa \cle{validité} à l'aide des règles étudiées ;
\item repérer et nommer les \cle{sophismes} (homme de paille, généralisation hâtive, argument d'autorité) ;
\item rédiger et justifier une conclusion personnelle en construisant un syllogisme valide, conformément au savoir-faire « rédiger et justifier une démarche, une réponse ou une conclusion ».
\end{itemize}

\begin{methode}[Organisation de la séance]
\begin{enumerate}
\item Travail en groupes de 4 à 6 élèves (45 minutes) : analyse du dossier de 4 documents.
\item Rédaction de l'avis motivé (20 minutes).
\item Restitution orale sous forme de plaidoiries (2 minutes par groupe).
\item Correction collective au tableau à l'aide de la fiche de contrôle en 5 questions.
\end{enumerate}
\end{methode}

\courssub{Situation}

\textbf{Le tribunal des arguments : un débat de quartier à Yaoundé.}

Dans le quartier Nkol-Eton, à Yaoundé, le conseil de quartier débat de l'installation d'un \cle{forage d'eau} financé par la commune. Deux emplacements sont en concurrence : le \emph{terrain A}, situé près de l'école publique, et le \emph{terrain B}, situé derrière le marché. Tu es membre de la commission des jeunes chargée de rendre un avis motivé. Voici le dossier dont tu disposes.

\textbf{Document 1 — Retranscription du débat au conseil de quartier (extraits).}

\emph{M. Abena (président du conseil) :} « Tous les forages installés près des écoles réduisent les absences des élèves pour corvée d'eau. Or le terrain A est situé près d'une école. Donc un forage sur le terrain A réduira les absences des élèves. »

\emph{Mme Etoundi (commerçante) :} « Ceux qui défendent le terrain B veulent supprimer le marché et chasser les mamans de leurs étals. Nous ne pouvons pas accepter la destruction du marché ! Refusons donc le terrain B. »

\emph{M. Nkoulou (habitant) :} « Mon voisin a creusé un puits derrière sa concession et l'eau est sortie abondante. Le quartier voisin de Mvog-Ada a aussi trouvé de l'eau à faible profondeur. Donc l'eau jaillira forcément partout chez nous, peu importe l'emplacement. »

\textbf{Document 2 — Extrait d'un discours électoral.}

« J'ai visité deux villages de l'arrondissement et dans les deux, les chefs m'ont dit qu'ils manquaient d'eau potable. C'est la preuve que tous les villages de la région manquent d'eau potable. Votez pour moi : je suis le seul à connaître vos problèmes ! »

\textbf{Document 3 — Publicité diffusée sur WhatsApp.}

« EAU-VITALITÉ, la solution miracle ! Le célèbre Dr M., spécialiste reconnu, l'affirme : notre produit purifie toutes les eaux et guérit la fatigue. Un grand médecin ne peut pas se tromper. Commandez vite au 6XX XX XX XX ! »

\textbf{Document 4 — Données chiffrées : fréquentation du marché de Nkol-Eton.}

\begin{center}
\begin{tabularx}{\linewidth}{|l|X|X|}\hline
\rowcolor{popblueL} \textbf{Période} & \textbf{Moyenne de clients par jour} & \textbf{Recettes journalières moyennes (FCFA)} \\ \hline
Avant réhabilitation (2023) & 850 & 1\,200\,000 \\ \hline
Après réhabilitation (2024) & 1\,275 & 1\,920\,000 \\ \hline
\end{tabularx}
\end{center}

Le conseil municipal affirme : « La réhabilitation du marché a fait augmenter la fréquentation ; donc toute réhabilitation d'infrastructure augmente automatiquement la fréquentation, et le forage au terrain B attirera forcément plus d'usagers. »

\textbf{Problème posé à la commission des jeunes :} le forage doit-il être installé à l'emplacement proposé (terrain A) ? Ton avis devra distinguer les arguments solides des arguments trompeurs.

\courssub{Consignes}

\begin{enumerate}
\item \textbf{Relever.} Pour chacun des trois arguments du document 1, souligne les prémisses et encadre la conclusion. Rédige chaque argument sous la forme : « Prémisse 1 ; Prémisse 2 ; donc Conclusion ».
\item \textbf{Classer.} Indique, pour chaque raisonnement des documents 1, 2 et 4, s'il s'agit d'un raisonnement \cle{déductif}, \cle{inductif} ou par \cle{analogie}. Justifie ta réponse en une phrase.
\item \textbf{Tester.} L'argument de M. Abena est un syllogisme. Identifie le grand terme, le moyen terme et le petit terme, puis vérifie sa validité à l'aide des règles du syllogisme (nombre de termes, distribution du moyen terme, passage des prémisses à la conclusion). Refais le même travail avec l'argument de M. Nkoulou : est-il valide ?
\item \textbf{Nommer.} Repère le sophisme commis par Mme Etoundi (document 1), celui du discours électoral (document 2) et celui de la publicité (document 3). Nomme chacun précisément et explique en quoi il trompe.
\item \textbf{Exploiter les données.} Calcule le taux d'augmentation de la fréquentation du marché entre 2023 et 2024. Le raisonnement du conseil municipal (document 4) est-il valide ? Quel type de raisonnement utilise-t-il et quelle erreur commet-il ?
\item \textbf{Rédiger.} Rédige un avis motivé de 10 lignes répondant à la question : « Le forage doit-il être installé à l'emplacement proposé ? » Ton avis doit contenir \textbf{ton propre syllogisme valide}, clairement identifiable, et écarter au moins un sophisme du dossier.
\end{enumerate}

\courssub{Ressources à mobiliser}

\begin{aretenir}[Boîte à outils de la leçon]
\begin{itemize}
\item \textbf{Structure du syllogisme :} une \cle{majeure} (contient le grand terme), une \cle{mineure} (contient le petit terme), une \cle{conclusion} ; le \cle{moyen terme} relie les deux prémisses et disparaît de la conclusion.
\item \textbf{Règles de validité :} trois termes exactement ; le moyen terme doit être pris au moins une fois dans toute son extension ; un terme ne peut être universel dans la conclusion s'il est particulier dans les prémisses ; de deux prémisses négatives on ne tire rien ; de deux prémisses particulières on ne tire rien ; la conclusion suit la prémisse la plus faible.
\item \textbf{Types de raisonnement :} la \cle{déduction} va du général au particulier (conclusion nécessaire si les prémisses sont vraies) ; l'\cle{induction} va du particulier au général (conclusion seulement probable) ; l'\cle{analogie} va du particulier au particulier par ressemblance.
\item \textbf{Sophismes étudiés :} \cle{homme de paille} (déformer la thèse adverse pour la réfuter plus facilement), \cle{généralisation hâtive} (induction fondée sur trop peu de cas), \cle{argument d'autorité} (tenir pour vrai ce que dit un prestige sans preuve).
\item \textbf{Fiche de contrôle en 5 questions :} 1) Quelles sont les prémisses et la conclusion ? 2) Quel type de raisonnement ? 3) Les règles du syllogisme sont-elles respectées ? 4) Y a-t-il un sophisme, et lequel ? 5) La conclusion est-elle justifiée ?
\end{itemize}
\end{aretenir}

\courssub{Démarche et corrigé commenté}

\begin{exoresolu}[Corrigé commenté du tribunal des arguments]
\textbf{Étape 1 — Relever prémisses et conclusions (document 1).}

\emph{Argument de M. Abena :} P1 : « Tous les forages installés près des écoles réduisent les absences pour corvée d'eau. » P2 : « Le terrain A est situé près d'une école. » C : « Un forage sur le terrain A réduira les absences des élèves. »

\emph{Argument de Mme Etoundi :} P1 : « Les partisans du terrain B voudraient supprimer le marché » (thèse attribuée aux adversaires). P2 : « La destruction du marché est inacceptable. » C : « Il faut refuser le terrain B. »

\emph{Argument de M. Nkoulou :} P1 : « Mon voisin a trouvé de l'eau abondante. » P2 : « Le quartier voisin a trouvé de l'eau à faible profondeur. » C : « L'eau jaillira forcément partout chez nous. »

\tcblower

\textbf{Étape 2 — Classer les raisonnements.}

L'argument de M. Abena est \cle{déductif} : il part d'une loi générale (tous les forages près des écoles...) pour conclure sur un cas particulier (le terrain A). Celui de M. Nkoulou est \cle{inductif} : de deux cas particuliers, il prétend conclure à une loi générale (« partout »). Le discours électoral (document 2) est également inductif : deux villages visités, conclusion sur tous les villages de la région. Le raisonnement du document 4 est un raisonnement par \cle{analogie} (ou induction abusive) : d'un cas de réhabilitation réussi, on conclut que toute infrastructure réussira de même.

\textbf{Étape 3 — Tester les syllogismes.}

\emph{Syllogisme de M. Abena.} Grand terme : « réduire les absences » ; petit terme : « forage sur le terrain A » ; moyen terme : « forage installé près d'une école ». Il y a trois termes ; le moyen terme est pris dans toute son extension dans la majeure (« tous les forages... ») ; les deux prémisses sont affirmatives et universelles, la conclusion est affirmative. \textbf{Toutes les règles sont respectées : le syllogisme est valide.} Si les prémisses sont vraies, la conclusion est nécessaire. (On peut toutefois discuter la \emph{vérité} de la majeure : est-il établi que \emph{tous} ces forages réduisent les absences ? La validité formelle ne garantit pas la vérité matérielle.)

\emph{Raisonnement de M. Nkoulou.} Ce n'est pas un syllogisme mais une induction. Elle est \textbf{invalide} car elle conclut « forcément » à partir de deux cas seulement : c'est une \cle{généralisation hâtive}. Deux observations ne suffisent pas à établir une loi géologique ; seule une étude du sous-sol permettrait de conclure avec probabilité.

\textbf{Étape 4 — Nommer les sophismes.}

\emph{Mme Etoundi} commet un \cle{sophisme de l'homme de paille} : elle déforme la position adverse (installer le forage au terrain B) en une thèse extrême et fabriquée (« supprimer le marché »), facile à réfuter. Elle réfute un épouvantail, pas l'argument réel.

\emph{Le discours électoral} commet une \cle{généralisation hâtive} : deux villages ne représentent pas toute une région. Il ajoute un glissement vers l'argument \emph{ad hominem} (« je suis le seul à connaître vos problèmes »), qui substitue la personne à la preuve.

\emph{La publicité WhatsApp} repose sur un \cle{argument d'autorité} : le prestige du « Dr M. » remplace toute démonstration scientifique. Qu'un médecin l'affirme ne prouve ni l'efficacité du produit ni l'impossibilité de l'erreur (« un grand médecin ne peut pas se tromper » est une prémisse fausse).

\textbf{Étape 5 — Exploiter les données chiffrées.}

Taux d'augmentation de la fréquentation :
\[
\frac{1\,275-850}{850}\times 100 = \frac{425}{850}\times 100 = 50\,\%.
\]
La fréquentation a augmenté de 50~\% et les recettes de 60~\% ($\frac{720\,000}{1\,200\,000}\times 100$). Mais le conseil municipal commet une \cle{induction abusive} : d'\emph{un seul} cas (le marché), il conclut que \emph{toute} réhabilitation augmente la fréquentation, puis que le forage au terrain B attirera plus d'usagers. Un cas ne fait pas une loi ; de plus, un marché et un forage ne sont pas des infrastructures comparables (analogie faible).

\textbf{Étape 6 — Modèle d'avis motivé (10 lignes).}

« La commission des jeunes, après examen des arguments, émet un avis favorable à l'installation du forage sur le terrain A. Notre raisonnement est le suivant : tout équipement qui rapproche l'eau des élèves réduit les corvées et l'absentéisme scolaire ; or un forage près de l'école publique rapproche l'eau des élèves ; donc un forage sur le terrain A réduira les corvées et l'absentéisme. Ce syllogisme respecte les règles de validité : trois termes, moyen terme distribué, prémisses affirmatives. En revanche, nous écartons l'argument de Mme Etoundi, qui prête aux adversaires une intention de détruire le marché qu'ils n'ont jamais exprimée : c'est un homme de paille. Nous écartons également l'induction de M. Nkoulou : deux puits réussis ne prouvent pas la présence d'eau partout ; une étude technique préalable reste nécessaire. Enfin, l'argument tiré de la réhabilitation du marché, bien que les chiffres soient réels (+50~\% de fréquentation), ne peut être généralisé sans preuve. Nous recommandons donc le terrain A, sous réserve d'une étude du sous-sol. »
\end{exoresolu}

\begin{attention}[Pièges fréquents dans cette activité]
\begin{itemize}
\item Confondre \textbf{validité} (correction de la forme) et \textbf{vérité} (exactitude des prémisses) : un syllogisme peut être valide avec des prémisses fausses.
\item Croire qu'une induction peut conclure « forcément » : la conclusion inductive n'est jamais que \emph{probable}.
\item Nommer un sophisme sans l'expliquer : au Probatoire et au Baccalauréat technique, il faut toujours dire \emph{en quoi} le raisonnement trompe.
\end{itemize}
\end{attention}

\courssub{Bilan de l'activité}

Cette activité a permis de mettre en évidence trois acquis essentiels de la leçon. D'abord, les arguments de la vie courante — débats de quartier, discours politiques, publicités — peuvent tous être \emph{reconstruits} en prémisses et conclusions, ce qui rend leur examen rationnel possible. Ensuite, la distinction des trois types de raisonnement éclaire la force de chaque argument : seule la déduction correcte donne une conclusion nécessaire ; l'induction et l'analogie ne donnent que du probable, et deviennent des sophismes lorsqu'elles prétendent au nécessaire. Enfin, les règles du syllogisme et le répertoire des sophismes constituent une véritable \emph{fiche de contrôle} transférable : l'élève est désormais capable, face à n'importe quel discours, de poser les 5 questions de la fiche et de rédiger une conclusion justifiée — compétence directement exigible dans la dissertation au Probatoire et au Baccalauréat technique, et indispensable au citoyen face à la désinformation.

\newpage

\coursec{Méthodes et exercices résolus}

\coursec{Le raisonnement}

\courssub{Méthode 1 : Analyser un raisonnement dans une situation concrète}

\begin{methode}[Identifier et analyser un raisonnement]
Pour appliquer les notions de l'unité à une situation de la vie courante :
\begin{enumerate}
\item \textbf{Repérer le raisonnement} : chercher les articulations logiques (donc, car, puisque, par conséquent) qui relient des affirmations.
\item \textbf{Dégager la conclusion} : c'est la thèse défendue, ce que l'on veut démontrer.
\item \textbf{Identifier les prémisses} : ce sont les affirmations qui servent de point de départ et de justification.
\item \textbf{Classer le raisonnement} : est-il \cle{déductif} (du général au particulier), \cle{inductif} (du particulier au général) ou par \cle{analogie} (d'un cas semblable à un cas semblable) ?
\item \textbf{Évaluer la validité} : la conclusion découle-t-elle nécessairement des prémisses ? Les prémisses sont-elles vraies ?
\item \textbf{Conclure par une phrase} qui juge le raisonnement : valide ou non, et pourquoi.
\end{enumerate}
\textbf{Réflexe à retenir} : un raisonnement est \emph{valide} quand la forme est correcte (la conclusion découle nécessairement des prémisses) ; il est \emph{sain} (ou \emph{correct}) quand, de plus, les prémisses sont vraies.
\end{methode}

\begin{exoresolu}[Le raisonnement du commerçant de Mokolo]
\textbf{Énoncé.} Un commerçant du marché Mokolo à Yaoundé déclare : « Tous les clients qui achètent en gros bénéficient d'une réduction. Or Madame Fotso achète en gros. Donc elle bénéficiera d'une réduction. » Identifier les prémisses et la conclusion, préciser le type de raisonnement, puis dire s'il est valide.
\tcblower
\textbf{Solution détaillée.}
\begin{enumerate}
\item \textbf{Repérage} : le mot « donc » signale la conclusion ; les affirmations précédentes sont les prémisses.
\item \textbf{Prémisses} : (P1) « Tous les clients qui achètent en gros bénéficient d'une réduction » (prémisse générale, majeure) ; (P2) « Madame Fotso achète en gros » (prémisse particulière, mineure).
\item \textbf{Conclusion} : « Madame Fotso bénéficiera d'une réduction » (conclusion particulière).
\item \textbf{Type de raisonnement} : on passe d'une loi générale (P1) à un cas particulier (Madame Fotso) : c'est un raisonnement \cle{déductif}, et plus précisément un \cle{syllogisme}.
\item \textbf{Validité} : le terme moyen « acheter en gros » relie correctement les deux prémisses ; si P1 et P2 sont vraies, la conclusion est nécessairement vraie. Le raisonnement est donc \textbf{valide}. Si les prémisses sont exactes dans la réalité, il est aussi \textbf{sain}.
\end{enumerate}
\textbf{Conclusion} : le raisonnement du commerçant est un syllogisme déductif valide : la conclusion s'impose nécessairement à quiconque accepte les prémisses.
\end{exoresolu}

\courssub{Méthode 2 : Construire et vérifier un syllogisme}

\begin{methode}[Vérifier la structure d'un syllogisme]
Un \cle{syllogisme} est un raisonnement déductif composé de deux prémisses et d'une conclusion, comportant trois termes.
\begin{enumerate}
\item \textbf{Identifier les trois termes} : le \cle{grand terme} (attribut de la conclusion), le \cle{petit terme} (sujet de la conclusion), le \cle{terme moyen} (présent dans les deux prémisses, absent de la conclusion).
\item \textbf{Nommer les prémisses} : la \cle{majeure} contient le grand terme ; la \cle{mineure} contient le petit terme.
\item \textbf{Vérifier qu'il n'y a que trois termes} : quatre termes rendent le syllogisme invalide (équivoque ou quaternio terminorum).
\item \textbf{Vérifier le rôle du terme moyen} : il doit servir de lien entre le petit et le grand terme, puis disparaître de la conclusion. Il doit être pris au moins une fois dans toute son extension (être distribué).
\item \textbf{Vérifier le passage du général au particulier} : de deux prémisses particulières on ne peut rien conclure de nécessaire.
\item \textbf{Formuler le jugement} : syllogisme correct ou paralogisme.
\end{enumerate}
\textbf{Formule à retenir} : majeure + mineure $\Rightarrow$ conclusion ; 3 termes, ni plus ni moins.
\end{methode}

\begin{exoresolu}[Le syllogisme de l'élève de Première technique]
\textbf{Énoncé.} Un élève affirme : « Tous les élèves sérieux révisent leurs leçons. Or je révise mes leçons. Donc je suis un élève sérieux. » Construire l'analyse complète de ce syllogisme (termes, prémisses) et dire s'il est correct.
\tcblower
\textbf{Solution détaillée.}
\begin{enumerate}
\item \textbf{Les trois termes} :
\begin{itemize}
\item \textbf{Petit terme} (sujet de la conclusion) : « je » (l'élève qui parle).
\item \textbf{Grand terme} (attribut de la conclusion) : « élève sérieux ».
\item \textbf{Terme moyen} (présent dans les deux prémisses, absent de la conclusion) : « réviser ses leçons ».
\end{itemize}
\item \textbf{Majeure} (contient le grand terme) : « Tous les élèves sérieux révisent leurs leçons ». \textbf{Mineure} (contient le petit terme) : « Je révise mes leçons ».
\item \textbf{Test de structure} : il y a bien trois termes ; la forme est un syllogisme en 2ème figure (le terme moyen est attribut dans les deux prémisses).
\item \textbf{Test de validité} : la majeure dit que \emph{tous} les élèves sérieux révisent (le terme moyen « réviser » n'est pas distribué dans la majeure), mais elle ne dit pas que \emph{seuls} les élèves sérieux révisent. Un élève paresseux peut aussi réviser la veille du Probatoire par peur de l'échec. Le terme moyen n'est pas pris dans toute son extension dans la majeure : il ne sert pas de lien nécessaire entre « je » et « élève sérieux ». On ne peut donc pas conclure nécessairement.
\item \textbf{Jugement} : le raisonnement est un \cle{paralogisme} (erreur involontaire) : la conclusion est possible mais non nécessaire. Pour être valide, la majeure aurait dû être : « Seuls les élèves sérieux révisent leurs leçons » (ce qui distribue le terme moyen) ou « Tous ceux qui révisent leurs leçons sont des élèves sérieux » (conversion de la proposition).
\end{enumerate}
\textbf{Conclusion} : ce syllogisme est formellement incorrect : de « les élèves sérieux révisent » et « je révise », on ne peut pas déduire nécessairement « je suis sérieux ».
\end{exoresolu}

\courssub{Méthode 3 : Distinguer les types de raisonnement}

\begin{methode}[Reconnaître déduction, induction et analogie]
\begin{enumerate}
\item \textbf{Repérer le sens du passage} : du général vers le particulier = \cle{déduction} ; du particulier vers le général = \cle{induction} ; d'un cas particulier vers un cas particulier semblable = raisonnement par \cle{analogie}.
\item \textbf{Mesurer la force de la conclusion} : la déduction donne une conclusion \emph{nécessaire} (si prémisses vraies) ; l'induction et l'analogie donnent des conclusions \emph{probables} seulement.
\item \textbf{Vérifier l'induction} : le nombre de cas observés est-il suffisant ? Les cas sont-ils variés et représentatifs (induction complète vs incomplète) ?
\item \textbf{Vérifier l'analogie} : la ressemblance entre les deux cas porte-t-elle sur des points essentiels (causaux) ou seulement superficiels (accidentels) ?
\item \textbf{Nommer le raisonnement} dans la copie avec le vocabulaire exact du cours.
\end{enumerate}
\textbf{Réflexe} : « tous... donc celui-ci » = déduction ; « ceux-ci... donc tous » = induction ; « celui-ci ressemble à celui-là, donc... » = analogie.
\end{methode}

\begin{exoresolu}[Trois raisonnements à classer]
\textbf{Énoncé.} Classer chacun des raisonnements suivants et évaluer sa force. (a) « Tous les métaux se dilatent sous l'effet de la chaleur ; or le fer est un métal ; donc le fer se dilate sous l'effet de la chaleur. » (b) « J'ai observé que le maïs pousse bien dans les champs de Bafou, de Mbouda et de Foumban ; donc le maïs pousse bien dans tout l'Ouest camerounais. » (c) « Le quartier A de Douala, construit sur un terrain marécageux, a été inondé pendant la saison des pluies ; le quartier B est aussi construit sur un terrain marécageux ; il sera donc probablement inondé. »
\tcblower
\textbf{Solution détaillée.}
\begin{enumerate}
\item \textbf{Raisonnement (a)} : on part d'une loi générale (tous les métaux) pour aboutir à un cas particulier (le fer). C'est une \cle{déduction}, sous forme de syllogisme (Barbara : AAA-1). Le terme moyen « métal » relie correctement les prémisses : la conclusion est \textbf{nécessaire}. Raisonnement valide (et sain si les prémisses sont vraies).
\item \textbf{Raisonnement (b)} : on part de trois cas particuliers observés (Bafou, Mbouda, Foumban) pour conclure à une loi générale (tout l'Ouest). C'est une \cle{induction} (incomplète). Sa conclusion n'est que \textbf{probable} : trois localités ne représentent pas toute la région (diversité des sols, altitudes, microclimats) ; il faudrait multiplier et varier les observations pour renforcer la conclusion.
\item \textbf{Raisonnement (c)} : on part d'un cas particulier (quartier A) pour conclure à un autre cas particulier (quartier B) en s'appuyant sur une ressemblance (terrain marécageux). C'est un raisonnement par \cle{analogie}. La conclusion est \textbf{probable} : elle sera d'autant plus solide que la ressemblance porte sur un point essentiel (ici la nature du sol, cause efficiente des inondations) et non sur un détail accessoire.
\end{enumerate}
\textbf{Conclusion} : (a) est une déduction nécessaire ; (b) est une induction probable ; (c) est une analogie probable. Seule la déduction garantit sa conclusion par sa forme seule.
\end{exoresolu}

\courssub{Méthode 4 : Appliquer les règles du raisonnement correct}

\begin{methode}[Contrôler un raisonnement par les règles]
Pour rédiger et justifier une conclusion, soumettre tout raisonnement à quatre contrôles :
\begin{enumerate}
\item \textbf{Règle des termes} : les termes doivent garder le même sens du début à la fin (pas d'\cle{équivoque}) ; le syllogisme comporte exactement trois termes.
\item \textbf{Règle des prémisses} : les prémisses doivent être vraies et clairement établies ; on ne peut pas prouver une conclusion par elle-même (\cle{pétition de principe}).
\item \textbf{Règle de l'enchaînement} : la conclusion doit découler des prémisses ; de deux prémisses négatives ou de deux prémisses particulières, on ne peut rien conclure de nécessaire ; le terme moyen doit être distribué au moins une fois.
\item \textbf{Règle de la conclusion} : la conclusion ne doit rien affirmer de plus que ce que les prémisses contiennent ; elle suit toujours la partie la plus faible (si une prémisse est particulière, la conclusion est particulière ; si une prémisse est négative, la conclusion est négative).
\end{enumerate}
\textbf{Réflexe de rédaction} : annoncer les prémisses, expliciter l'enchaînement (« or », « donc »), puis formuler la conclusion sans l'élargir.
\end{methode}

\begin{exoresolu}[Contrôle d'un raisonnement fautif]
\textbf{Énoncé.} Un candidat au Probatoire technique écrit : « Aucun élève paresseux ne réussit au Probatoire. Or aucun élève de mon village n'est paresseux. Donc tous les élèves de mon village réussiront au Probatoire. » Appliquer les règles du raisonnement correct et dire si la conclusion est justifiée.
\tcblower
\textbf{Solution détaillée.}
\begin{enumerate}
\item \textbf{Règle des termes} : trois termes (« élève paresseux », « réussir au Probatoire », « élève de mon village ») : la règle est respectée, pas d'équivoque apparente.
\item \textbf{Règle des prémisses} : la première prémisse est discutable (certains élèves paresseux réussissent parfois par chance ou tricherie), mais admettons-la pour l'analyse formelle.
\item \textbf{Règle de l'enchaînement} : c'est ici que le raisonnement échoue formellement. Les deux prémisses sont \textbf{négatives} (« aucun... ne », « aucun... n'est »). Or de deux prémisses négatives on ne peut tirer \emph{aucune} conclusion nécessaire : le terme moyen « paresseux » exclut les deux autres termes sans jamais les relier positivement (il n'est distribué dans aucune prémisse de manière à unir le petit et le grand terme).
\item \textbf{Règle de la conclusion} : en outre, la conclusion est affirmative (A : « tous... réussissent ») alors que les prémisses sont négatives, et elle affirme plus que les prémisses ne permettent : ne pas être paresseux est une condition nécessaire mais non suffisante pour réussir (il faut aussi travailler, être bien encadré, etc.).
\end{enumerate}
\textbf{Conclusion} : le raisonnement viole la règle de l'enchaînement (deux prémisses négatives) et la règle de la conclusion (affirmative d'après négatives) : la conclusion n'est pas justifiée. Le candidat devait écrire, par exemple : « Tous les élèves travailleurs ont des chances de réussir ; or les élèves de mon village sont travailleurs ; donc ils ont des chances de réussir. »
\end{exoresolu}

\courssub{Méthode 5 : Démasquer les erreurs de raisonnement}

\begin{methode}[Identifier paralogismes et sophismes]
\begin{enumerate}
\item \textbf{Distinguer} : le \cle{paralogisme} est une erreur de raisonnement involontaire (faute de logique commise de bonne foi) ; le \cle{sophisme} est un raisonnement fallacieux construit \emph{volontairement} pour tromper ou persuader indûment.
\item \textbf{Repérer l'équivoque} : un même mot est-il employé avec deux sens différents (ex: « rien n'est éternel ; or l'éternité n'est rien ; donc l'éternité est éternelle ») ?
\item \textbf{Repérer la pétition de principe} : la conclusion est-elle déjà contenue, déguisée, dans les prémisses (raisonnement circulaire) ?
\item \textbf{Repérer la fausse généralisation (généralisation hâtive)} : conclut-on trop vite d'un ou deux cas à tous les cas (induction insuffisante) ?
\item \textbf{Repérer le faux raisonnement par analogie} : la comparaison porte-t-elle sur des ressemblances superficielles ou non pertinentes (ex: « l'horloge a un cadran, l'homme a un visage, donc l'homme est une horloge ») ?
\item \textbf{Corriger} : reformuler le raisonnement de manière valide et expliquer pourquoi la version initiale trompe.
\end{enumerate}
\end{methode}

\begin{exoresolu}[Le sophisme du vendeur de téléphones]
\textbf{Énoncé.} À Douala, un vendeur affirme à un client : « Ce téléphone est le meilleur, car aucun autre téléphone n'est meilleur que lui. » Identifier l'erreur de raisonnement commise, dire s'il s'agit d'un paralogisme ou d'un sophisme, et corriger.
\tcblower
\textbf{Solution détaillée.}
\begin{enumerate}
\item \textbf{Analyse} : la conclusion est « ce téléphone est le meilleur ». L'unique prémisse est « aucun autre téléphone n'est meilleur que lui ».
\item \textbf{Identification de l'erreur} : la prémisse ne fait que répéter la conclusion sous une forme négative : dire qu'aucun téléphone n'est meilleur que celui-ci, c'est déjà dire qu'il est le meilleur (au sens de « non surpassé »). On suppose donc dans la prémisse ce qu'on prétend démontrer : c'est une \cle{pétition de principe} (ou raisonnement circulaire).
\item \textbf{Paralogisme ou sophisme ?} : le vendeur cherche à convaincre le client pour vendre ; l'erreur paraît volontaire et destinée à tromper en donnant l'illusion d'un argument : c'est un \cle{sophisme}. (Si un élève commettait la même erreur de bonne foi dans une dissertation, ce serait un paralogisme.)
\item \textbf{Correction} : pour prouver honnêtement sa thèse, le vendeur devrait avancer des prémisses indépendantes et vérifiables : « Ce téléphone a une batterie de longue durée, un appareil photo performant et un prix inférieur aux concurrents ; or ce sont les critères d'un bon téléphone ; donc c'est un bon téléphone. » La conclusion reposerait alors sur des preuves réelles.
\end{enumerate}
\textbf{Conclusion} : le vendeur commet une pétition de principe, sophisme classique du discours publicitaire : il affirme au lieu de démontrer.
\end{exoresolu}

\courssub{Méthode 6 : Rédiger une argumentation complète}

\begin{methode}[Rédiger et justifier une démarche ou une conclusion]
Pour la dissertation ou l'exercice au Probatoire et au Baccalauréat technique :
\begin{enumerate}
\item \textbf{Annoncer la thèse} à démontrer en une phrase claire.
\item \textbf{Poser les prémisses} : une prémisse générale (la règle, la définition, la loi, le principe) puis une prémisse particulière (le cas concret, le fait observé).
\item \textbf{Expliciter l'enchaînement} avec les connecteurs logiques : « or », « donc », « par conséquent », « dès lors ».
\item \textbf{Formuler la conclusion} sans la dépasser : elle ne doit rien affirmer de plus que les prémisses ne contiennent implicitement.
\item \textbf{Justifier la démarche} : nommer le type de raisonnement utilisé (déduction, induction, analogie) et dire pourquoi il est valide ou seulement probable.
\item \textbf{Illustrer par un exemple concret} tiré de la vie courante camerounaise pour ancrer l'abstraction.
\end{enumerate}
\end{methode}

\begin{exoresolu}[Rédaction d'une argumentation sur la formation technique]
\textbf{Énoncé.} Sujet : « Un élève de Première technique soutient que la formation technique est utile à la société. » Rédiger en huit à dix lignes une argumentation rigoureuse défendant cette thèse, en justifiant la démarche employée.
\tcblower
\textbf{Solution détaillée (corrigé type).}

\emph{Toute formation qui répond aux besoins réels d'une société est utile à cette société. Or la formation technique répond aux besoins réels de la société camerounaise : elle forme les mécaniciens, électriciens, maçons et informaticiens dont le pays a besoin pour construire ses routes, ses usines et ses entreprises, comme on le voit dans les chantiers de Yaoundé et de Douala. Par conséquent, la formation technique est utile à la société camerounaise.}

\textbf{Justification de la démarche} : ce raisonnement est une \cle{déduction} construite en \cle{syllogisme} (mode Barbara). La majeure pose une règle générale (toute formation adaptée aux besoins est utile) ; la mineure constate un fait particulier vérifiable (la formation technique forme les professionnels dont le Cameroun a besoin) ; le terme moyen « répondre aux besoins réels de la société » relie les deux prémisses et disparaît de la conclusion. La conclusion ne dépasse pas les prémisses : elle est donc \textbf{nécessaire} et le raisonnement est \textbf{valide}. L'exemple des chantiers de Yaoundé et de Douala ancre l'argumentation dans la réalité concrète du pays, ce qui renforce la vérité de la prémisse particulière et rend le raisonnement \textbf{sain}.

\textbf{Conclusion} : en combinant une structure déductive correcte, des prémisses vraies et un exemple camerounais précis, on obtient une argumentation complète et justifiée, conforme aux exigences de l'examen.
\end{exoresolu}

\begin{attention}[Les 5 erreurs qui coûtent des points au Bac]
\begin{enumerate}
\item \textbf{Confondre validité et vérité (santé)} : un raisonnement peut être valide (forme correcte) avec une conclusion fausse si une prémisse est fausse. Toujours vérifier séparément la forme (validité) ET le contenu des prémisses (vérité) pour juger de la santé du raisonnement.
\item \textbf{Mal nommer le type de raisonnement} : écrire « déduction » pour un passage du particulier au général est une erreur grave. Retenir : général $\rightarrow$ particulier = déduction ; particulier $\rightarrow$ général = induction ; particulier $\rightarrow$ particulier = analogie.
\item \textbf{Oublier le terme moyen} : dans l'analyse d'un syllogisme, beaucoup de candidats ne désignent pas le terme moyen ni ne vérifient qu'il disparaît de la conclusion et qu'il est distribué au moins une fois. C'est pourtant la clé de la validité formelle.
\item \textbf{Confondre paralogisme et sophisme} : le paralogisme est une erreur \emph{involontaire} (incompétence, précipitation), le sophisme est une ruse \emph{volontaire} (intention de tromper). Préciser toujours l'intention dans la justification.
\item \textbf{Affirmer la conclusion sans la justifier} : poser « donc » ne suffit pas. Il faut expliciter l'enchaînement (prémisses, règle appliquée) et vérifier que la conclusion ne dépasse pas les prémisses, sous peine de commettre soi-même un paralogisme dans la copie.
\end{enumerate}
\end{attention}

\newpage

\coursec{Exercices}

\courssub{Je vérifie mes connaissances}

\begin{exercice}[QCM : les bases du raisonnement]
Pour chaque question, une seule réponse est exacte. Entoure la lettre correspondante.
\begin{enumerate}
\item Un raisonnement est :
\begin{enumerate}
\item[a)] une opinion exprimée avec force ;
\item[b)] un enchaînement de jugements dont l'un (la conclusion) découle des autres (les prémisses) ;
\item[c)] un souvenir que l'on raconte ;
\item[d)] une émotion ressentie devant un problème.
\end{enumerate}
\item Dans le syllogisme, le moyen terme est celui qui :
\begin{enumerate}
\item[a)] figure dans la conclusion ;
\item[b)] figure dans les deux prémisses mais pas dans la conclusion ;
\item[c)] figure dans les deux prémisses et dans la conclusion ;
\item[d)] ne figure nulle part.
\end{enumerate}
\item Le raisonnement qui va du général au particulier est :
\begin{enumerate}
\item[a)] l'induction ; \item[b)] l'analogie ; \item[c)] la déduction ; \item[d)] l'opinion.
\end{enumerate}
\item Un sophisme est :
\begin{enumerate}
\item[a)] un raisonnement valide à conclusion vraie ;
\item[b)] un raisonnement faux présenté volontairement comme valide pour tromper ;
\item[c)] une vérité évidente ;
\item[d)] une question sans réponse.
\end{enumerate}
\end{enumerate}
\end{exercice}

\begin{exercice}[Vrai ou faux ? Justifie]
Réponds par Vrai ou Faux, puis justifie ta réponse en une ou deux phrases.
\begin{enumerate}
\item Dans un syllogisme valide, la conclusion contient toujours le moyen terme.
\item L'induction permet de passer de l'observation de cas particuliers à une loi générale.
\item Un raisonnement peut être valide (logiquement correct) tout en ayant une conclusion fausse, si une prémisse est fausse.
\item L'analogie consiste à conclure que deux choses semblables sur certains points le sont sur tous les points.
\end{enumerate}
\end{exercice}

\begin{exercice}[Définitions]
Définis en deux ou trois phrases chacune des notions suivantes, en donnant pour chacune un exemple tiré de la vie courante camerounaise :
\begin{enumerate}
\item la prémisse ;
\item la conclusion ;
\item le paralogisme ;
\item le syllogisme.
\end{enumerate}
\end{exercice}

\begin{exercice}[Repérage des termes]
On considère le syllogisme suivant : « Tous les élèves sérieux réussissent au Probatoire ; or Amina est une élève sérieuse ; donc Amina réussira au Probatoire. »
\begin{enumerate}
\item Identifie la majeure, la mineure et la conclusion.
\item Relève le grand terme, le petit terme et le moyen terme.
\item Ce syllogisme est-il valide ? Justifie en une phrase.
\end{enumerate}
\end{exercice}

\courssub{Je m'entraîne}

\begin{exercice}[Complète le syllogisme]
On donne le raisonnement suivant : « Tous les fonctionnaires sont ponctuels ; or certains habitants de ce quartier sont fonctionnaires ; donc \ldots »
\begin{enumerate}
\item Complète la conclusion.
\item Identifie les trois termes du syllogisme (grand terme, petit terme, moyen terme).
\item Vérifie la validité du raisonnement à l'aide des règles du syllogisme vues en cours.
\end{enumerate}
\end{exercice}

\begin{exercice}[Classe les raisonnements]
Classe chacun des huit énoncés suivants selon son type (déduction, induction ou analogie) et justifie chaque classement en une phrase.
\begin{enumerate}
\item « Tous les métaux se dilatent sous l'effet de la chaleur ; or le fer est un métal ; donc le fer se dilate sous l'effet de la chaleur. »
\item « J'ai observé que les manguiers de mon village fleurissent chaque année en saison sèche ; donc tous les manguiers fleurissent en saison sèche. »
\item « Douala et Yaoundé sont deux grandes villes camerounaises ; à Douala, les embouteillages ralentissent les transports ; donc à Yaoundé aussi, les embouteillages ralentissent les transports. »
\item « Tous les mammifères allaitent leurs petits ; or la chèvre est un mammifère ; donc la chèvre allaite ses petits. »
\item « Les trois commerçants du marché que j'ai interrogés vendent leurs produits plus cher le samedi ; donc tous les commerçants vendent plus cher le samedi. »
\item « Mon frère aîné réussit bien en travaillant régulièrement ; mon frère cadet, qui lui ressemble, réussira donc aussi en travaillant régulièrement. »
\item « Aucun élève non inscrit ne peut composer ; or Jean n'est pas inscrit ; donc Jean ne peut pas composer. »
\item « Chaque fois qu'il a plu abondamment dans la plaine de Mbo, les récoltes de maïs ont été bonnes ; donc les pluies abondantes donnent de bonnes récoltes de maïs. »
\end{enumerate}
\end{exercice}

\begin{exercice}[Valide ou invalide ?]
Pour chacun des syllogismes suivants, dis s'il est valide ou invalide et justifie ta réponse en citant la règle respectée ou enfreinte.
\begin{enumerate}
\item « Tous les élèves de Première technique étudient la philosophie ; or Bienvenue est élève de Première technique ; donc Bienvenue étudie la philosophie. »
\item « Tous les footballeurs sont sportifs ; or tous les sportifs sont footballeurs ; donc tous les sportifs sont footballeurs. »
\item « Certains élèves sont ponctuels ; or certains élèves sont doués en mathématiques ; donc certains élèves ponctuels sont doués en mathématiques. »
\item « Aucun tricheur n'est honnête ; or Paul est honnête ; donc Paul n'est pas un tricheur. »
\end{enumerate}
\end{exercice}

\begin{exercice}[Construction : syllogisme valide et sophisme]
Sur le thème « l'usage du téléphone portable au lycée » :
\begin{enumerate}
\item Rédige un syllogisme valide, en indiquant clairement la majeure, la mineure et la conclusion.
\item Rédige ensuite volontairement un sophisme sur le même thème.
\item Explique en quelques lignes pourquoi le second raisonnement est fallacieux (quelle règle est enfreinte ? quelle apparence de validité le rend trompeur ?).
\end{enumerate}
\end{exercice}

\courssub{Je me prépare au Bac}

\begin{exercice}[Analyse de texte : le discours publicitaire]
Lis attentivement l'extrait suivant, tiré d'un discours publicitaire diffusé dans une ville camerounaise :

\emph{« Tous les jeunes modernes utilisent le téléphone "StarPhone". Or vous êtes un jeune moderne : achetez donc le "StarPhone" ! D'ailleurs, tous ceux qui ont acheté ce téléphone dans notre quartier en sont satisfaits ; il est donc certain que tous les acheteurs du Cameroun en seront satisfaits. Et puis, ce téléphone est comme un ordinateur : or un ordinateur rend intelligent ; le "StarPhone" vous rendra donc intelligent. »}

\begin{enumerate}
\item Releve dans ce texte deux raisonnements distincts et recopie-les.
\item Pour chacun, précise son type (déduction, induction ou analogie) en justifiant ta réponse.
\item Décele l'erreur de raisonnement commise dans chaque cas : s'agit-il d'un paralogisme ou d'un sophisme ? Quelle règle du raisonnement correct est enfreinte ?
\item Rédige un paragraphe d'une quinzaine de lignes dans lequel tu montres pourquoi il faut se méfier des arguments publicitaires, en t'appuyant sur l'analyse précédente et sur un exemple de la vie courante camerounaise.
\end{enumerate}
\end{exercice}

\begin{exercice}[Raisonner dans un débat de la vie courante]
Au cours d'un débat au foyer socio-éducatif d'un lycée technique de Yaoundé, un élève déclare : « Tous les élèves qui réussissent au Baccalauréat technique travaillent régulièrement ; or mon voisin travaille régulièrement ; donc mon voisin réussira forcément. » Un autre élève répond : « J'ai vu trois élèves de notre lycée réussir sans jamais ouvrir leurs cahiers ; donc aucun élève n'a besoin de travailler pour réussir. »
\begin{enumerate}
\item Reformule le raisonnement du premier élève sous la forme d'un syllogisme (majeure, mineure, conclusion) et identifie ses trois termes.
\item Ce raisonnement est-il valide ? Justifie ta réponse en citant la règle concernée.
\item Quel type de raisonnement le second élève utilise-t-il ? Pourquoi sa conclusion est-elle abusive ?
\item Propose, pour chacun des deux élèves, une version corrigée de son raisonnement, de sorte qu'il devienne acceptable.
\item Rédige et justifie ta propre conclusion sur la question débattue, en une dizaine de lignes.
\end{enumerate}
\end{exercice}

\begin{exercice}[Dissertation guidée : sujet type Baccalauréat]
Sujet : \emph{« Un raisonnement valide garantit-il toujours la vérité de la conclusion ? »}

Travail à faire :
\begin{enumerate}
\item Analyse du sujet : dégage le sens des mots-clés (\emph{raisonnement valide}, \emph{garantir}, \emph{vérité}, \emph{conclusion}) et formule la problématique en une question précise.
\item Rédige une introduction complète (accroche, définition des termes, problématique, annonce du plan).
\item Première partie : montre, à l'aide d'un syllogisme valide dont les prémisses sont vraies (exemple camerounais), que la validité formelle conduit à une conclusion vraie.
\item Deuxième partie : montre, à l'aide d'un raisonnement valide construit sur une prémisse fausse (exemple camerounais), que la validité ne suffit pas à garantir la vérité ; distingue clairement \emph{validité} et \emph{vérité}.
\item Rédige une conclusion qui répond directement à la problématique et ouvre sur une question (par exemple : quelles conditions faut-il réunir pour qu'un raisonnement soit à la fois valide et vrai ?).
\end{enumerate}
\end{exercice}

\newpage

\coursec{Corrigés des exercices}

\begin{corrige}[Exercice 1 — QCM : les bases du raisonnement]
\begin{enumerate}
\item Réponse \textbf{b)}. Un raisonnement est un enchaînement de jugements dont l'un, la \cle{conclusion}, découle nécessairement des autres, les \cle{prémisses}. Les réponses a), c) et d) désignent respectivement une opinion, un récit et une émotion, qui ne comportent aucun enchaînement logique.
\item Réponse \textbf{b)}. Le \cle{moyen terme} figure dans les deux prémisses (il relie le grand terme au petit terme) mais disparaît de la conclusion. S'il figurait dans la conclusion, le syllogisme serait fautif (règle : un syllogisme n'a que trois termes).
\item Réponse \textbf{c)}. La \cle{déduction} va du général au particulier : d'une loi universelle, elle tire une conséquence particulière. L'induction fait le trajet inverse (du particulier au général) et l'analogie va du particulier au particulier.
\item Réponse \textbf{b)}. Le \cle{sophisme} est un raisonnement faux présenté \emph{volontairement} comme valide, dans l'intention de tromper ou de persuader. Si l'erreur est involontaire, on parle de \cle{paralogisme}.
\end{enumerate}
\end{corrige}

\begin{attention}[Points de vigilance]
Ne pas confondre sophisme (intention de tromper) et paralogisme (erreur de bonne foi) : c'est la question la plus fréquente au QCM du Probatoire technique. Bien mémoriser aussi la place du moyen terme : \emph{dans les deux prémisses, jamais dans la conclusion}.
\end{attention}

\begin{corrige}[Exercice 2 — Vrai ou faux ? Justifie]
\begin{enumerate}
\item \textbf{Faux.} Le moyen terme figure dans les deux prémisses mais \emph{jamais} dans la conclusion. La conclusion ne contient que le petit terme (sujet) et le grand terme (attribut). Exemple : « Tous les hommes sont mortels ; or Socrate est un homme ; donc Socrate est mortel » : le moyen terme « homme » a disparu de la conclusion.
\item \textbf{Vrai.} L'\cle{induction} est précisément le raisonnement qui part de l'observation répétée de cas particuliers pour formuler une loi générale. Exemple : après avoir observé que plusieurs manguiers fleurissent en saison sèche, on conclut que tous les manguiers fleurissent en saison sèche.
\item \textbf{Vrai.} La \cle{validité} est une propriété de la \emph{forme} du raisonnement, indépendante de la vérité des prémisses. Exemple : « Tous les Camerounais parlent anglais ; or Amina est camerounaise ; donc Amina parle anglais » : la forme est correcte, mais la majeure étant fausse, la conclusion peut être fausse.
\item \textbf{Vrai.} L'\cle{analogie} conclut que deux réalités semblables sur certains points le sont aussi sur un autre point. Mais cette conclusion n'est que \emph{probable} : la ressemblance partielle ne garantit pas la ressemblance totale, d'où le danger des analogies abusives.
\end{enumerate}
\end{corrige}

\begin{attention}[Point de vigilance]
La validité ne garantit la vérité de la conclusion que si les prémisses sont vraies (raisonnement \emph{solide}). Un raisonnement valide à prémisses fausses peut mener à une conclusion fausse.
\end{attention}

\begin{corrige}[Exercice 3 — Définitions]
\begin{enumerate}
\item \textbf{La prémisse} : proposition (jugement) posée au départ d'un raisonnement et qui sert de fondement à la conclusion. Un syllogisme comporte deux prémisses : la majeure et la mineure. \emph{Exemple} : dans « Tous les élèves inscrits peuvent composer ; or Jean est inscrit ; donc Jean peut composer », les deux premières propositions sont les prémisses.
\item \textbf{La conclusion} : proposition nouvelle qui découle des prémisses et que le raisonnement vise à établir. \emph{Exemple} : dans le raisonnement précédent, « Jean peut composer » est la conclusion, annoncée par le mot « donc ».
\item \textbf{Le paralogisme} : raisonnement faux commis \emph{involontairement}, par ignorance ou inattention, sans intention de tromper. \emph{Exemple} : un élève de Ngaoundéré conclut « tous les chauffeurs de moto roulent vite » après en avoir vu deux seulement : il généralise trop vite, de bonne foi.
\item \textbf{Le syllogisme} : raisonnement déductif composé de deux prémisses (majeure et mineure) et d'une conclusion, mettant en jeu exactement trois termes : le grand terme, le petit terme et le moyen terme. \emph{Exemple} : « Tous les commerçants du marché paient une taxe ; or Mme Etoundi est commerçante au marché ; donc Mme Etoundi paie une taxe. »
\end{enumerate}
\end{corrige}

\begin{attention}[Point de vigilance]
Une définition correcte comporte le genre prochain (un raisonnement, une proposition\ldots) et la différence spécifique (ce qui distingue la notion des notions voisines). Évite de définir un mot par lui-même (« la conclusion est ce qui conclut »).
\end{attention}

\begin{corrige}[Exercice 4 — Repérage des termes]
Syllogisme : « Tous les élèves sérieux réussissent au Probatoire ; or Amina est une élève sérieuse ; donc Amina réussira au Probatoire. »
\begin{enumerate}
\item \textbf{Majeure} (prémisse contenant le grand terme) : « Tous les élèves sérieux réussissent au Probatoire. » \textbf{Mineure} (prémisse contenant le petit terme) : « Amina est une élève sérieuse. » \textbf{Conclusion} : « Amina réussira au Probatoire. »
\item \textbf{Grand terme} (attribut de la conclusion) : « réussir au Probatoire ». \textbf{Petit terme} (sujet de la conclusion) : « Amina ». \textbf{Moyen terme} (présent dans les deux prémisses, absent de la conclusion) : « élève sérieux ».
\item Le syllogisme est \textbf{valide} : le moyen terme « élève sérieux » est pris universellement dans la majeure (« \emph{tous} les élèves sérieux »), il relie correctement le petit terme au grand terme, et il n'y a que trois termes. La conclusion suit donc nécessairement des prémisses.
\end{enumerate}
\end{corrige}

\begin{attention}[Rappel de règle]
Dans un syllogisme valide (mode Barbara : AAA-1), la majeure est universelle affirmative, la mineure universelle affirmative, la conclusion universelle affirmative. Le moyen terme est distribué (pris universellement) dans la majeure.
\end{attention}

\begin{corrige}[Exercice 5 — Complète le syllogisme]
\begin{enumerate}
\item Conclusion : « \ldots donc \textbf{certains habitants de ce quartier sont ponctuels}. »
\item \textbf{Grand terme} : « ponctuel » (attribut de la conclusion) ; \textbf{petit terme} : « habitants de ce quartier » (sujet de la conclusion) ; \textbf{moyen terme} : « fonctionnaires » (dans les deux prémisses, absent de la conclusion).
\item Vérification des règles :
\begin{itemize}
\item Trois termes exactement : respectée (fonctionnaires, ponctuels, habitants du quartier).
\item Le moyen terme est pris universellement au moins une fois : respectée (« \emph{tous} les fonctionnaires » dans la majeure).
\item La conclusion ne contient pas le moyen terme : respectée.
\item La conclusion est particulière (« certains ») parce que la mineure est particulière : respectée (la conclusion suit toujours la prémisse la plus faible).
\end{itemize}
Le raisonnement est donc \textbf{valide} (mode Darii : AII-1).
\end{enumerate}
\end{corrige}

\begin{corrige}[Exercice 6 — Classe les raisonnements]
\begin{enumerate}
\item \textbf{Déduction} : on applique une loi générale (tous les métaux se dilatent) à un cas particulier (le fer).
\item \textbf{Induction} : de l'observation de cas particuliers (les manguiers du village), on passe à une loi générale (tous les manguiers). La généralisation reste ici fragile car l'échantillon est limité à un seul village.
\item \textbf{Analogie} : on va du particulier au particulier (de Douala à Yaoundé) en s'appuyant sur une ressemblance (deux grandes villes camerounaises). Conclusion seulement probable.
\item \textbf{Déduction} : du général (tous les mammifères) au particulier (la chèvre).
\item \textbf{Induction} : de trois cas observés à une loi générale sur tous les commerçants ; généralisation hâtive car l'échantillon est trop restreint.
\item \textbf{Analogie} : on transfère au frère cadet une propriété constatée chez le frère aîné, en raison de leur ressemblance supposée.
\item \textbf{Déduction} : d'une règle générale (aucun élève non inscrit ne compose), on tire la conséquence particulière pour Jean.
\item \textbf{Induction} : de plusieurs observations répétées (pluies abondantes suivies de bonnes récoltes), on conclut à une loi générale.
\end{enumerate}
\end{corrige}

\begin{attention}[Point de vigilance]
Critère de classement : regarde le \emph{sens du trajet} — général vers particulier (déduction), particulier vers général (induction), particulier vers particulier (analogie). Ne classe pas selon la vérité de la conclusion : une induction peut être vraie et une déduction construite sur une prémisse fausse.
\end{attention}

\begin{corrige}[Exercice 7 — Valide ou invalide ?]
\begin{enumerate}
\item \textbf{Valide.} Majeure universelle, moyen terme « élève de Première technique » pris universellement et absent de la conclusion, trois termes seulement : toutes les règles sont respectées (mode Barbara).
\item \textbf{Invalide.} La mineure « tous les sportifs sont footballeurs » inverse abusivement la majeure : de « tous les footballeurs sont sportifs », on ne peut pas conclure que tous les sportifs sont footballeurs (un athlète ou un handballeur est sportif sans être footballeur). La règle enfreinte : on ne peut pas étendre dans la conclusion un terme qui n'était pas pris universellement dans les prémisses (procès illicite du grand terme ou conversion simple abusive).
\item \textbf{Invalide.} Les deux prémisses sont particulières (« certains\ldots ») : or de deux prémisses particulières on ne peut rien conclure. Les élèves ponctuels et les élèves doués en mathématiques peuvent être deux groupes entièrement distincts (moyen terme non distribué).
\item \textbf{Valide.} « Aucun tricheur n'est honnête » équivaut à « tous les honnêtes sont non-tricheurs » ; Paul étant honnête, il n'est pas tricheur. La conclusion négative est permise car l'une des prémisses est négative, et le moyen terme « honnête » est pris universellement dans la majeure (« \emph{aucun} tricheur » distribue le sujet et l'attribut). Mode Celarent (EAE-1).
\end{enumerate}
\end{corrige}

\begin{attention}[Règles de validité à retenir]
1. Trois termes exactement. 2. Le moyen terme doit être pris universellement au moins une fois. 3. Un terme non distribué dans les prémisses ne peut l'être dans la conclusion. 4. Deux prémisses négatives $\rightarrow$ pas de conclusion. 5. Deux prémisses particulières $\rightarrow$ pas de conclusion. 6. Si une prémisse est négative, la conclusion est négative (et inversement).
\end{attention}

\begin{corrige}[Exercice 8 — Construction : syllogisme valide et sophisme]
\begin{enumerate}
\item \textbf{Syllogisme valide.} Majeure : « Tous les appareils qui perturbent les cours doivent être éteints en classe. » Mineure : « Or le téléphone portable allumé est un appareil qui perturbe les cours. » Conclusion : « Donc le téléphone portable doit être éteint en classe. » (Trois termes, moyen terme « appareil qui perturbe les cours » pris universellement, absent de la conclusion : raisonnement valide, mode Barbara.)
\item \textbf{Sophisme (exemple).} « Tous les jeunes branchés ont un smartphone ; or toi tu as un smartphone ; donc tu es forcément un jeune branché — et un jeune branché doit laisser son téléphone allumé en classe. »
\item \textbf{Explication.} Ce raisonnement est fallacieux parce qu'il inverse la majeure : de « tous les jeunes branchés ont un smartphone », on ne peut pas conclure que tous les possesseurs de smartphone sont branchés (règle enfreinte : un terme non universel dans la prémisse ne peut devenir universel dans la conclusion ; c'est le sophisme de l'\emph{affirmation du conséquent} ou conversion simple abusive). L'apparence de validité vient de la structure en trois propositions qui imite le syllogisme correct, et de la majeure vraie, qui endort la vigilance de l'auditeur.
\end{enumerate}
\end{corrige}

\begin{attention}[Sophisme vs Paralogisme]
Ici l'intention est de manipuler pour obtenir un privilège (laisser le téléphone allumé) : c'est bien un \cle{sophisme}. Si un élève croyait sincèrement que posséder un smartphone suffit à être « branché », ce serait un \cle{paralogisme}.
\end{attention}

\begin{corrige}[Exercice 9 — Analyse de texte : le discours publicitaire]
\begin{enumerate}
\item Deux raisonnements relevés :
\begin{itemize}
\item R1 : « Tous les jeunes modernes utilisent le téléphone "StarPhone". Or vous êtes un jeune moderne : achetez donc le "StarPhone" ! »
\item R2 : « Tous ceux qui ont acheté ce téléphone dans notre quartier en sont satisfaits ; il est donc certain que tous les acheteurs du Cameroun en seront satisfaits. »
\item R3 (complémentaire) : « Ce téléphone est comme un ordinateur : or un ordinateur rend intelligent ; le "StarPhone" vous rendra donc intelligent. »
\end{itemize}
\item \textbf{Types} : R1 est une \textbf{déduction} (d'une prétendue loi générale sur les jeunes modernes, on tire une conséquence pour un individu). R2 est une \textbf{induction} (des acheteurs d'un quartier à tous les acheteurs du Cameroun). R3 est une \textbf{analogie} (du téléphone à l'ordinateur, par ressemblance supposée).
\item \textbf{Erreurs} : dans R1, la majeure est \emph{fausse} (tous les jeunes modernes n'utilisent pas ce téléphone) et la conclusion glisse du constat (« utilisent ») à l'injonction (« achetez ») : c'est un \textbf{sophisme}, car l'erreur est volontaire, destinée à vendre. R2 est une \textbf{généralisation hâtive} : un quartier n'est pas représentatif de tout le Cameroun ; présentée comme « certaine », c'est aussi un sophisme (ou un paralogisme si le vendeur est de bonne foi). R3 repose sur une \textbf{analogie abusive} : la ressemblance technique entre téléphone et ordinateur ne garantit nullement l'effet « rendre intelligent » — d'ailleurs la prémisse « un ordinateur rend intelligent » est elle-même fausse.
\item \textbf{Paragraphe rédigé (modèle)} : Il faut se méfier des arguments publicitaires parce qu'ils visent à persuader, non à établir la vérité. Le discours analysé enfourne trois procédés fallacieux : une déduction bâtie sur une majeure fausse (« tous les jeunes modernes\ldots »), une induction qui généralise à tout le pays la satisfaction de quelques acheteurs d'un quartier, et une analogie abusive entre téléphone et ordinateur. Ces raisonnements ont l'\emph{apparence} de la rigueur — prémisses, connecteur « donc », conclusion — mais enfreignent les règles du raisonnement correct : prémisses vraies, échantillon suffisant, ressemblances pertinentes. Dans la vie courante camerounaise, on entend souvent à la radio qu'un produit « utilisé par toutes les grandes familles de la ville » doit être adopté par tous : c'est le même sophisme d'appel à la majorité. L'esprit critique consiste à vérifier les prémisses, à mesurer l'échantillon et à repérer les glissements de sens. Raisonner correctement, c'est ainsi se protéger de la manipulation commerciale.
\end{enumerate}
\textbf{Barème indicatif (20 points)} : relevé des raisonnements (4 pts) ; types justifiés (4 pts) ; erreurs identifiées avec la règle enfreinte et la distinction sophisme/paralogisme (6 pts) ; paragraphe argumenté avec exemple camerounais (6 pts).
\end{corrige}

\begin{attention}[Points de vigilance]
Cite toujours la \emph{règle} enfreinte (prémisse fausse, généralisation hâtive, analogie abusive), pas seulement le nom de l'erreur. Distingue sophisme (intention de tromper) et paralogisme (bonne foi). Le paragraphe doit mobiliser l'analyse précédente et un exemple personnel : un simple résumé du texte ne suffit pas.
\end{attention}

\begin{corrige}[Exercice 10 — Raisonner dans un débat de la vie courante]
\begin{enumerate}
\item \textbf{Syllogisme du premier élève.} Majeure : « Tous les élèves qui réussissent au Baccalauréat technique travaillent régulièrement. » Mineure : « Or mon voisin travaille régulièrement. » Conclusion : « Donc mon voisin réussira. » Grand terme : « réussir au Baccalauréat technique » (attribut de la conclusion) ; petit terme : « mon voisin » (sujet de la conclusion) ; moyen terme : « travailler régulièrement ».
\item \textbf{Validité} : le raisonnement est \textbf{invalide}. Le moyen terme « travailler régulièrement » n'est pris universellement dans \emph{aucune} des deux prémisses : dans la majeure, il est l'attribut d'une proposition affirmative (A), donc pris particulièrement ; dans la mineure, il est l'attribut d'une proposition affirmative (A), donc pris particulièrement. C'est la faute du \emph{moyen terme non distribué} : de « tous les A sont B » et « C est B », on ne peut pas conclure « C est A ».
\item Le second élève utilise une \textbf{induction} : de trois cas observés, il conclut à une loi générale négative (« aucun élève n'a besoin de travailler »). Sa conclusion est abusive parce que l'échantillon est dérisoire (trois élèves sur des centaines) et qu'il généralise hâtivement : trois exceptions apparentes ne détruisent pas une régularité générale.
\item \textbf{Versions corrigées.} Premier élève : « Tous les élèves qui travaillent régulièrement mettent toutes les chances de leur côté ; or mon voisin travaille régulièrement ; donc mon voisin met toutes les chances de son côté » (conclusion ramenée à ce que les prémisses permettent). Second élève : « J'ai vu trois élèves réussir sans travailler ; donc \emph{quelques} élèves peuvent réussir sans travailler » (conclusion particulière, conforme à l'échantillon).
\item \textbf{Conclusion personnelle (modèle)} : Le travail régulier n'est pas une garantie absolue de réussite, mais il en est la condition la plus sûre. Le premier élève confond condition nécessaire et condition suffisante : travailler ne suffit pas toujours (il faut aussi des méthodes, de la santé, de bonnes conditions d'examen), mais ne pas travailler rend l'échec presque certain. Le second élève tire de rares exceptions une règle générale, ce qui est une généralisation hâtive. La position raisonnable est donc intermédiaire : le travail régulier augmente considérablement les probabilités de réussite au Baccalauréat technique, sans les porter à la certitude. C'est pourquoi le conseil sage reste de travailler régulièrement tout en soignant sa méthode.
\end{enumerate}
\textbf{Barème indicatif (20 points)} : reformulation et termes (4 pts) ; invalidité justifiée par la règle du moyen terme (4 pts) ; type d'induction et caractère abusif (4 pts) ; deux corrections acceptables (4 pts) ; conclusion personnelle argumentée (4 pts).
\end{corrige}

\begin{attention}[Erreur fréquente]
Ne pas dire « le syllogisme est faux » mais « le syllogisme est \emph{invalide} » (vice de forme). La fausseté concerne les propositions (prémisses ou conclusion), l'invalidité concerne l'enchaînement.
\end{attention}

\begin{corrige}[Exercice 11 — Dissertation guidée : « Un raisonnement valide garantit-il toujours la vérité de la conclusion ? »]
\begin{enumerate}
\item \textbf{Analyse du sujet.} Un \emph{raisonnement valide} est un raisonnement dont la forme est correcte : si les prémisses sont vraies, la conclusion ne peut pas être fausse. \emph{Garantir} signifie assurer avec certitude. La \emph{vérité} est l'accord entre ce qu'on affirme et la réalité. La \emph{conclusion} est la proposition obtenue à partir des prémisses. \textbf{Problématique} : la correction formelle d'un raisonnement suffit-elle à assurer que sa conclusion est vraie, ou la vérité exige-t-elle autre chose que la validité ?
\item \textbf{Introduction (modèle).} Au marché comme en classe, nous raisonnons sans cesse pour convaincre et pour décider. Mais tout raisonnement qui « sonne juste » dit-il pour autant la vérité ? On appelle raisonnement valide un enchaînement de jugements dont la conclusion découle nécessairement des prémisses ; la vérité, elle, concerne le rapport des énoncés à la réalité. Dès lors, un problème se pose : un raisonnement valide garantit-il toujours la vérité de la conclusion ? Nous verrons d'abord que la validité conduit à la vérité lorsque les prémisses sont vraies, puis qu'elle ne suffit pas à elle seule, car un raisonnement correct peut reposer sur des prémisses fausses.
\item \textbf{Première partie (modèle).} Considérons le syllogisme : « Tous les élèves inscrits au Probatoire technique peuvent composer ; or Chantal, élève du lycée technique de Douala, est inscrite ; donc Chantal peut composer. » La forme est valide (trois termes, moyen terme « inscrit » pris universellement) et les deux prémisses sont vraies. La conclusion est alors \emph{nécessairement} vraie : nul ne peut admettre les prémisses et nier la conclusion sans se contredire. La validité formelle est donc bien un chemin vers la vérité : elle transmet la vérité des prémisses à la conclusion. C'est la force de la déduction, illustrée par Aristote, inventeur du syllogisme.
\item \textbf{Deuxième partie (modèle).} Mais considérons maintenant : « Tous les élèves de Première technique sont des garçons ; or Divine est élève de Première technique ; donc Divine est un garçon. » La forme du raisonnement est exactement la même, donc parfaitement valide ; pourtant la conclusion est fausse, parce que la majeure est fausse. La validité ne crée pas la vérité : elle ne fait que la \emph{transporter}. Si la matière du raisonnement (les prémisses) est fausse, la machine logique, pourtant correcte, produit une conclusion fausse. Il faut donc distinguer rigoureusement la \emph{validité}, propriété de la forme, et la \emph{vérité}, propriété du contenu : un raisonnement n'est pleinement satisfaisant — on dit \emph{solide} ou \emph{correct} — que s'il est à la fois valide et construit sur des prémisses vraies.
\item \textbf{Conclusion (modèle).} Un raisonnement valide ne garantit donc pas toujours la vérité de la conclusion : il la garantit seulement \emph{si} les prémisses sont vraies. La logique formelle est un instrument puissant mais neutre, qui transmet aussi fidèlement l'erreur que la vérité. Pour qu'un raisonnement soit irréfutable, il faut réunir deux conditions : une forme valide et des prémisses vraies. Reste alors une question plus vaste : comment s'assurer, dans les sciences comme dans la vie courante, de la vérité des prémisses elles-mêmes ?
\end{enumerate}
\textbf{Barème indicatif (20 points)} : analyse des termes et problématique (3 pts) ; introduction complète en quatre temps (4 pts) ; première partie avec exemple valide et vrai (4 pts) ; deuxième partie avec distinction validité/vérité et exemple (5 pts) ; conclusion répondant à la problématique avec ouverture (2 pts) ; clarté et orthographe (2 pts).
\end{corrige}

\begin{attention}[Points de vigilance]
La distinction \cle{validité}/\cle{vérité} est le cœur du sujet : elle doit apparaître explicitement, avec les deux exemples camerounais demandés. Évite deux erreurs classiques : affirmer qu'un raisonnement valide est « toujours vrai » (contre-sens) ou qu'un raisonnement à prémisses fausses est « invalide » (il est valide mais non solide). L'introduction doit contenir les quatre temps : accroche, définitions, problématique, annonce du plan.
\end{attention}

\newpage

\coursec{Fiche bilan}

\courssub{Carte mentale de la leçon}
\begin{popfigure}
\begin{tikzpicture}[
  mindmap,
  concept color=popblue,
  level 1/.append style={level distance=4.5cm, sibling angle=60},
  level 2/.append style={level distance=3cm, sibling angle=45},
  every node/.style={font=\small, text=popdark},
  branch1/.style={concept color=popblue, grow=90},
  branch2/.style={concept color=poporange, grow=30},
  branch3/.style={concept color=popgreen, grow=-30},
  branch4/.style={concept color=poppurple, grow=-90},
  branch5/.style={concept color=poppink, grow=-150},
  branch6/.style={concept color=popgold, grow=150},
  text width=2.8cm, align=center
]
\node[concept, text width=3.2cm] {Le Raisonnement\\ \textit{(Logique formelle)}}
  child[branch1] {node[concept, text width=2.8cm] {1. Définition \& Opérations\\ -- Concept, Jugement\\ -- Inférence, Argumentation\\ -- Enthymème, Épichérème}}
  child[branch2] {node[concept, text width=2.8cm] {2. Le Syllogisme\\ -- Structure (3 props, 3 termes)\\ -- Figures (1 à 4) \& Modes\\ -- Règles de validité (4 classiques + prémisses)}}
  child[branch3] {node[concept, text width=2.8cm] {3. Types de raisonnement\\ -- Déduction (nécessaire)\\ -- Induction (probable)\\ -- Analogie / Abduction (heuristique)}}
  child[branch4] {node[concept, text width=2.8cm] {4. Règles du bon raisonnement\\ -- Identité, Non-contradiction\\ -- Tiers exclu, Raison suffisante\\ -- Validité formelle vs Vérité matérielle}}
  child[branch5] {node[concept, text width=2.8cm] {5. Erreurs : Paralogismes \& Sophismes\\ -- Formels (structure : 4 termes, moyen non distribué...)\\ -- Matériels (contenu)\\ -- Intention : Sophisme vs Paralogisme (inadvertance)}}
  child[branch6] {node[concept, text width=2.8cm] {6. Méthode \& Rédaction\\ -- Analyser / Formaliser (S, M, P)\\ -- Vérifier validité \& vérité\\ -- Classer \& Détecter erreur\\ -- Conclure rigoureusement}};
\end{tikzpicture}
\legende{Carte mentale synthétique de la leçon 14 : Le raisonnement (6 branches principales).}
\end{popfigure}

\courssub{Bilan des connaissances essentielles}
\begin{bilan}

\textbf{1. Définitions clés}
\begin{itemize}
  \item \cle{Raisonnement} : Opération mentale qui, à partir de propositions admises (\cle{prémisses}), en déduit une nouvelle (\cle{conclusion}) de manière nécessaire (déduction) ou probable (induction, analogie, abduction).
  \item \cle{Syllogisme} : Raisonnement \textbf{déductif} formé de trois propositions (deux prémisses, une conclusion) et trois termes (Majeur P, Mineur S, Moyen M).
  \item \cle{Déduction} : Passage du général au particulier (ou du général au général) ; conclusion \textbf{nécessaire} si prémisses vraies et forme valide (ex: Barbara).
  \item \cle{Induction} : Passage du particulier au général ; conclusion \textbf{probable} (généralisation à partir de cas observés), jamais certaine.
  \item \cle{Abduction} : Raisonnement d'hypothèse (inférence vers la meilleure explication) : fait surprenant C ; si A alors C ; donc A est plausible. Heuristique, non probante.
  \item \cle{Analogie} : Raisonnement par similarité (A a propriété p, B ressemble à A $\Rightarrow$ B a p) ; heuristique, non probante, force dépend de la pertinence de la ressemblance.
  \item \cle{Enthymème} : Syllogisme tronqué dont une prémisse (souvent la majeure, évidente) est implicite (ex: \og Il pleut, donc le sol est mouillé \fg $\leftarrow$ \og Si il pleut, le sol est mouillé \fg).
  \item \cle{Épichérème} : Syllogisme dont une prémisse est justifiée par une raison explicite (ex: \og Tout homme est mortel, car il est composé de parties corruptibles ; or Socrate est homme ; donc Socrate est mortel \fg).
  \item \cle{Paralogisme} : Erreur de raisonnement \textbf{involontaire} (inadvertance, manque de rigueur).
  \item \cle{Sophisme} : Erreur de raisonnement \textbf{intentionnelle} (visée de tromper, de persuader fallacieusement).
\end{itemize}

\textbf{2. Règles \& Théorèmes à connaître par cœur}
\begin{center}
\begin{tabularx}{\linewidth}{|l|X|}\hline
\textbf{Règle / Loi} & \textbf{Énoncé / Condition de validité} \\ \hline
\cle{Loi d'identité} & A est A. Un concept reste identique à lui-même dans le raisonnement (pas d'équivoque). \\ \hline
\cle{Loi de non-contradiction} & Impossible qu'A soit et ne soit pas en même temps et sous le même rapport. \\ \hline
\cle{Loi du tiers exclu} & Pour toute proposition, soit elle est vraie, soit sa négation est vraie (pas de milieu). \\ \hline
\cle{Principe de raison suffisante} & Rien n'est sans raison ; toute vérité a un fondement (cause ou motif). \\ \hline
\cle{Règles du syllogisme (4 classiques)} & 
\begin{enumerate}
  \item 3 termes exactement, ni plus ni moins (pas de \textit{quaternio terminorum}).
  \item Le moyen (M) ne doit pas figurer dans la conclusion.
  \item Le moyen (M) doit être universel (distribué) au moins une fois (pris en toute son étendue).
  \item Aucun terme (S ou P) n'est universel (distribué) dans la conclusion s'il ne l'était pas dans les prémisses.
\end{enumerate} \\ \hline
\cle{Règles des prémisses (qualité/quantité)} & 
\begin{itemize}
  \item Deux prémisses négatives $\Rightarrow$ pas de conclusion valide.
  \item Une prémisse négative $\Rightarrow$ conclusion obligatoirement négative.
  \item Deux prémisses particulières $\Rightarrow$ pas de conclusion valide.
  \item La conclusion suit la prémisse la plus faible : si une prémisse est négative, conclusion négative ; si une prémisse est particulière, conclusion particulière.
\end{itemize} \\ \hline
\end{tabularx}
\end{center}

\textbf{3. Structure du syllogisme (à mémoriser)}
\begin{itemize}
  \item \textbf{Prémisse majeure} : Contient le Majeur (P, prédicat de la conclusion) et le Moyen (M). Notation usuelle : M-P ou P-M.
  \item \textbf{Prémisse mineure} : Contient le Mineur (S, sujet de la conclusion) et le Moyen (M). Notation usuelle : S-M ou M-S.
  \item \textbf{Conclusion} : Relie S à P (S est P / S n'est pas P).
  \item \textbf{Figures} (position du terme moyen M) :
  \begin{enumerate}
    \item Figure 1 : M-P / S-M $\Rightarrow$ S-P (Le M est sujet de la majeure, prédicat de la mineure).
    \item Figure 2 : P-M / S-M $\Rightarrow$ S-P (Le M est prédicat des deux prémisses).
    \item Figure 3 : M-P / M-S $\Rightarrow$ S-P (Le M est sujet des deux prémisses).
    \item Figure 4 : P-M / M-S $\Rightarrow$ S-P (Le M est prédicat de la majeure, sujet de la mineure).
  \end{enumerate}
  \item \textbf{Mode} : Qualité (A, E, I, O) et Quantité des 3 propositions (ex: AAA-1 = \textit{Barbara}, EAE-1 = \textit{Celarent}, AII-1 = \textit{Darii}, EIO-1 = \textit{Ferio}). Seuls 19 modes (sur 256) sont valides (réduits à 24 avec les modes affaiblis).
\end{itemize}

\textbf{4. Distinction cruciale : Validité vs Vérité}
\begin{itemize}
  \item \textbf{Validité (forme)} : Respect des règles logiques. Si prémisses vraies $\Rightarrow$ conclusion \textbf{nécessairement} vraie. Ne garantit pas la vérité des prémisses.
  \item \textbf{Vérité (contenu)} : Conformité des prémisses à la réalité.
  \item \cle{Raisonnement sound (correct / démonstratif)} : Raisonnement \textbf{valide} ET aux \textbf{prémisses vraies}. Seule cette conjonction assure la certitude de la conclusion.
\end{itemize}

\textbf{5. Méthode express : Analyser \& Rédiger un raisonnement}
\begin{enumerate}
  \item \textbf{Repérer} la conclusion (indices : \og donc \fg, \og par conséquent \fg, \og c'est pourquoi \fg) et les prémisses (indices : \og car \fg, \og puisque \fg, \og en effet \fg, \og car \fg).
  \item \textbf{Formaliser} : Identifier S (sujet conclusion), P (prédicat conclusion), M (terme commun aux prémisses) ; écrire sous forme canonique (Majeure, Mineure, Conclusion). Faire apparaître l'enthymème si prémisse implicite.
  \item \textbf{Vérifier la validité formelle} : Compter les termes (3), vérifier distribution du M (règle 3), non-distribution de S/P en conclusion si non distribués en prémisses (règle 4), règles des prémisses (négatives, particulières), identifier Figure/Mode valide.
  \item \textbf{Vérifier la vérité matérielle} : Les prémisses sont-elles vraies dans la réalité ? (Valide + Vrai = \cle{Sound / Correct / Démonstratif}).
  \item \textbf{Classifier} : Déduction / Induction / Analogie / Abduction / Enthymème / Épichérème.
  \item \textbf{Détecter l'erreur} le cas échéant :
  \begin{itemize}
    \item \textit{Paralogismes formels} : 4 termes (équivoque), moyen non distribué, prémisses négatives/particulières, terme distribué en conclusion seulement.
    \item \textit{Paralogismes matériels} : Pétition de principe, cause fausse (\textit{post hoc ergo propter hoc}), généralisation hâtive.
    \item \textit{Sophismes} : Équivoque (volontaire), \textit{ad hominem}, \textit{ad populum}, \textit{ad verecundiam}, homme de paille, faux dilemme, \textit{argumentum ad baculum}.
  \end{itemize}
  \item \textbf{Rédiger la justification} : Structure type : \og Le raisonnement est un syllogisme en Figure 1 Mode AAA (\textit{Barbara}). Il respecte les 4 règles classiques et les règles des prémisses : il est donc \textbf{valide}. Par ailleurs, les prémisses sont vraies (fait établi / définition), donc le raisonnement est \textbf{sound (correct)} : la conclusion est certaine. \fg \textbf{OU} \og Le raisonnement commet un paralogisme de [nom précis, ex: moyen non distribué] car [règle violée explicite], il est donc \textbf{invalide}. \fg
\end{enumerate}

\end{bilan}

\begin{aretenir}[Checklist avant le Probatoire / Bac Tech]
$\square$ Savoir définir précisément : raisonnement, syllogisme, déduction, induction, abduction, analogie, paralogisme, sophisme, enthymème, épichérème, validité, vérité, soundness.
$\square$ Citer les 3 lois de la pensée (Identité, Non-contradiction, Tiers exclu) et le Principe de raison suffisante ; expliquer leur rôle dans le raisonnement.
$\square$ Décomposer un syllogisme quelconque : identifier Majeure, Mineure, Conclusion, termes M, S, P ; écrire en forme canonique.
$\square$ Appliquer les 4 règles classiques du syllogisme (3 termes, M absent conclusion, M universel 1 fois, pas d'extension en conclusion).
$\square$ Appliquer les règles des prémisses (négatives, particulières, conclusion suit la plus faible).
$\square$ Différencier validité (forme) et vérité (contenu) ; définir un raisonnement \og sound \fg (valide + prémisses vraies) et donner un exemple valide mais non sound.
$\square$ Classer un raisonnement présenté : Déductif (nécessaire), Inductif (probable), Analogique/Abductif (plausible/heuristique).
$\square$ Identifier les 4 figures du syllogisme et nommer les modes valides principaux de la Figure 1 (Barbara, Celarent, Darii, Ferio) + savoir qu'il existe 19 modes valides au total.
$\square$ Repérer et nommer les paralogismes formels (ex: 4 termes/équivoque, moyen non distribué, prémisses négatives, terme distribué en conclusion seulement) et matériels (pétition de principe, cause fausse, généralisation hâtive).
$\square$ Repérer et nommer les sophismes courants (équivoque volontaire, \textit{ad hominem}, \textit{ad populum}, \textit{ad verecundiam}, \textit{post hoc}, homme de paille, faux dilemme, \textit{ad baculum}).
$\square$ Savoir formaliser un énoncé courant en syllogisme explicite (faire apparaître l'enthymème : majeure ou mineure implicite).
$\square$ Rédiger une justification structurée, complète et terminologiquement exacte : Formalisation $\rightarrow$ Test validité (règles citées) $\rightarrow$ Test vérité $\rightarrow$ Conclusion finale (Valide/Invalide, Sound/Non sound, Type d'erreur le cas échéant).
\end{aretenir}

\findecours

\end{document}
